\input{preamble}

% OK, start here.
%
\begin{document}

\title{Compléments sur les morphismes de champs}


\maketitle

\phantomsection
\label{section-phantom}

\tableofcontents

\section{Introduction}
\label{section-introduction}

\noindent
Dans ce chapitre, nous poursuivons l'étude des propriétés des morphismes de
champs algébriques. Une référence, dans le cas des champs algébriques quasi-séparés
à diagonale représentable, est \cite{LM-B}.




\section{Conventions et abus de langage}
\label{section-conventions}

\noindent
Nous continuons d'employer les conventions et les abus de langage
introduits dans
Propriétés des champs, Section \ref{stacks-properties-section-conventions}.






\section{Épaississements}
\label{section-thickenings}

\noindent
La terminologie suivante n'est peut-être pas tout à fait standard, mais elle est commode.
Si $\mathcal{Y}$ est un sous-champ fermé d'un champ algébrique $\mathcal{X}$,
alors le morphisme $\mathcal{Y} \to \mathcal{X}$ est représentable.

\begin{definition}
\label{definition-thickening}
Épaississements.
\begin{enumerate}
\item On dit qu'un champ algébrique $\mathcal{X}'$ est un {\it épaississement}
d'un champ algébrique $\mathcal{X}$ si $\mathcal{X}$ est un sous-champ fermé
de $\mathcal{X}'$ et si les espaces topologiques associés sont égaux.
\item Étant donnés deux épaississements $\mathcal{X} \subset \mathcal{X}'$ et
$\mathcal{Y} \subset \mathcal{Y}'$, un {\it morphisme d'épaississements}
est un morphisme $f' : \mathcal{X}' \to \mathcal{Y}'$ de champs algébriques
tel que $f'|_\mathcal{X}$ se factorise par le
sous-champ fermé $\mathcal{Y}$. Dans cette situation, on pose
$f = f'|_\mathcal{X} : \mathcal{X} \to \mathcal{Y}$ et on dit que
$(f, f') : (\mathcal{X} \subset \mathcal{X}') \to
(\mathcal{Y} \subset \mathcal{Y}')$ est un morphisme d'épaississements.
\item Soit $\mathcal{Z}$ un champ algébrique. On définit de même les
{\it épaississements sur $\mathcal{Z}$} et les
{\it morphismes d'épaississements sur $\mathcal{Z}$}.
Cela signifie que les champs algébriques
$\mathcal{X}'$ et $\mathcal{Y}'$
sont munis d'un
morphisme structural vers $\mathcal{Z}$ et que $f'$ s'insère dans un
diagramme $2$-commutatif convenable de champs algébriques.
\end{enumerate}
\end{definition}

\noindent
Soit $\mathcal{X} \subset \mathcal{X}'$ un épaississement de champs algébriques.
Soit $U'$ un schéma et soit $U' \to \mathcal{X}'$ un morphisme lisse
surjectif. En posant $U = \mathcal{X} \times_{\mathcal{X}'} U'$, on obtient
un morphisme d'épaississements
$$
(U \subset U') \longrightarrow (\mathcal{X} \subset \mathcal{X}')
$$
et $U \to \mathcal{X}$ est un morphisme lisse surjectif. On peut souvent
déduire des propriétés de l'épaississement $\mathcal{X} \subset \mathcal{X}'$
des propriétés correspondantes de l'épaississement $U \subset U'$.
Parfois, par abus de langage, on dit qu'un morphisme
$\mathcal{X} \to \mathcal{X}'$ est un épaississement s'il s'agit d'une immersion
fermée induisant une bijection $|\mathcal{X}| \to |\mathcal{X}'|$.

\begin{lemma}
\label{lemma-thickening}
Soit $i : \mathcal{X} \to \mathcal{X}'$ un morphisme de champs algébriques.
Les conditions suivantes sont équivalentes :
\begin{enumerate}
\item $i$ est un épaississement de champs algébriques (au sens de l'abus de langage ci-dessus), et
\item $i$ est représentable par des espaces algébriques et
est un épaississement au sens de Propriétés des champs, Section
\ref{stacks-properties-section-properties-morphisms}.
\end{enumerate}
Dans ce cas, $i$ est une immersion fermée et un homéomorphisme universel.
\end{lemma}

\begin{proof}
D'après Compléments sur les morphismes d'espaces, Lemmes
\ref{spaces-more-morphisms-lemma-descending-property-thickening} et
\ref{spaces-more-morphisms-lemma-base-change-thickening}
la propriété $P$ qu'un morphisme d'espaces algébriques soit un
épaississement (du premier ordre) est locale pour la topologie fpqc sur la base et stable par
changement de base. La discussion de Propriétés des champs, Section
\ref{stacks-properties-section-properties-morphisms} s'applique donc bien.
Cela étant dit, l'équivalence de (1) et (2) résulte du
fait que $P = P_1 + P_2$, où $P_1$ est la propriété d'être
une immersion fermée et $P_2$ celle d'être surjectif.
(À strictement parler, le lecteur devrait également consulter
Compléments sur les morphismes d'espaces, Définition
\ref{spaces-more-morphisms-definition-thickening},
Propriétés des champs, Définition \ref{stacks-properties-definition-immersion}
et la discussion qui suit, Morphismes d'espaces, Lemme
\ref{spaces-morphisms-lemma-surjective-representable},
Propriétés des champs, Section \ref{stacks-properties-section-surjective},
pour constater que toutes les notions coïncident.)
La dernière assertion est claire d'après ce qui précède.
\end{proof}

\noindent
Nous utiliserons le lemme sans autre mention. Les mêmes références,
Compléments sur les morphismes d'espaces, Lemmes
\ref{spaces-more-morphisms-lemma-descending-property-thickening} et
\ref{spaces-more-morphisms-lemma-base-change-thickening}
permettent de définir comme suit un épaississement du premier ordre.

\begin{definition}
\label{definition-first-order-thickening}
On dit qu'un champ algébrique $\mathcal{X}'$ est un {\it épaississement du premier ordre}
d'un champ algébrique $\mathcal{X}$ si $\mathcal{X}$ est un sous-champ fermé
de $\mathcal{X}'$ et si $\mathcal{X} \to \mathcal{X}'$ est un épaississement
du premier ordre au sens de Propriétés des champs, Section
\ref{stacks-properties-section-properties-morphisms}.
\end{definition}

\noindent
Si $(U \subset U') \to (\mathcal{X} \subset \mathcal{X}')$ est un
recouvrement lisse par un schéma comme ci-dessus, cela signifie simplement que $U \subset U'$
est un épaississement du premier ordre. Formulons maintenant les lemmes indispensables.

\begin{lemma}
\label{lemma-base-change-thickening}
Soit $\mathcal{Y} \subset \mathcal{Y}'$ un épaississement de champs algébriques.
Soit $\mathcal{X}' \to \mathcal{Y}'$ un morphisme de champs algébriques
et posons $\mathcal{X} = \mathcal{Y} \times_{\mathcal{Y}'} \mathcal{X}'$.
Alors
$(\mathcal{X} \subset \mathcal{X}') \to (\mathcal{Y} \subset \mathcal{Y}')$
est un morphisme d'épaississements. Si $\mathcal{Y} \subset \mathcal{Y}'$ est un
épaississement du premier ordre, alors $\mathcal{X} \subset \mathcal{X}'$ est un
épaississement du premier ordre.
\end{lemma}

\begin{proof}
Voir la discussion ci-dessus, Propriétés des champs, Section
\ref{stacks-properties-section-properties-morphisms}, et
et Compléments sur les morphismes d'espaces, Lemme
\ref{spaces-more-morphisms-lemma-base-change-thickening}.
\end{proof}

\begin{lemma}
\label{lemma-composition-thickening}
Si $\mathcal{X} \subset \mathcal{X}'$ et $\mathcal{X}' \subset \mathcal{X}''$
sont des épaississements de champs algébriques, alors
$\mathcal{X} \subset \mathcal{X}''$ en est également un.
\end{lemma}

\begin{proof}
Voir la discussion ci-dessus, Propriétés des champs, Section
\ref{stacks-properties-section-properties-morphisms}, et
et Compléments sur les morphismes d'espaces, Lemme
\ref{spaces-more-morphisms-lemma-composition-thickening}
\end{proof}

\begin{example}
\label{example-reduction-thickening}
Soit $\mathcal{X}'$ un champ algébrique. Alors $\mathcal{X}'$ est un épaississement
de son réduit $\mathcal{X}'_{red}$, voir
Propriétés des champs, Définition
\ref{stacks-properties-definition-reduced-induced-stack}.
De plus, si $\mathcal{X} \subset \mathcal{X}'$ est un épaississement
de champs algébriques, alors
$\mathcal{X}'_{red} = \mathcal{X}_{red} \subset \mathcal{X}$.
Autrement dit, $\mathcal{X} = \mathcal{X}'_{red}$ si et seulement
si $\mathcal{X}$ est un champ algébrique réduit.
\end{example}

\begin{lemma}
\label{lemma-reduced-diagonal}
Soit $(f, f') : (\mathcal{X} \subset \mathcal{X}') \to
(\mathcal{Y} \subset \mathcal{Y}')$ un morphisme d'épaississements
de champs algébriques. Alors
$\mathcal{X} \times_\mathcal{Y} \mathcal{X} \to
\mathcal{X}' \times_{\mathcal{Y}'} \mathcal{X}'$
est un épaississement, et le diagramme canonique
$$
\xymatrix{
\mathcal{X} \ar[r]_-\Delta \ar[d] &
\mathcal{X} \times_\mathcal{Y} \mathcal{X} \ar[d] \\
\mathcal{X}' \ar[r]^-{\Delta'} &
\mathcal{X}' \times_{\mathcal{Y}'} \mathcal{X}'
}
$$
est cartésien.
\end{lemma}

\begin{proof}
Puisque $\mathcal{X} \to \mathcal{Y}'$ se factorise par le
sous-champ fermé $\mathcal{Y}$, on a
$\mathcal{X} \times_\mathcal{Y} \mathcal{X} =
\mathcal{X} \times_{\mathcal{Y}'} \mathcal{X}$.
Par conséquent,
$\mathcal{X} \times_\mathcal{Y} \mathcal{X} \to
\mathcal{X}' \times_{\mathcal{Y}'} \mathcal{X}'$
est isomorphe au composé
$$
\mathcal{X} \times_{\mathcal{Y}'} \mathcal{X} \to
\mathcal{X} \times_{\mathcal{Y}'} \mathcal{X}' \to
\mathcal{X}' \times_{\mathcal{Y}'} \mathcal{X}'
$$
dont les deux flèches sont des épaississements, puisqu'elles sont obtenues par changement de base
d'épaississements (Lemme \ref{lemma-base-change-thickening}).
Le composé est donc lui aussi un épaississement
(Lemme \ref{lemma-composition-thickening}).
Puisque $\mathcal{X} \to \mathcal{X}'$ est un monomorphisme,
la dernière assertion du lemme résulte de
Propriétés des champs, Lemme
\ref{stacks-properties-lemma-monomorphism-diagonal}
appliqué à $\mathcal{X} \to \mathcal{X}' \to \mathcal{Y}'$.
\end{proof}

\begin{lemma}
\label{lemma-thickening-diagonals}
Soit $(f, f') : (\mathcal{X} \subset \mathcal{X}') \to
(\mathcal{Y} \subset \mathcal{Y}')$ un morphisme d'épaississements
de champs algébriques.
Soient $\Delta : \mathcal{X} \to \mathcal{X} \times_\mathcal{Y} \mathcal{X}$ et
$\Delta' : \mathcal{X}' \to \mathcal{X}' \times_{\mathcal{Y}'} \mathcal{X}'$
les morphismes diagonaux correspondants.
Alors chacune des propriétés de la liste suivante est satisfaite par $\Delta$ si
et seulement si elle est satisfaite par $\Delta'$ :
(a) représentable par des schémas, (b) affine, (c) surjectif, (d) quasi-compact,
(e) universellement fermé, (f) entier, (g) quasi-séparé, (h) séparé,
(i) universellement injectif, (j) universellement ouvert, (k) localement quasi-fini,
(l) fini, (m) non ramifié, (n) monomorphisme, (o) immersion,
(p) immersion fermée et (q) propre.
\end{lemma}

\begin{proof}
Observons que
$$
(\Delta, \Delta') :
(\mathcal{X} \subset \mathcal{X}')
\longrightarrow
(\mathcal{X} \times_\mathcal{Y} \mathcal{X} \subset
\mathcal{X}' \times_{\mathcal{Y}'} \mathcal{X}')
$$
est un morphisme d'épaississements (Lemme \ref{lemma-reduced-diagonal}).
De plus, $\Delta$ et $\Delta'$ sont
représentables par des espaces algébriques d'après
Morphismes de champs, Lemme \ref{stacks-morphisms-lemma-properties-diagonal}.
Par conséquent, via la discussion de
Propriétés des champs, Section
\ref{stacks-properties-section-properties-morphisms}
le lemme résulte, dans les cas (a), (b), (c), (d),
(e), (f), (g), (h), (i) et (j), du
Lemme de Compléments sur les morphismes d'espaces
\ref{spaces-more-morphisms-lemma-thicken-property-morphisms}.

\medskip\noindent
Le Lemme \ref{lemma-reduced-diagonal} nous dit que
$\mathcal{X} = (\mathcal{X} \times_\mathcal{Y} \mathcal{X})
\times_{(\mathcal{X}' \times_{\mathcal{Y}'} \mathcal{X}')} \mathcal{X}'$.
De plus, $\Delta$ et $\Delta'$ sont localement de type fini d'après
le Lemme de Morphismes de champs déjà cité,
Morphismes de champs, Lemme \ref{stacks-morphisms-lemma-properties-diagonal}.
Le résultat pour les cas (k), (l), (m), (n), (o), (p) et (q) découle donc du
Lemme de Compléments sur les morphismes d'espaces
\ref{spaces-more-morphisms-lemma-properties-that-extend-over-thickenings}.
\end{proof}

\noindent
On obtient comme conséquence le résultat agréable suivant.

\begin{lemma}
\label{lemma-thickening-properties}
\begin{reference}
\cite[Théorème 2.2.5]{Conrad-moduli}
\end{reference}
Soit $\mathcal{X} \subset \mathcal{X}'$ un épaississement de
champs algébriques. Alors
\begin{enumerate}
\item $\mathcal{X}$ est un espace algébrique si et seulement si $\mathcal{X}'$
est un espace algébrique,
\item $\mathcal{X}$ est un schéma si et seulement si $\mathcal{X}'$ est un schéma,
\item $\mathcal{X}$ est DM si et seulement si $\mathcal{X}'$ est DM,
\item $\mathcal{X}$ est quasi-DM si et seulement si $\mathcal{X}'$ est quasi-DM,
\item $\mathcal{X}$ est séparé si et seulement si $\mathcal{X}'$ est séparé,
\item $\mathcal{X}$ est quasi-séparé si et seulement si $\mathcal{X}'$ est
quasi-séparé, et
\item en ajouter ici.
\end{enumerate}
\end{lemma}

\begin{proof}
Dans chaque cas, on se ramène à une question portant sur la diagonale, puis
on applique le Lemme \ref{lemma-thickening-diagonals} au
morphisme d'épaississements
$$
(\mathcal{X} \subset \mathcal{X}') \to
\left(\Spec(\mathbf{Z}) \subset \Spec(\mathbf{Z})\right)
$$
On procède ainsi après avoir considéré
$\mathcal{X} \subset \mathcal{X}'$ comme un épaississement de champs algébriques
sur $\Spec(\mathbf{Z})$ au moyen de
Champs algébriques, Définition \ref{algebraic-definition-viewed-as}.

\medskip\noindent
Cas (1). Un champ algébrique est un espace algébrique si et seulement si sa
diagonale est un monomorphisme, voir
Morphismes de champs, Lemme \ref{stacks-morphisms-lemma-hierarchy}
(cela résulte aussi immédiatement de Champs algébriques,
Proposition \ref{algebraic-proposition-algebraic-stack-no-automorphisms}).

\medskip\noindent
Cas (2). D'après (1), on peut supposer que $\mathcal{X}$ et $\mathcal{X}'$
sont des espaces algébriques, puis appliquer
Compléments sur les morphismes d'espaces, Lemme
\ref{spaces-more-morphisms-lemma-thickening-scheme}.

\medskip\noindent
Cas (3) -- (6). Chacun de ces cas correspond à une condition
sur la diagonale, voir Morphismes de champs, Définitions
\ref{stacks-morphisms-definition-separated} et
\ref{stacks-morphisms-definition-absolute-separated}.
\end{proof}









\section{Morphismes d'épaississements}
\label{section-morphisms-thickenings}

\noindent
Si $(f, f') : (\mathcal{X} \subset \mathcal{X}') \to
(\mathcal{Y} \subset \mathcal{Y}')$ est un morphisme
d'épaississements de champs algébriques, alors les propriétés du morphisme
$f$ se transmettent souvent à $f'$. Il existe plusieurs variantes.

\begin{lemma}
\label{lemma-thicken-property-morphisms}
Soit $(f, f') : (\mathcal{X} \subset \mathcal{X}') \to
(\mathcal{Y} \subset \mathcal{Y}')$
un morphisme d'épaississements de champs algébriques. Alors
\begin{enumerate}
\item $f$ est un morphisme affine si et seulement si $f'$ est un morphisme affine,
\item $f$ est un morphisme surjectif si et seulement si $f'$ est un morphisme surjectif,
\item $f$ est quasi-compact si et seulement si $f'$ est quasi-compact,
\item $f$ est universellement fermé si et seulement si $f'$ est universellement fermé,
\item $f$ est entier si et seulement si $f'$ est entier,
\item $f$ est universellement injectif si et seulement si $f'$ est universellement injectif,
\item $f$ est universellement ouvert si et seulement si $f'$ est universellement ouvert,
\item $f$ est quasi-DM si et seulement si $f'$ est quasi-DM,
\item $f$ est DM si et seulement si $f'$ est DM,
\item $f$ est (quasi-)séparé si et seulement si $f'$ est (quasi-)séparé,
\item $f$ est représentable si et seulement si $f'$ est représentable,
\item $f$ est représentable par des espaces algébriques si et seulement si $f'$ est
représentable par des espaces algébriques,
\item en ajouter ici.
\end{enumerate}
\end{lemma}

\begin{proof}
D'après le Lemme \ref{lemma-thickening}, les morphismes $\mathcal{X} \to \mathcal{X}'$
et $\mathcal{Y} \to \mathcal{Y}'$ sont des homéomorphismes universels. Ainsi, toute
condition portant sur $|f| : |\mathcal{X}| \to |\mathcal{Y}|$ équivaut à
la condition correspondante sur $|f'| : |\mathcal{X}'| \to |\mathcal{Y}'|$,
et il en va de même après un changement de base arbitraire par un morphisme
$\mathcal{Z}' \to \mathcal{Y}'$. Cela démontre les
énoncés (2), (3), (4), (6) et (7).

\medskip\noindent
Dans les cas (8), (9), (10) et (12), on peut traduire les conditions sur
$f$ et $f'$ en conditions sur les diagonales $\Delta$ et $\Delta'$
comme dans le Lemme \ref{lemma-thickening-diagonals}. Voir
Morphismes de champs, Définition \ref{stacks-morphisms-definition-separated} et
Lemme \ref{stacks-morphisms-lemma-hierarchy}.
Ces cas résultent donc du Lemme \ref{lemma-thickening-diagonals}.

\medskip\noindent
Démonstration de (11). Si $f'$ est représentable, alors $f$ l'est aussi, car
pour un schéma $T$ et un morphisme $T \to \mathcal{Y}$, on a
$\mathcal{X} \times_\mathcal{Y} T =
\mathcal{X} \times_{\mathcal{X}'} (\mathcal{X}' \times_{\mathcal{Y}'} T)$
et $\mathcal{X} \to \mathcal{X}'$ est une immersion fermée (donc représentable).
Réciproquement, supposons $f$ représentable et soit $T' \to \mathcal{Y}'$
un morphisme, où $T'$ est un schéma. Alors
$$
\mathcal{X} \times_{\mathcal{Y}}
(\mathcal{Y} \times_{\mathcal{Y}'} T') =
\mathcal{X} \times_{\mathcal{X}'}
(\mathcal{X}' \times_{\mathcal{Y}'} T') \to
\mathcal{X}' \times_{\mathcal{Y}'} T'
$$
est un épaississement (d'après le Lemme \ref{lemma-base-change-thickening})
et la source est un schéma. La cible est donc un schéma d'après le
Lemme \ref{lemma-thickening-properties}.

\medskip\noindent
Dans les cas (1) et (5), si $f$ ou $f'$ possède la propriété indiquée,
alors $f$ et $f'$ sont tous deux représentables d'après (11). Dans ce cas,
choisissons un espace algébrique $V'$ et un morphisme lisse surjectif
$V' \to \mathcal{Y}'$. Posons $V = \mathcal{Y} \times_{\mathcal{Y}'} V'$,
$U' = \mathcal{X}' \times_{\mathcal{Y}'} V'$ et
$U = \mathcal{X} \times_{\mathcal{Y}'} V'$. Les résultats cherchés
découlent alors des résultats correspondants pour
le morphisme d'épaississements $(U \subset U') \to (V \subset V')$
d'espaces algébriques, au moyen du principe de Propriétés des champs, Lemme
\ref{stacks-properties-lemma-check-property-covering}.
Voir Compléments sur les morphismes d'espaces, Lemme
\ref{spaces-more-morphisms-lemma-thicken-property-morphisms}
pour les résultats correspondants dans le cas des espaces algébriques.
\end{proof}





\section{Déformations infinitésimales des champs algébriques}
\label{section-deform}

\noindent
Cette section est l'analogue de
Compléments sur les morphismes d'espaces, Section
\ref{spaces-more-morphisms-section-deform}.

\begin{lemma}
\label{lemma-flatness-morphism-thickenings-fp-over-ft}
Considérons un diagramme commutatif
$$
\xymatrix{
(\mathcal{X} \subset \mathcal{X}') \ar[rr]_{(f, f')} \ar[rd] & &
(\mathcal{Y} \subset \mathcal{Y}') \ar[ld] \\
& (\mathcal{B} \subset \mathcal{B}')
}
$$
d'épaississements de champs algébriques. Supposons que
\begin{enumerate}
\item $\mathcal{Y}' \to \mathcal{B}'$ soit localement de type fini,
\item $\mathcal{X}' \to \mathcal{B}'$ soit
plat et localement de présentation finie,
\item $f$ soit plat, et
\item $\mathcal{X} = \mathcal{B} \times_{\mathcal{B}'} \mathcal{X}'$ et
$\mathcal{Y} = \mathcal{B} \times_{\mathcal{B}'} \mathcal{Y}'$.
\end{enumerate}
Alors $f'$ est plat et, pour tout $y' \in |\mathcal{Y}'|$ appartenant à l'image de $|f'|$,
le morphisme $\mathcal{Y}' \to \mathcal{B}'$ est plat en $y'$.
\end{lemma}

\begin{proof}
Choisissons un espace algébrique $U'$ et un morphisme lisse surjectif
$U' \to \mathcal{B}'$.
Choisissons un espace algébrique $V'$ et un morphisme lisse surjectif
$V' \to U' \times_{\mathcal{B}'} \mathcal{Y}'$.
Choisissons un espace algébrique $W'$ et un morphisme lisse surjectif
$W' \to V' \times_{\mathcal{Y}'} \mathcal{X}'$. Soient $U, V, W$
les changements de base de $U', V', W'$ par $\mathcal{B} \to \mathcal{B}'$.
Alors la platitude de $f'$ équivaut à celle de $W' \to V'$, et
on sait que $W \to V$ est plat. On peut donc appliquer le lemme
dans le cas des espaces algébriques au diagramme
$$
\xymatrix{
(W \subset W') \ar[rr] \ar[rd] & & (V \subset V') \ar[ld] \\
& (U \subset U')
}
$$
d'épaississements d'espaces algébriques. Voir
Compléments sur les morphismes d'espaces, Lemme
\ref{spaces-more-morphisms-lemma-flatness-morphism-thickenings-fp-over-ft}.
L'assertion concernant la platitude de $\mathcal{Y}'/\mathcal{B}'$ aux points de
l'image de $|f'|$ s'en déduit de la même manière.
\end{proof}

\begin{lemma}
\label{lemma-deform-property-fp-over-ft}
Considérons un diagramme commutatif
$$
\xymatrix{
(\mathcal{X} \subset \mathcal{X}') \ar[rr]_{(f, f')} \ar[rd] & &
(\mathcal{Y} \subset \mathcal{Y}') \ar[ld] \\
& (\mathcal{B} \subset \mathcal{B}')
}
$$
d'épaississements de champs algébriques.
Supposons que $\mathcal{Y}' \to \mathcal{B}'$ soit localement de type fini,
que $\mathcal{X}' \to \mathcal{B}'$ soit plat et localement de présentation finie,
que $\mathcal{X} = \mathcal{B} \times_{\mathcal{B}'} \mathcal{X}'$ et
que $\mathcal{Y} = \mathcal{B} \times_{\mathcal{B}'} \mathcal{Y}'$. Alors
\begin{enumerate}
\item $f$ est plat si et seulement si $f'$ est plat,
\label{item-flat-fp-over-ft}
\item $f$ est un isomorphisme si et seulement si $f'$ est un isomorphisme,
\label{item-isomorphism-fp-over-ft}
\item $f$ est une immersion ouverte si et seulement si $f'$ est une immersion ouverte,
\label{item-open-immersion-fp-over-ft}
\item $f$ est un monomorphisme si et seulement si $f'$ est un monomorphisme,
\label{item-monomorphism-fp-over-ft}
\item $f$ est localement quasi-fini si et seulement si $f'$ est localement quasi-fini,
\label{item-quasi-finite-fp-over-ft}
\item $f$ est syntomique si et seulement si $f'$ est syntomique,
\label{item-syntomic-fp-over-ft}
\item $f$ est lisse si et seulement si $f'$ est lisse,
\label{item-smooth-fp-over-ft}
\item $f$ est non ramifié si et seulement si $f'$ est non ramifié,
\label{item-unramified-fp-over-ft}
\item $f$ est \'etale si et seulement si $f'$ est \'etale,
\label{item-etale-fp-over-ft}
\item $f$ est fini si et seulement si $f'$ est fini, et
\label{item-finite-fp-over-ft}
\item en ajouter ici.
\end{enumerate}
\end{lemma}

\begin{proof}
Dans le cas (\ref{item-flat-fp-over-ft}), cela résulte du
Lemme \ref{lemma-flatness-morphism-thickenings-fp-over-ft}.

\medskip\noindent
Dans les cas
(\ref{item-syntomic-fp-over-ft}), (\ref{item-smooth-fp-over-ft})
cela se démontre par la méthode employée dans la démonstration du
Lemme \ref{lemma-flatness-morphism-thickenings-fp-over-ft}.
Plus précisément, choisissons un espace algébrique $U'$ et un morphisme lisse surjectif
$U' \to \mathcal{B}'$.
Choisissons un espace algébrique $V'$ et un morphisme lisse surjectif
$V' \to U' \times_{\mathcal{B}'} \mathcal{Y}'$.
Choisissons un espace algébrique $W'$ et un morphisme lisse surjectif
$W' \to V' \times_{\mathcal{Y}'} \mathcal{X}'$. Soient $U, V, W$
les changements de base de $U', V', W'$ par $\mathcal{B} \to \mathcal{B}'$.
Alors la propriété pour $f$, resp.\ $f'$,
équivaut à la propriété pour $W \to V$, resp.\ $W' \to V'$.
On peut donc appliquer le lemme, dans le cas des espaces algébriques, au
diagramme
$$
\xymatrix{
(W \subset W') \ar[rr] \ar[rd] & & (V \subset V') \ar[ld] \\
& (U \subset U')
}
$$
d'épaississements d'espaces algébriques. Voir
Compléments sur les morphismes d'espaces, Lemme
\ref{spaces-more-morphisms-lemma-deform-property-fp-over-ft}.

\medskip\noindent
Dans les cas (\ref{item-unramified-fp-over-ft}) et (\ref{item-etale-fp-over-ft}),
on constate d'abord que l'hypothèse portant sur $f$ ou $f'$ implique que
$f$ et $f'$ sont tous deux des morphismes DM de champs algébriques, voir
le Lemme \ref{lemma-thicken-property-morphisms}. On peut alors choisir
un espace algébrique $U'$ et un morphisme lisse surjectif
$U' \to \mathcal{B}'$.
Choisissons un espace algébrique $V'$ et un morphisme lisse surjectif
$V' \to U' \times_{\mathcal{B}'} \mathcal{Y}'$.
Choisissons un espace algébrique $W'$ et un morphisme \'etale surjectif (!)
$W' \to V' \times_{\mathcal{Y}'} \mathcal{X}'$. Soient $U, V, W$
les changements de base de $U', V', W'$ par $\mathcal{B} \to \mathcal{B}'$.
Alors $W \to V \times_\mathcal{Y} \mathcal{X}$ est lui aussi
\'etale et surjectif. La propriété pour $f$, resp.\ $f'$,
équivaut donc à la propriété pour $W \to V$, resp.\ $W' \to V'$.
On peut donc appliquer le lemme, dans le cas des espaces algébriques, au
diagramme
$$
\xymatrix{
(W \subset W') \ar[rr] \ar[rd] & & (V \subset V') \ar[ld] \\
& (U \subset U')
}
$$
d'épaississements d'espaces algébriques. Voir
Compléments sur les morphismes d'espaces, Lemme
\ref{spaces-more-morphisms-lemma-deform-property-fp-over-ft}.

\medskip\noindent
Dans les cas (\ref{item-isomorphism-fp-over-ft}),
(\ref{item-open-immersion-fp-over-ft}),
(\ref{item-monomorphism-fp-over-ft}),
(\ref{item-finite-fp-over-ft})
on déduit d'abord du Lemme \ref{lemma-thicken-property-morphisms}
que $f$ et $f'$ sont représentables par des espaces algébriques. On peut donc choisir
un espace algébrique $U'$ et un morphisme lisse surjectif
$U' \to \mathcal{B}'$,
un espace algébrique $V'$ et un morphisme lisse surjectif
$V' \to U' \times_{\mathcal{B}'} \mathcal{Y}'$, puis
$W' = V' \times_{\mathcal{Y}'} \mathcal{X}'$ sera un espace algébrique.
Soient $U, V, W$ les changements de base de
$U', V', W'$ par $\mathcal{B} \to \mathcal{B}'$.
On a également $W = V \times_\mathcal{Y} \mathcal{X}$.
Il faut alors montrer que $W' \to V'$ est un
isomorphisme, resp.\ une immersion ouverte, resp.\ un monomorphisme,
resp.\ un morphisme fini, si et seulement si $W \to V$ possède la même propriété.
Voir Propriétés des champs, Lemme
\ref{stacks-properties-lemma-check-property-covering}.
On conclut ainsi en appliquant les résultats
précédents relatifs aux espaces algébriques.

\medskip\noindent
Dans le cas (\ref{item-quasi-finite-fp-over-ft}), observons d'abord
que $f$ et $f'$ sont localement de type fini d'après
Morphismes de champs, Lemme \ref{stacks-morphisms-lemma-finite-type-permanence}.
D'autre part, $f$ est quasi-DM si et seulement si
$f'$ l'est, d'après le
Lemme \ref{lemma-thicken-property-morphisms}.
La dernière condition à vérifier pour que $f$ ou $f'$ soit localement quasi-fini
(Morphismes de champs, Définition
\ref{stacks-morphisms-definition-locally-quasi-finite})
est une condition portant sur les espaces topologiques sous-jacents;
elle est satisfaite par $f$ si et seulement si elle l'est par $f'$ d'après
la discussion du premier paragraphe de la démonstration.
\end{proof}









\section{Relèvement des schémas affines}
\label{section-lifting-affines}

\noindent
Considérons un diagramme aux flèches pleines
$$
\xymatrix{
W \ar[d] \ar@{..>}[r] & W' \ar@{..>}[d] \\
\mathcal{X} \ar[r] & \mathcal{X}'
}
$$
où $\mathcal{X} \subset \mathcal{X}'$ est un épaississement de champs algébriques,
$W$ est un schéma affine et $W \to \mathcal{X}$ est lisse. La question
étudiée dans cette section est de savoir si l'on peut trouver $W'$ et les flèches
en pointillé de sorte que le carré soit cartésien et que $W' \to \mathcal{X}'$ soit lisse.
Nous n'en connaissons pas la réponse en général, mais si $\mathcal{X} \subset \mathcal{X}'$
est un épaississement du premier ordre, nous démontrerons que la réponse est affirmative.

\medskip\noindent
Pour étudier ce problème, introduisons la catégorie suivante.

\begin{remark}[Catégorie des relèvements]
\label{remark-gerbe-of-lifts}
Considérons un diagramme
$$
\xymatrix{
W \ar[d]_x \\
\mathcal{X} \ar[r] & \mathcal{X}'
}
$$
où $\mathcal{X} \subset \mathcal{X}'$ est un épaississement de champs algébriques,
$W$ est un espace algébrique et $W \to \mathcal{X}$ est lisse.
Nous allons construire une catégorie $\mathcal{C}$ et un foncteur
$$
p : \mathcal{C} \longrightarrow W_{spaces, \etale}
$$
(voir Propriétés des espaces, Définition
\ref{spaces-properties-definition-spaces-etale-site} pour la notation)
comme suit. Un objet de $\mathcal{C}$ sera un système
$(U, U', a, i, x', \alpha)$
qui forme un diagramme commutatif
\begin{equation}
\label{equation-object}
\vcenter{
\xymatrix{
U \ar[d]_a \ar[r]_i & U' \ar[dd]^{x'} \\
W \ar[d]_x & \\
\mathcal{X} \ar[r] & \mathcal{X}'
}
}
\end{equation}
dont la commutativité est matérialisée par le $2$-morphisme
$\alpha : x \circ a \to x' \circ i$, et tel que
$U$ et $U'$ soient des espaces algébriques,
$a : U \to W$ soit \'etale, $x' : U' \to \mathcal{X}'$ soit lisse
et $U = \mathcal{X} \times_{\mathcal{X}'} U'$.
En particulier, $U \subset U'$ est un épaississement.
Un morphisme
$$
(V, V', b, j, y', \beta) \to (U, U', a, i, x', \alpha)
$$
est la donnée de $(f, f', \gamma)$, où $f : V \to U$ est un morphisme
au-dessus de $W$, $f' : V' \to U'$ est un morphisme dont la restriction
à $V$ donne $f$, et $\gamma : x' \circ f' \to y'$ est un $2$-morphisme
qui matérialise la commutativité dans le triangle de droite du diagramme ci-dessous
\begin{equation}
\label{equation-morphism}
\vcenter{
\xymatrix{
& V \ar[ld]_f \ar[ldd]^b \ar[rr]_j & & V' \ar[ld]_{f'} \ar[lddd]^{y'} \\
U \ar[d]_a \ar[rr]_i & & U' \ar[dd]_{x'} \\
W \ar[d]_x & \\
\mathcal{X} \ar[rr] & & \mathcal{X}'
}
}
\end{equation}
Enfin, on exige que $\gamma$ soit compatible avec $\alpha$ et $\beta$ :
dans le calcul des $2$-catégories de Catégories, Sections
\ref{categories-section-formal-cat-cat} et
\ref{categories-section-2-categories}, cela s'écrit
$$
\beta = (\gamma \star \text{id}_j) \circ (\alpha \star \text{id}_f)
$$
(plus brièvement : $\beta = j^*\gamma \circ f^*\alpha$).
Une autre formulation consiste à dire que les objets sont les diagrammes commutatifs
(\ref{equation-object}) satisfaisant certaines propriétés supplémentaires, et que les
morphismes sont les diagrammes commutatifs
(\ref{equation-morphism}) dans la catégorie $\textit{Spaces}/\mathcal{X}'$
introduite dans Propriétés des champs, Remarque
\ref{stacks-properties-remark-representable-over}.
Il est alors clair que $\mathcal{C}$ est une catégorie
et que la règle $p : \mathcal{C} \to W_{spaces, \etale}$ qui
envoie $(U, U', a, i, x', \alpha)$ sur $a : U \to W$
est un foncteur.
\end{remark}

\begin{lemma}
\label{lemma-morphisms-lifts-etale}
Pour tout morphisme (\ref{equation-morphism}), l'application $f' : V' \to U'$ est \'etale.
\end{lemma}

\begin{proof}
En effet, $f : V \to U$ est \'etale en tant que morphisme de $W_{spaces, \etale}$,
et l'on peut appliquer le
Lemme \ref{lemma-deform-property-fp-over-ft}, puisque $U' \to \mathcal{X}'$
et $V' \to \mathcal{X}'$ sont lisses et que
$U = \mathcal{X} \times_{\mathcal{X}'} U'$ et
$V = \mathcal{X} \times_{\mathcal{X}'} V'$.
\end{proof}

\begin{lemma}
\label{lemma-gerbe-of-lifts-fibred}
La catégorie $p : \mathcal{C} \to W_{spaces, \etale}$ construite
dans la Remarque \ref{remark-gerbe-of-lifts} est fibrée en groupoïdes.
\end{lemma}

\begin{proof}
Nous affirmons que les catégories fibres de $p$ sont des groupoïdes.
Si $(f, f', \gamma')$, comme dans (\ref{equation-morphism}),
est un morphisme tel que $f : V \to U$ soit un isomorphisme, alors
$f'$ est un isomorphisme d'après le Lemme \ref{lemma-deform-property-fp-over-ft},
et $(f, f', \gamma')$ est donc un isomorphisme.

\medskip\noindent
Considérons un morphisme $f : V \to U$ dans $W_{spaces, \etale}$
et un objet $\xi = (U, U', a, i, x', \alpha)$
de $\mathcal{C}$ au-dessus de $U$. Nous allons construire l'« image inverse »
$f^*\xi$ au-dessus de $V$.
Posons $b = a \circ f$. Soit $f' : V' \to U'$ le morphisme \'etale
dont la restriction à $V$ est $f$ (Compléments sur les morphismes d'espaces,
Lemme \ref{spaces-more-morphisms-lemma-topological-invariance}).
Notons $j : V \to V'$ l'épaississement correspondant.
Soient $y' = x' \circ f'$ et $\gamma = \text{id} : x' \circ f' \to y'$.
Posons
$$
\beta = \alpha \star \text{id}_f :
x \circ b = x \circ a \circ f \to
x' \circ i \circ f = x' \circ f' \circ j = y' \circ j
$$
Il est clair que $(f, f', \gamma) : (V, V', b, j, y', \beta) \to
(U, U', a, i, x', \alpha)$ est un morphisme comme dans
(\ref{equation-morphism}). Les morphismes $(f, f', \gamma)$
ainsi construits sont fortement cartésiens
(Catégories, Définition \ref{categories-definition-cartesian-over-C}).
Nous omettons la démonstration détaillée; essentiellement, la raison est que,
étant donné un morphisme
$(g, g', \epsilon) : (Y, Y', c, k, z', \delta) \to (U, U', a, i, x', \alpha)$
dans $\mathcal{C}$ tel que $g$ se factorise sous la forme $g = f \circ h$ pour un certain
$h : Y \to V$, on obtient une unique factorisation $g' = f' \circ h'$
d'après Compléments sur les morphismes d'espaces,
Lemme \ref{spaces-more-morphisms-lemma-topological-invariance}
puis on peut produire le $\zeta$ nécessaire tel que
$(h, h', \zeta) :  (Y, Y', c, k, z', \delta) \to
(V, V', b, j, y', \beta)$ soit un morphisme de $\mathcal{C}$
vérifiant $(g, g', \epsilon) = (f, f', \gamma) \circ (h, h', \zeta)$.

\medskip\noindent
Par conséquent, $p : \mathcal{C} \to W_\etale$ est une catégorie
fibrée (Catégories, Définition \ref{categories-definition-fibred-category}).
En combinant cela au fait établi ci-dessus que les catégories fibres sont des groupoïdes,
on conclut que $p : \mathcal{C} \to W_\etale$ est
fibrée en groupoïdes d'après Catégories, Lemme
\ref{categories-lemma-fibred-groupoids}.
\end{proof}

\begin{lemma}
\label{lemma-gerbe-of-lifts-stack}
La catégorie $p : \mathcal{C} \to W_{spaces, \etale}$ construite
dans la Remarque \ref{remark-gerbe-of-lifts} est un champ en groupoïdes.
\end{lemma}

\begin{proof}
D'après le Lemme \ref{lemma-gerbe-of-lifts-fibred}, la première condition
de Champs, Définition \ref{stacks-definition-stack-in-groupoids}, est satisfaite.
Comme il est d'usage, nous vérifions la descente des objets et laissons au lecteur
le soin de vérifier celle des morphismes. Supposons donc donnés $a : U \to W$
dans $W_{spaces, \etale}$, un recouvrement $\{U_k \to U\}_{k \in K}$ dans
$W_{spaces, \etale}$, des objets $\xi_k = (U_k, U'_k, a_k, i_k, x'_k, \alpha_k)$
de $\mathcal{C}$ au-dessus de $U_k$, et des morphismes
$$
\varphi_{kk'} = (f_{kk'}, f'_{kk'}, \gamma_{kk'}) :
\xi_k|_{U_k \times_U U_{k'}} \to
\xi_{k'}|_{U_k \times_U U_{k'}}
$$
entre les restrictions, satisfaisant à la condition de cocycle. Pour démontrer
l'effectivité, on peut d'abord raffiner le recouvrement. On peut donc supposer que chaque
$U_k$ est un schéma (même un schéma affine si l'on veut). Écrivons
$$
\xi_k|_{U_k \times_U U_{k'}} =
(U_k \times_U U_{k'}, U'_{kk'}, a_{kk'}, x'_{kk'}, \alpha_{kk'})
$$
On obtient alors un morphisme \'etale (d'après le Lemme \ref{lemma-morphisms-lifts-etale})
$s_{kk'} : U'_{kk'} \to U'_k$
comme seconde composante du morphisme
$\xi_k|_{U_k \times_U U_{k'}} \to \xi_k$ de $\mathcal{C}$.
De même, on obtient un morphisme \'etale $t_{kk'} : U'_{kk'} \to U'_{k'}$
en considérant la seconde composante du composé
$$
\xi_k|_{U_k \times_U U_{k'}} \xrightarrow{\varphi_{kk'}}
\xi_{k'}|_{U_k \times_U U_{k'}} \to \xi_{k'}
$$
Nous affirmons que
$$
j :
\coprod\nolimits_{(k, k') \in K \times K} U'_{kk'}
\xrightarrow{(\coprod s_{kk'}, \coprod t_{kk'})}
(\coprod\nolimits_{k \in K} U'_k) \times (\coprod\nolimits_{k \in K} U'_k)
$$
est une relation d'équivalence \'etale. Premièrement, nous avons déjà vu
que les composantes $s,t$ du morphisme affiché sont \'etales.
Le changement de base du morphisme $j$ par
$(\coprod U_k) \times (\coprod U_k) \to (\coprod U'_k) \times (\coprod U'_k)$
est un monomorphisme, car c'est l'application
$$
\coprod\nolimits_{(k, k') \in K \times K} U_k \times_U U_{k'}
\longrightarrow
(\coprod\nolimits_{k \in K} U_k) \times (\coprod\nolimits_{k \in K} U_k)
$$
Par conséquent, $j$ est un monomorphisme d'après Compléments sur les morphismes, Lemme
\ref{more-morphisms-lemma-properties-that-extend-over-thickenings}.
Enfin, la symétrie de la relation $j$ provient du fait que
$\varphi_{kk'}^{-1}$ est le « renversé » de $\varphi_{k'k}$ (voir
Champs, Remarques \ref{stacks-remarks-definition-descent-datum}),
et la transitivité résulte de la condition de cocycle (détails omis).
Le quotient de $\coprod U'_k$ par $j$ est donc un espace algébrique $U'$
(Espaces, Théorème \ref{spaces-theorem-presentation}).
Nous avons déjà montré ci-dessus qu'il existe un épaississement
$i : U \to U'$, puisque la restriction de $j$ à
$\coprod U_k$ donne $(\coprod U_k) \times_U (\coprod U_k)$.
Enfin, si l'on considère temporairement les $1$-morphismes
$x'_k : U'_k \to \mathcal{X}'$ comme des objets du champ
$\mathcal{X}'$ au-dessus de $U'_k$, on voit qu'ils sont munis d'une
donnée de descente relativement au recouvrement \'etale
$\{U'_k \to U'\}$, donnée par la troisième composante $\gamma_{kk'}$
des morphismes $\varphi_{kk'}$ dans $\mathcal{C}$.
Puisque $\mathcal{X}'$ est un champ,
cette donnée de descente est effective; en revenant à la formulation initiale, on obtient
un $1$-morphisme $x' : U' \to \mathcal{X}'$ tel que les composés
$U'_k \to U' \to \mathcal{X}'$ soient munis d'isomorphismes vers $x'_k$
compatibles avec $\gamma_{kk'}$. Cela signifie que les morphismes
$\alpha_k : x \circ a_k \to x'_k \circ i_k$ se recollent en un morphisme
$\alpha : x \circ a \to x' \circ i$. Alors $\xi = (U, U', a, i, x', \alpha)$
est l'objet cherché au-dessus de $U$.
\end{proof}

\begin{lemma}
\label{lemma-etale-local-lifts}
Soit $\mathcal{X} \subset \mathcal{X}'$ un épaississement de champs algébriques.
Soit $W$ un espace algébrique et soit $W \to \mathcal{X}$ un morphisme lisse.
Il existe un recouvrement \'etale $\{W_i \to W\}_{i \in I}$ et, pour tout $i$,
un diagramme cartésien
$$
\xymatrix{
W_i \ar[r] \ar[d] & W_i' \ar[d] \\
\mathcal{X} \ar[r] & \mathcal{X}'
}
$$
où $W_i' \to \mathcal{X}'$ est lisse.
\end{lemma}

\begin{proof}
Choisissons un schéma $U'$ et un morphisme lisse surjectif $U' \to \mathcal{X}'$.
Comme d'habitude, posons $U = \mathcal{X} \times_{\mathcal{X}'} U'$. Alors
$U \to \mathcal{X}$ est un morphisme lisse surjectif. Par conséquent, le changement de base
$$
V = W \times_{\mathcal{X}} U \longrightarrow W
$$
est un morphisme lisse surjectif d'espaces algébriques.
D'après Topologies sur les espaces, Lemme
\ref{spaces-topologies-lemma-etale-dominates-smooth}
on peut trouver un recouvrement \'etale $\{W_i \to W\}$ tel
que $W_i \to W$ se factorise par $V \to W$.
Après avoir recouvert $W_i$ par des ouverts affines (Propriétés des espaces, Lemme
\ref{spaces-properties-lemma-cover-by-union-affines})
on peut supposer chaque $W_i$ affine. On peut, et l'on va, remplacer $W$ par $W_i$,
ce qui ramène à la situation étudiée au paragraphe suivant.

\medskip\noindent
Supposons $W$ affine et le morphisme donné $W \to \mathcal{X}$ factorisé
par $U$. On a le diagramme
$$
W \xrightarrow{i} U \to \mathcal{X}
$$
Puisque $W$ et $U$ sont lisses sur $\mathcal{X}$, on voit que
$i$ est localement de type fini (Morphismes de champs, Lemme
\ref{stacks-morphisms-lemma-finite-type-permanence}).
Après avoir remplacé $U$ par $\mathbf{A}^n_U$, on peut supposer
que $i$ est une immersion, voir
Morphismes, Lemme \ref{morphisms-lemma-quasi-affine-finite-type-over-S}.
D'après Morphismes de champs,
Lemme \ref{stacks-morphisms-lemma-lci-permanence}
le morphisme $i$ est d'intersection complète locale.
Par conséquent, $i$ est une immersion Koszul-régulière (au sens de
Diviseurs, Définition \ref{divisors-definition-regular-immersion}),
d'après Compléments sur les morphismes, Lemme \ref{more-morphisms-lemma-lci}.

\medskip\noindent
On peut encore remplacer $W$ par un recouvrement ouvert affine.
Pour tout point $w \in W$, on peut choisir un ouvert affine
$U'_w \subset U'$ tel que, si $U_w \subset U$ est
l'ouvert affine correspondant, alors $w \in i^{-1}(U_w)$ et
$i^{-1}(U_w) \to U_w$ est une immersion fermée définie
par une suite Koszul-régulière
$f_1, \ldots, f_r \in \Gamma(U_w, \mathcal{O}_{U_w})$.
Cela résulte de la définition des immersions Koszul-régulières
et de Diviseurs, Lemme \ref{divisors-lemma-regular-ideal-sheaf-scheme}.
Posons $W_w = i^{-1}(U_w)$; c'est un voisinage ouvert affine
de $w \in W$.
Choisissons des relèvements $f'_1, \ldots, f'_r \in \Gamma(U'_w, \mathcal{O}_{U'_w})$
de $f_1, \ldots, f_r$. Cela est possible puisque $U_w \to U'_w$
est une immersion fermée de schémas affines.
Soit $W'_w \subset U'_w$ le sous-schéma fermé défini par
$f'_1, \ldots, f'_r$.
Nous affirmons que $W'_w \to \mathcal{X}'$ est lisse.
Cette affirmation achève la démonstration, car
$W_w = \mathcal{X} \times_{\mathcal{X}'} W'_w$
par construction.

\medskip\noindent
Pour vérifier l'affirmation, il suffit de vérifier que le changement de base
$W'_w \times_{\mathcal{X}'} X' \to X'$ est lisse pour tout
schéma affine $X'$ lisse sur $\mathcal{X}'$. Choisissons un
morphisme \'etale
$$
Y' \to U'_w \times_{\mathcal{X}'} X'
$$
avec $Y'$ affine. Puisque $U'_w \times_{\mathcal{X}'} X'$ est recouvert
par les images de tels morphismes, il suffit de montrer que le sous-schéma fermé
$Z'$ de $Y'$ défini par $f'_1, \ldots, f'_r$ est lisse sur $X'$.
Considérons le diagramme
$$
\xymatrix{
Z' \ar[r] \ar[d] & Y' \ar[d] \\
W'_w \times_{\mathcal{X}'} X' \ar[d] \ar[r] &
U'_w \times_{\mathcal{X}'} X' \ar[d] \ar[r] & X' \\
W'_w = V(f'_1, \ldots, f'_r) \ar[r] & U'_w
}
$$
Posons $X = \mathcal{X} \times_{\mathcal{X}'} X'$,
$Y = X \times_{X'} Y' = \mathcal{X} \times_{\mathcal{X}'} Y'$ et
$Z = Y \times_{Y'} Z' = X \times_{X'} Z' =
\mathcal{X} \times_{\mathcal{X}'} Z'$.
Alors $(Z \subset Z') \to (Y \subset Y') \subset (X \subset X')$
sont des morphismes (cartésiens) d'épaississements de schémas affines, et
on sait que $Z \to X$ et $Y' \to X'$ sont lisses.
Enfin, la suite de fonctions $f'_1, \ldots, f'_r$
a pour image une suite Koszul-régulière dans $\Gamma(Y', \mathcal{O}_{Y'})$ d'après
Compléments d'algèbre, Lemme \ref{more-algebra-lemma-koszul-regular-flat-base-change},
car $Y' \to U'_w$ est lisse, donc plat.
D'après Compléments d'algèbre, Lemme \ref{more-algebra-lemma-cut-by-koszul}
(et le fait que les suites Koszul-régulières sont des suites quasi-régulières
d'après Compléments d'algèbre, Lemmes \ref{more-algebra-lemma-regular-koszul-regular},
\ref{more-algebra-lemma-koszul-regular-H1-regular} et
\ref{more-algebra-lemma-H1-regular-quasi-regular})
on conclut que $Z' \to X'$ est lisse, comme voulu.
\end{proof}

\begin{lemma}
\label{lemma-etale-local-lifts-isomorphic}
Soit $\mathcal{X} \subset \mathcal{X}'$ un épaississement de champs algébriques.
Considérons un diagramme commutatif
$$
\xymatrix{
W'' \ar[d]_{x''} & W \ar[l] \ar[r] \ar[d]_x & W' \ar[d]^{x'} \\
\mathcal{X}' & \mathcal{X} \ar[l] \ar[r] & \mathcal{X}'
}
$$
dont les carrés sont cartésiens, où $W', W, W''$ sont des espaces algébriques et
les flèches verticales sont lisses. Alors il existe
\begin{enumerate}
\item un recouvrement \'etale $\{f'_k : W'_k \to W'\}_{k \in K}$,
\item des morphismes \'etales $f''_k : W'_k \to W''$, et
\item des $2$-morphismes $\gamma_k : x' \circ f'_k \to x'' \circ f''_k$
\end{enumerate}
tels que (a) $(f'_k)^{-1}(W) = (f''_k)^{-1}(W)$, (b)
$f'_k|_{(f'_k)^{-1}(W)} = f''_k|_{(f''_k)^{-1}(W)}$ et
(c) l'image inverse de $\gamma_k$ sur le sous-schéma fermé de (a)
coïncide avec le $2$-morphisme donné par la commutativité du
diagramme initial au-dessus de $W$.
\end{lemma}

\begin{proof}
Notons $i : W \to W'$ et $i'' : W \to W''$ les épaississements donnés.
La commutativité du diagramme de l'énoncé du lemme
signifie qu'il existe un $2$-morphisme $\delta : x' \circ i \to x'' \circ i''$.
C'est le $2$-morphisme auquel renvoie la partie (c) de l'énoncé.
Considérons l'espace algébrique
$$
I' = W' \times_{x', \mathcal{X}', x''} W''
$$
muni des projections $p' : I' \to W'$ et $q' : I' \to W''$.
Observons qu'il existe un $2$-morphisme « universel »
$\gamma : x' \circ p' \to x'' \circ q'$ (nous l'utiliserons plus loin).
Le choix de $\delta$ définit un morphisme
$$
\xymatrix{
W \ar[rr]_\delta & & I' \ar[ld]^{p'} \ar[rd]_{q'} \\
& W' & & W''
}
$$
tel que les composés $W \to I' \to W'$ et $W \to I' \to W''$
soient respectivement $i : W \to W'$ et $i'' : W \to W''$.
Puisque $x''$ est lisse, le morphisme $p' : I' \to W'$ est lisse
en tant que changement de base de $x''$.

\medskip\noindent
Supposons que l'on puisse trouver un recouvrement \'etale $\{f'_k : W'_k \to W'\}$
et des morphismes $\delta_k : W'_k \to I'$ tels que la restriction
de $\delta_k$ à $W_k = (f'_k)^{-1}(W)$ soit égale à $\delta \circ f_k$,
où $f_k = f'_k|_{W_k}$. On a le diagramme
$$
\xymatrix{
W_k \ar[r]^{f_k} \ar[d] & W \ar[r]^\delta & I' \ar[d]^{p'} \\
W'_k \ar[rr]^{f'_k} \ar[rru]^{\delta_k} & & W'
}
$$
Autrement dit, on veut pouvoir prolonger la section donnée
$\delta : W \to I'$ de $p'$ en une section au-dessus de $W'$, après avoir éventuellement
remplacé $W'$ par un recouvrement \'etale.

\medskip\noindent
Si cela est possible, on peut poser $f''_k = q' \circ \delta_k$
et $\gamma_k = \gamma \star \text{id}_{\delta_k}$ (plus brièvement,
$\gamma_k = \delta_k^*\gamma$). Il reste seulement à montrer
que le morphisme $f''_k$ est \'etale.
Par construction, le morphisme $x' \circ p'$ est $2$-isomorphe
à $x'' \circ q'$. Par conséquent, $x'' \circ f''_k$ est $2$-isomorphe
à $x' \circ f'_k$. On en déduit que le composé
$$
W'_k \xrightarrow{f''_k} W'' \xrightarrow{x''} \mathcal{X}'
$$
est lisse, puisque $x' \circ f'_k$ l'est.
Puisque $f_k$ est \'etale, on conclut que $f''_k$ est \'etale
d'après le Lemme \ref{lemma-deform-property-fp-over-ft}.

\medskip\noindent
Si l'épaississement est du premier ordre, on peut
choisir un recouvrement \'etale quelconque $\{W'_k \to W'\}$ avec $W'_k$ affine.
En effet, puisque $p'$ est lisse, on voit que $p'$ est formellement lisse d'après le
critère infinitésimal de relèvement (Compléments sur les morphismes d'espaces, Lemme
\ref{spaces-more-morphisms-lemma-smooth-formally-smooth}).
Comme $W_k$ est affine et que $W_k \to W'_k$ est un épaississement du premier ordre
(en tant que changement de base de $\mathcal{X} \to \mathcal{X}'$, voir
le Lemme \ref{lemma-base-change-thickening}), on obtient le $\delta_k$
cherché.

\medskip\noindent
Dans le cas général, l'existence du recouvrement et des morphismes
$\delta_k$ résulte de Compléments sur les morphismes d'espaces, Lemme
\ref{spaces-more-morphisms-lemma-smooth-strong-lift}.
\end{proof}

\begin{lemma}
\label{lemma-gerbe-of-lifts}
La catégorie $p : \mathcal{C} \to W_{spaces, \etale}$ construite
dans la Remarque \ref{remark-gerbe-of-lifts} est une gerbe.
\end{lemma}

\begin{proof}
Dans le Lemme \ref{lemma-gerbe-of-lifts-stack},
nous avons vu qu'il s'agit d'un champ en groupoïdes.
Il reste donc à vérifier les conditions (2) et (3) de
Champs, Définition \ref{stacks-definition-gerbe}.
La condition (2) résulte du
Lemme \ref{lemma-etale-local-lifts}.
La condition (3) résulte du
Lemme \ref{lemma-etale-local-lifts-isomorphic}.
\end{proof}

\begin{lemma}
\label{lemma-gerbe-of-lifts-first-order}
Dans la Remarque \ref{remark-gerbe-of-lifts}, supposons que $\mathcal{X} \subset \mathcal{X}'$
soit un épaississement du premier ordre. Alors
\begin{enumerate}
\item les faisceaux d'automorphismes des objets de la gerbe
$p : \mathcal{C} \to W_{spaces, \etale}$ construite
dans la Remarque \ref{remark-gerbe-of-lifts} sont abéliens, et
\item le faisceau de groupes $\mathcal{G}$ construit dans
Champs, Lemme \ref{stacks-lemma-gerbe-abelian-auts},
est un $\mathcal{O}_W$-module quasi-cohérent.
\end{enumerate}
\end{lemma}

\begin{proof}
Nous allons démontrer simultanément les deux assertions. Plus précisément, étant donné
un objet $\xi = (U, U', a, i, x', \alpha)$, nous munirons
$\mathit{Aut}(\xi)$ d'une structure de
$\mathcal{O}_U$-module quasi-cohérent sur $U_{spaces, \etale}$,
et nous montrerons que cette structure est compatible aux images inverses.
Cela suffira par recollement des faisceaux
(Sites, Section \ref{sites-section-glueing-sheaves}),
et par la construction de $\mathcal{G}$ dans la démonstration de
Champs, Lemme \ref{stacks-lemma-gerbe-abelian-auts},
comme recollement des faisceaux d'automorphismes $\mathit{Aut}(\xi)$,
ainsi que par le fait qu'il suffit, pour vérifier qu'un module est
quasi-cohérent, de passer à un recouvrement \'etale
(Propriétés des espaces, Lemme
\ref{spaces-properties-lemma-characterize-quasi-coherent}).

\medskip\noindent
Nous allons décrire le faisceau $\mathit{Aut}(\xi)$ par la
même méthode que dans la démonstration du
Lemme \ref{lemma-etale-local-lifts-isomorphic}.
Considérons l'espace algébrique
$$
I' = U' \times_{x', \mathcal{X}', x'} U'
$$
muni des projections $p' : I' \to U'$ et $q' : I' \to U'$.
Sur $I'$, il existe un $2$-morphisme universel
$\gamma : x' \circ p' \to x' \circ q'$.
L'identité $x' \to x'$ définit un morphisme diagonal
$$
\xymatrix{
U' \ar[rr]_{\Delta'} & & I' \ar[ld]^{p'} \ar[rd]_{q'} \\
& U' & & U'
}
$$
tel que les composés $U' \to I' \to U'$ et $U' \to I' \to U'$
soient les morphismes identité. Nous noterons les changements de base de
$U', I', p', q', \Delta'$ à $\mathcal{X}$ par $U, I, p, q, \Delta$.
Puisque $U' \to \mathcal{X}'$ est lisse, on voit que $p' : I' \to U'$
est lisse par changement de base.

\medskip\noindent
Une section de $\mathit{Aut}(\xi)$ au-dessus de $U$ est un morphisme $\delta' : U' \to I'$
tel que $\delta'|_U = \Delta$ et
$p' \circ \delta' = \text{id}_{U'}$. Explicitement,
$(\text{id}_U, q' \circ \delta', (\delta')^*\gamma) : \xi \to \xi$
est une formule pour l'automorphisme correspondant.
Plus généralement, si
$f : V \to U$ est un morphisme \'etale, alors il existe un épaississement
$j : V \to V'$ et un morphisme \'etale
$f' : V' \to U'$ dont la restriction à $V$ est $f$, et $f^*\xi$
correspond à $(V, V', a \circ f, j, x' \circ f', f^*\alpha)$, voir la démonstration du
Lemme \ref{lemma-gerbe-of-lifts-fibred}.
Une section de $\mathit{Aut}(\xi)$ au-dessus de $V$ est un morphisme
$\delta' : V' \to I'$
tel que $\delta'|_V = \Delta \circ f$
et $p' \circ \delta' = f'$\footnote{Une formule pour l'automorphisme
correspondant est $(\text{id}_V, h', (\delta')^*\gamma)$.
Ici, $h' : V' \to V'$ est l'unique (iso)morphisme tel que
$h'|_V = \text{id}_V$ et que le diagramme
$$
\xymatrix{
V' \ar[r]_{h'} \ar[rd]_{q' \circ \delta'} & V' \ar[d]^{f'} \\
& U'
}
$$
commute. L'unicité et l'existence de $h'$ résultent de l'invariance topologique
du site \'etale, voir Compléments sur les morphismes d'espaces,
Théorème \ref{spaces-more-morphisms-theorem-topological-invariance}.
Le lecteur pourrait estimer qu'il faudrait plutôt considérer les morphismes
$\delta'' : V' \to V' \times_{\mathcal{X}'} V'$ tels que
$\delta'' \circ j = \Delta_{V'/\mathcal{X}'}$ et
$\text{pr}_1 \circ \delta'' = \text{id}_{V'}$.
Cela conviendrait également : puisque $V' \times_{\mathcal{X}'} V' \to I'$ est \'etale,
la même invariance topologique montre que l'application qui envoie
$\delta''$ sur $\delta' = (V' \times_{\mathcal{X}'} V' \to I') \circ \delta''$
est une bijection entre les deux
ensembles de morphismes.}.

\medskip\noindent
On conclut que le faisceau d'ensembles $\mathit{Aut}(\xi)$ coïncide avec
le faisceau défini dans
Compléments sur les morphismes d'espaces, Remarque
\ref{spaces-more-morphisms-remark-another-special-case}
pour les épaississements $(U \subset U')$ et $(I \subset I')$ sur
$(U \subset U')$, via $\text{id}_{U'}$ et $p'$.
La diagonale $\Delta'$ est une section de ce faisceau et, en faisant
agir sur cette section Compléments sur les morphismes d'espaces, Lemme
\ref{spaces-more-morphisms-lemma-action-sheaf}
on obtient un isomorphisme
\begin{equation}
\label{equation-isomorphism}
\SheafHom_{\mathcal{O}_U}(\Delta^*\Omega_{I/U}, \mathcal{C}_{U/U'})
\longrightarrow
\mathit{Aut}(\xi)
\end{equation}
sur $U_{spaces, \etale}$. Il reste trois choses à vérifier :
\begin{enumerate}
\item la construction de (\ref{equation-isomorphism})
commute à la localisation \'etale,
\item $\SheafHom_{\mathcal{O}_U}(\Delta^*\Omega_{I/U}, \mathcal{C}_{U/U'})$
est un module quasi-cohérent sur $U$,
\item la composition dans $\mathit{Aut}(\xi)$ correspond
à l'addition des sections de ce module quasi-cohérent.
\end{enumerate}
Nous allons les vérifier dans cet ordre.

\medskip\noindent
Pour démontrer (1), il faut montrer que, si $f : V \to U$ est \'etale,
alors (\ref{equation-isomorphism}), construit à partir de $\xi$ sur $U$,
se restreint à l'application (\ref{equation-isomorphism})
$$
\SheafHom_{\mathcal{O}_V}(
\Delta_V^*\Omega_{V \times_\mathcal{X} V/V}, \mathcal{C}_{V/V'}) \to
\mathit{Aut}(\xi|_V)
$$
construite à partir de $\xi|_V$ sur $V$, dans $V_{spaces, \etale}$.
Cela résulte de la discussion dans la note ci-dessus
et de Compléments sur les morphismes d'espaces, Lemme
\ref{spaces-more-morphisms-lemma-action-by-derivations-etale-localization}.

\medskip\noindent
Démonstration de (2). Puisque $p'$ est lisse, le morphisme $I \to U$ est lisse,
et le module des différentielles relatives $\Omega_{I/U}$ est donc
localement libre de type fini (Compléments sur les morphismes d'espaces, Lemme
\ref{spaces-more-morphisms-lemma-smooth-omega-finite-locally-free}).
D'autre part, $\mathcal{C}_{U/U'}$ est quasi-cohérent
(Compléments sur les morphismes d'espaces, Définition
\ref{spaces-more-morphisms-definition-conormal-sheaf}).
D'après Propriétés des espaces, Lemme
\ref{spaces-properties-lemma-properties-quasi-coherent}
on conclut.

\medskip\noindent
Démonstration de (3). Il existe un morphisme $c' : I' \times_{p', U', q'} I' \to I'$
tel que $(U', I', p', q', c')$ soit un groupoïde en espaces algébriques
d'identité $\Delta'$. Voir par exemple
Champs algébriques, Lemme \ref{algebraic-lemma-map-space-into-stack}.
La composition dans $\mathit{Aut}(\xi)$ est induite par le morphisme
$c'$ comme suit. Supposons donnés deux morphismes
$$
\delta'_1, \delta'_2 : U' \longrightarrow I'
$$
correspondant, comme ci-dessus, à des sections de $\mathit{Aut}(\xi)$ au-dessus de $U$;
autrement dit, on a $\delta'_i|_U = \Delta_U$ et
$p' \circ \delta'_i = \text{id}_{U'}$. Alors la composition
dans $\mathit{Aut}(\xi)$ est
$$
\delta'_1 \circ \delta'_2 = c'(\delta'_1 \circ q' \circ \delta'_2, \delta'_2)
$$
Nous omettons la vérification détaillée\footnote{Le lecteur voit immédiatement
qu'il est nécessaire de précomposer
$\delta'_1$ par $q' \circ \delta'_2$ pour obtenir un point à valeurs dans $U'$
bien défini du produit fibré $I' \times_{p', U', q'} I'$.}.
Nous sommes donc dans la situation décrite dans
Compléments sur les groupoïdes en espaces, Section
\ref{spaces-more-groupoids-section-groupoid-sections}
et le résultat cherché découle du
Lemme de Compléments sur les groupoïdes en espaces
\ref{spaces-more-groupoids-lemma-composition-is-addition}.
\end{proof}

\begin{proposition}[Emerton]
\label{proposition-affine-smooth-lift-to-first-order}
\begin{reference}
Courriel de Matthew Emerton daté du 27 avril 2016.
\end{reference}
Soit $\mathcal{X} \subset \mathcal{X}'$ un épaississement du premier ordre
de champs algébriques. Soit $W$ un schéma affine et soit
$W \to \mathcal{X}$ un morphisme lisse. Alors il existe
un diagramme cartésien
$$
\xymatrix{
W \ar[d] \ar[r] & W' \ar[d] \\
\mathcal{X} \ar[r] & \mathcal{X}'
}
$$
où $W' \to \mathcal{X}'$ est lisse et $W'$ affine.
\end{proposition}

\begin{proof}
Considérons la catégorie $p : \mathcal{C} \to W_{spaces, \etale}$
introduite dans la Remarque \ref{remark-gerbe-of-lifts}.
La proposition affirme qu'il existe un objet de $\mathcal{C}$
au-dessus de $W$. En effet, si l'on dispose d'un tel objet
$(W, W', a, i, y', \alpha)$, alors $W = \mathcal{X} \times_{\mathcal{X}'} W'$.
Par conséquent, $W \to W'$ est un épaississement d'espaces algébriques, donc
$W'$ est affine d'après
Compléments sur les morphismes d'espaces, Lemme
\ref{spaces-more-morphisms-lemma-thickening-scheme}
et Compléments sur les morphismes, Lemme
\ref{more-morphisms-lemma-thickening-affine-scheme}.

\medskip\noindent
Le Lemme \ref{lemma-gerbe-of-lifts} nous dit que $\mathcal{C}$ est une gerbe sur
$W_{spaces, \etale}$. Cela signifie que l'on peut trouver des solutions localement pour la topologie \'etale,
et que ces solutions locales sont localement isomorphes pour la topologie \'etale;
cette partie ne requiert pas que l'épaississement soit du premier ordre.
D'après le Lemme \ref{lemma-gerbe-of-lifts-first-order},
les faisceaux d'automorphismes des objets de notre gerbe sont abéliens et
se recollent pour former un module quasi-cohérent $\mathcal{G}$
sur $W_{spaces, \etale}$. Nous allons vérifier les conditions (1) et (2)
de Cohomologie sur les sites, Lemme \ref{sites-cohomology-lemma-existence},
afin de conclure à l'existence d'un objet de $\mathcal{C}$ au-dessus de $W$.
La condition (1) est satisfaite : les recouvrements \'etales $\{W_i \to W\}$
dont chaque $W_i$ est affine sont cofinaux dans l'ensemble de tous les recouvrements.
Pour un tel recouvrement, $W_i$ et $W_i \times_W W_j$ sont affines,
et $H^1(W_i, \mathcal{G})$ ainsi que $H^1(W_i \times_W W_j, \mathcal{G})$
sont nuls : la cohomologie d'un module quasi-cohérent sur un
espace algébrique affine est nulle, par exemple d'après Cohomologie des espaces, Proposition
\ref{spaces-cohomology-proposition-vanishing}.
Enfin, la condition (2) est que $H^2(W, \mathcal{G}) = 0$
pour notre faisceau quasi-cohérent $\mathcal{G}$, ce qui résulte encore
de Cohomologie des espaces, Proposition
\ref{spaces-cohomology-proposition-vanishing}.
Cela achève la démonstration.
\end{proof}






\section{Déformations infinitésimales}
\label{section-inf}

\noindent
Nous poursuivons la discussion de
Les axiomes d'Artin, Section \ref{artin-section-inf}.

\begin{lemma}
\label{lemma-inf-quasi-coherent}
Soit $\mathcal{X}$ un champ algébrique sur un schéma $S$.
Supposons que $\mathcal{I}_\mathcal{X} \to \mathcal{X}$ soit localement
de présentation finie. Soit $A \to B$ un homomorphisme plat de $S$-algèbres.
Soit $x$ un objet de $\mathcal{X}$ sur $A$, et posons $y = x|_B$.
Alors $\text{Inf}_x(M) \otimes_A B = \text{Inf}_y(M \otimes_A B)$.
\end{lemma}

\begin{proof}
Rappelons que $\text{Inf}_x(M)$ est l'ensemble des automorphismes de la
déformation triviale de $x$ à $A[M]$ qui induisent l'automorphisme
identique de $x$ sur $A$. La déformation triviale est
l'image inverse de $x$ sur $\Spec(A[M])$ par $\Spec(A[M]) \to \Spec(A)$.
Soit $G \to \Spec(A)$ l'espace algébrique groupe des automorphismes de $x$
(celui-ci existe puisque $\mathcal{X}$ est un champ algébrique).
Soit $e : \Spec(A) \to G$ l'élément neutre.
La discussion de Compléments sur les morphismes d'espaces, Section
\ref{spaces-more-morphisms-section-action-by-derivations}
donne
$$
\text{Inf}_x(M) = \Hom_A(e^*\Omega_{G/A}, M)
$$
De même,
$$
\text{Inf}_y(M \otimes_A B) = \Hom_B(e_B^*\Omega_{G_B/B}, M \otimes_A B)
$$
Puisque $G \to \Spec(A)$ est localement de présentation finie par
hypothèse, on voit que $\Omega_{G/A}$ est localement de
présentation finie, voir
Compléments sur les morphismes d'espaces, Lemme
\ref{spaces-more-morphisms-lemma-finite-presentation-differentials}.
Par conséquent, $e^*\Omega_{G/A}$ est un $A$-module de présentation finie.
De plus, $\Omega_{G_B/B}$ est l'image inverse de $\Omega_{G/A}$ d'après
Compléments sur les morphismes d'espaces, Lemme
\ref{spaces-more-morphisms-lemma-base-change-differentials}.
Ainsi, $e_B^*\Omega_{G_B/B} = e^*\Omega_{G/A} \otimes_A B$.
On conclut par Compléments d'algèbre, Lemme
\ref{more-algebra-lemma-pseudo-coherence-and-base-change-ext}.
\end{proof}

\begin{lemma}
\label{lemma-sheaf-of-infinitesimal-lifts}
Soit $\mathcal{X}$ un champ algébrique sur un schéma de base $S$.
Supposons que $\mathcal{I}_\mathcal{X} \to \mathcal{X}$ soit localement
de présentation finie. Soit $(A' \to A, x)$ une situation de déformation.
Alors le foncteur
$$
F : B' \longmapsto
\{\text{relèvements de }x|_{B' \otimes_{A'} A}\text{ à } B'\}/\text{isomorphismes}
$$
est un faisceau sur le site $(\textit{Aff}/\Spec(A'))_{fppf}$ de
Topologies, Définition \ref{topologies-definition-big-small-fppf}.
\end{lemma}

\begin{proof}
Soit $\{T'_i \to T'\}_{i = 1, \ldots, n}$ un recouvrement fppf standard
de schémas affines sur $A'$. Écrivons $T' = \Spec(B')$. Comme d'habitude, notons
$$
T'_{i_0 \ldots i_p} =
T'_{i_0} \times_{T'} \ldots \times_{T'} T'_{i_p} = \Spec(B'_{i_0 \ldots i_p})
$$
où l'anneau est un produit tensoriel convenable.
Posons $B = B' \otimes_{A'} A$ et
$B_{i_0 \ldots i_p} = B'_{i_0 \ldots i_p} \otimes_{A'} A$.
Notons $y = x|_B$ et $y_{i_0 \ldots i_p} = x|_{B_{i_0 \ldots i_p}}$.
Soient $\gamma_i \in F(B'_i)$ tels que $\gamma_{i_0}$ et $\gamma_{i_1}$
s'envoient sur le même élément de $F(B'_{i_0i_1})$.
Il faut trouver un unique $\gamma \in F(B')$ dont l'image soit
$\gamma_i$ dans $F(B'_i)$.

\medskip\noindent
Choisissons un objet véritable $y'_i$ de $\textit{Lift}(y_i, B'_i)$
dans la classe d'isomorphisme $\gamma_i$.
Choisissons des isomorphismes
$\varphi_{i_0i_1} : y'_{i_0}|_{B'_{i_0i_1}} \to y'_{i_1}|_{B'_{i_0i_1}}$
dans la catégorie $\textit{Lift}(y_{i_0i_1}, B'_{i_0i_1})$.
Si les applications $\varphi_{i_0i_1}$ satisfont à la condition de cocycle,
on obtient l'objet $\gamma$ cherché, puisque $\mathcal{X}$ est
un champ pour la topologie fppf. La condition de cocycle affirme que le
composé
$$
y'_{i_0}|_{B'_{i_0i_1i_2}}
\xrightarrow{\varphi_{i_0i_1}|_{B'_{i_0i_1i_2}}}
y'_{i_1}|_{B'_{i_0i_1i_2}}
\xrightarrow{\varphi_{i_1i_2}|_{B'_{i_0i_1i_2}}}
y'_{i_2}|_{B'_{i_0i_1i_2}}
\xrightarrow{\varphi_{i_2i_0}|_{B'_{i_0i_1i_2}}}
y'_{i_0}|_{B'_{i_0i_1i_2}}
$$
est l'identité. Sinon, ces applications donnent des éléments
$$
\delta_{i_0i_1i_2} \in
\text{Inf}_{y_{i_0i_1i_2}}(J_{i_0i_1i_2}) =
\text{Inf}_y(J) \otimes_B B_{i_0i_1i_2}
$$
Ici, $J = \Ker(B' \to B)$ et
$J_{i_0 \ldots i_p} = \Ker(B'_{i_0 \ldots i_p} \to B_{i_0 \ldots i_p})$.
L'égalité de la formule affichée résulte du
Lemme \ref{lemma-inf-quasi-coherent} appliqué à
$B' \to B'_{i_0 \ldots i_p}$, à $y$ et à $y_{i_0 \ldots i_p}$,
ainsi que de la platitude des applications $B' \to B'_{i_0 \ldots i_p}$,
qui garantit également que
$J_{i_0 \ldots i_p} = J \otimes_{B'} B'_{i_0 \ldots i_p}$.
Un calcul (omis) montre que $\delta_{i_0i_1i_2}$ définit
un $2$-cocycle dans le complexe de {\v C}ech
$$
\prod \text{Inf}_y(J) \otimes_B B_{i_0} \to
\prod \text{Inf}_y(J) \otimes_B B_{i_0i_1} \to
\prod \text{Inf}_y(J) \otimes_B B_{i_0i_1i_2} \to \ldots
$$
D'après Descente, Lemme \ref{descent-lemma-standard-covering-Cech-quasi-coherent},
ce complexe est acyclique en degrés positifs et vérifie $H^0 = \text{Inf}_y(J)$.
Puisque $\text{Inf}_{y_{i_0i_1}}(J_{i_0i_1})$ agit sur les
morphismes (Les axiomes d'Artin, Remarque \ref{artin-remark-automorphisms}),
cela signifie que l'on peut modifier le choix de $\varphi_{i_0i_1}$
pour se ramener au cas où $\delta_{i_0i_1i_2} = 0$.

\medskip\noindent
Unicité. Il reste à montrer qu'il existe au plus un $\gamma$ dont la restriction
soit $\gamma_i$ pour tout $i$. Supposons donnés des objets $y', z'$ de
$\textit{Lift}(y, B')$ et des isomorphismes $\psi_i : y'|_{B'_i} \to z'|_{B'_i}$
dans $\textit{Lift}(y_i, B'_i)$. On peut alors considérer
$$
\psi_{i_1}^{-1} \circ \psi_{i_0} \in
\text{Inf}_{y_{i_0i_1}}(J_{i_0i_1}) =
\text{Inf}_y(J) \otimes_B B_{i_0i_1}
$$
Comme précédemment, l'obstruction à l'existence d'un isomorphisme entre
$y'$ et $z'$ sur $B'$ est un élément du $H^1$ du
complexe de {\v C}ech affiché ci-dessus, lequel est nul.
\end{proof}

\begin{lemma}
\label{lemma-T-quasi-coherent}
Soit $\mathcal{X}$ un champ algébrique sur un schéma $S$, dont le
morphisme structural $\mathcal{X} \to S$ est localement de présentation finie.
Soit $A \to B$ un homomorphisme plat de $S$-algèbres.
Soit $x$ un objet de $\mathcal{X}$ sur $A$.
Alors $T_x(M) \otimes_A B = T_y(M \otimes_A B)$.
\end{lemma}

\begin{proof}
Choisissons un schéma $U$ et un morphisme lisse surjectif $U \to \mathcal{X}$.
Réduisons d'abord le lemme au cas où $x$ se relève à $U$.
Rappelons que $T_x(M)$ est l'ensemble des classes d'isomorphisme des relèvements de $x$
à $A[M]$. Par conséquent, le
Lemme \ref{lemma-sheaf-of-infinitesimal-lifts}\footnote{Ce lemme s'applique :
$\Delta : \mathcal{X} \to \mathcal{X} \times_S \mathcal{X}$
est localement de présentation finie d'après
Morphismes de champs, Lemme
\ref{stacks-morphisms-lemma-finite-presentation-permanence}
et l'hypothèse que $\mathcal{X} \to S$ est localement de présentation finie.
Par conséquent, $\mathcal{I}_\mathcal{X} \to \mathcal{X}$ est localement
de présentation finie en tant que changement de base de $\Delta$.}
affirme que la règle
$$
A_1 \mapsto T_{x|_{A_1}}(M \otimes_A A_1)
$$
est un faisceau sur le petit site \'etale de $\Spec(A)$; le produit tensoriel
est nécessaire pour que $A[M] \to A_1[M \otimes_A A_1]$ soit un homomorphisme plat.
On peut choisir un homomorphisme d'anneaux \'etale fidèlement plat $A \to A_1$
tel que $x|_{A_1}$ se relève en un morphisme $u_1 : \Spec(A_1) \to U$, voir
par exemple Faisceaux sur les champs, Lemme
\ref{stacks-sheaves-lemma-surjective-flat-locally-finite-presentation}.
Écrivons $A_2 = A_1 \otimes_A A_1$ et posons $B_1 = B \otimes_A A_1$
et $B_2 = B \otimes_A A_2$. Considérons le diagramme
$$
\xymatrix{
0 \ar[r] &
T_y(M \otimes_A B) \ar[r] &
T_{y|_{B_1}}(M \otimes_A B_1) \ar[r] &
T_{y|_{B_2}}(M \otimes_A B_2) \\
0 \ar[r] &
T_x(M) \ar[r] \ar[u] &
T_{x|_{A_1}}(M \otimes_A A_1) \ar[r] \ar[u] &
T_{x|_{A_2}}(M \otimes_A A_2) \ar[u]
}
$$
Les lignes sont exactes par la condition de faisceau. On a
$M \otimes_A B_i = (M \otimes_A A_i) \otimes_{A_i} B_i$.
Ainsi, si l'on démontre le résultat pour les flèches verticales du milieu et de droite,
le résultat en découle. On est donc ramené au cas étudié au
paragraphe suivant.

\medskip\noindent
Supposons que $x$ soit l'image d'un morphisme $u : \Spec(A) \to U$.
Observons que $T_u(M) \to T_x(M)$ est surjectif, puisque
$U \to \mathcal{X}$ est lisse et représentable par des
espaces algébriques, voir Critères de représentabilité, Lemme
\ref{criteria-lemma-representable-by-spaces-formally-smooth}
(voir la discussion qui précède pour une explication) et
Compléments sur les morphismes d'espaces, Lemme
\ref{spaces-more-morphisms-lemma-smooth-formally-smooth}.
Posons $R = U \times_\mathcal{X} U$. Rappelons que l'on obtient
un groupoïde $(U, R, s, t, c, e, i)$ en espaces algébriques avec
$\mathcal{X} = [U/R]$. D'après
Les axiomes d'Artin, Lemme \ref{artin-lemma-ses-inf-and-T},
on a une suite exacte
$$
T_{e \circ u}(M) \to T_u(M) \oplus T_u(M) \to T_x(M) \to 0
$$
où la nullité du terme de droite a été démontrée ci-dessus. Une suite analogue
vaut après changement de base à $B$. Le résultat cherché découle donc de
celui du lemme pour
$T_u(M)$ et $T_{e \circ u}(M)$. On est ainsi ramené au cas étudié
dans le paragraphe suivant.

\medskip\noindent
Supposons que $\mathcal{X} = X$ soit un espace algébrique localement de
présentation finie sur $S$. Alors
$$
T_x(M) = \Hom_A(x^*\Omega_{X/S}, M)
$$
d'après la discussion de Compléments sur les morphismes d'espaces, Section
\ref{spaces-more-morphisms-section-action-by-derivations}.
De même,
$$
T_y(M \otimes_A B) = \Hom_B(y^*\Omega_{X/S}, M \otimes_A B)
$$
Puisque $X \to S$ est localement de présentation finie, on voit que
$\Omega_{X/S}$ est localement de présentation finie, voir
Compléments sur les morphismes d'espaces, Lemme
\ref{spaces-more-morphisms-lemma-finite-presentation-differentials}.
Ainsi, $x^*\Omega_{X/S}$ est un $A$-module de présentation finie.
Évidemment, on a $y^*\Omega_{X/S} = x^*\Omega_{X/S} \otimes_A B$.
On conclut par Compléments d'algèbre, Lemme
\ref{more-algebra-lemma-pseudo-coherence-and-base-change-ext}.
\end{proof}

\begin{lemma}
\label{lemma-local-lift-enough}
Soit $\mathcal{X}$ un champ algébrique sur un schéma $S$, dont le
morphisme structural $\mathcal{X} \to S$ est localement de présentation finie.
Soit $(A' \to A, x)$ une situation de déformation. S'il existe une
$A'$-algèbre $B'$ fidèlement plate et de présentation finie, ainsi qu'un
objet $y'$ de $\mathcal{X}$ sur $B'$ relevant $x|_{B' \otimes_{A'} A}$,
alors il existe un objet $x'$ sur $A'$ relevant $x$.
\end{lemma}

\begin{proof}
Soit $I = \Ker(A' \to A)$. Posons $B'_1 = B' \otimes_{A'} B'$ et
$B'_2 = B' \otimes_{A'} B' \otimes_{A'} B'$. Soient
$J = IB'$, $J_1 = IB'_1$, $J_2 = IB'_2$ et
$B = B'/J$, $B_1 = B'_1/J_1$, $B_2 = B'_2/J_2$.
Posons $y = x|_B$, $y_1 = x|_{B_1}$, $y_2 = x|_{B_2}$.
Soit $F$ le faisceau fppf du
Lemme \ref{lemma-sheaf-of-infinitesimal-lifts}
(qui s'applique, voir la note dans la démonstration du
Lemme \ref{lemma-T-quasi-coherent}).
On a donc un diagramme d'égalisateur
$$
\xymatrix{
F(A') \ar[r] &
F(B') \ar@<1ex>[r] \ar@<-1ex>[r] &
F(B'_1)
}
$$
D'autre part, on a $F(B') = \text{Lift}(y, B')$,
$F(B'_1) = \text{Lift}(y_1, B'_1)$ et $F(B'_2) = \text{Lift}(y_2, B'_2)$,
dans la terminologie de Les axiomes d'Artin, Section \ref{artin-section-inf}.
Ces ensembles sont non vides et sont (canoniquement) des espaces principaux homogènes sous
$T_y(J)$, $T_{y_1}(J_1)$ et $T_{y_2}(J_2)$, voir
Les axiomes d'Artin, Lemme \ref{artin-lemma-properties-lift-RS-star}.
La différence des deux images de $y'$ dans
$F(B'_1)$ est donc un élément
$$
\delta_1 \in T_{y_1}(J_1) = T_x(I) \otimes_A B_1
$$
L'égalité de la formule affichée résulte du
Lemme \ref{lemma-T-quasi-coherent} appliqué à
$A' \to B'_1$, à $x$ et à $y_1$, ainsi que de la platitude de $A' \to B'_1$,
qui garantit également que $J_1 = I \otimes_{A'} B'_1$. On a des
égalités analogues pour $B'$ et $B'_2$.
Un calcul (omis) montre que $\delta_1$ définit
un $1$-cocycle dans le complexe de {\v C}ech
$$
T_x(I) \otimes_A B \to
T_x(I) \otimes_A B_1 \to
T_x(I) \otimes_A B_2 \to \ldots
$$
D'après Descente, Lemme \ref{descent-lemma-standard-covering-Cech-quasi-coherent},
ce complexe est acyclique en degrés positifs et vérifie $H^0 = T_x(I)$.
On peut donc choisir un élément de $T_x(I) \otimes_A B = T_y(J)$
dont le cobord est $\delta_1$. En remplaçant $y'$ par le résultat
de l'action de cet élément, on obtient un nouveau choix $y'$ tel que $\delta_1 = 0$.
Ainsi, $y'$ s'envoie sur le même élément par les deux applications
$F(B') \to F(B'_1)$, et on obtient un élément de $F(A')$ par
la condition de faisceau.
\end{proof}












\section{Morphismes formellement lisses}
\label{section-formally-smooth}

\noindent
Dans cette section, nous introduisons la notion de morphisme formellement lisse
$\mathcal{X} \to \mathcal{Y}$ de champs algébriques. Un tel morphisme est
caractérisé par la propriété que les points de $\mathcal{X}$ à valeurs dans $T$ se relèvent
aux épaississements infinitésimaux de $T$, pourvu que $T$ soit affine.
Le résultat principal affirme qu'un morphisme formellement lisse et
localement de présentation finie est lisse, voir le
Lemme \ref{lemma-smooth-formally-smooth}.
Il s'avère que ce critère est souvent plus facile à utiliser que le
critère jacobien.

\begin{definition}
\label{definition-formally-smooth}
Un morphisme $f : \mathcal{X} \to \mathcal{Y}$ de champs algébriques est dit
{\it formellement lisse} s'il est formellement lisse sur les objets en tant que
$1$-morphisme de catégories fibrées en groupoïdes, comme expliqué dans
Critères de représentabilité, Section \ref{criteria-section-formally-smooth}.
\end{definition}

\noindent
Traduisons la condition de la définition dans le langage que nous employons
actuellement (voir
Propriétés des champs, Section \ref{stacks-properties-section-conventions}).
Soit $f : \mathcal{X} \to \mathcal{Y}$ un morphisme de champs algébriques.
Considérons un diagramme $2$-commutatif aux flèches pleines
\begin{equation}
\label{equation-diagram}
\vcenter{
\xymatrix{
T \ar[r]_-x \ar[d]_i & \mathcal{X} \ar[d]^f \\
T' \ar[r]^-y \ar@{..>}[ru] & \mathcal{Y}
}
}
\end{equation}
où $i : T \to T'$ est un épaississement du premier ordre de schémas affines.
Soit
$$
\gamma : y \circ i \longrightarrow f \circ x
$$
un $2$-morphisme matérialisant la $2$-commutativité du diagramme.
(Notation comme dans Catégories, Sections \ref{categories-section-formal-cat-cat}
et \ref{categories-section-2-categories}.)
Étant donnés (\ref{equation-diagram}) et $\gamma$,
une {\it flèche pointillée} est un triplet $(x', \alpha, \beta)$ formé d'un
morphisme $x' : T' \to \mathcal{X}$ et de $2$-flèches
$\alpha : x' \circ i \to x$, $\beta : y \to f \circ x'$
tels que
$\gamma = (\text{id}_f \star \alpha) \circ (\beta \star \text{id}_i)$,
autrement dit tel que
$$
\xymatrix{
& f \circ x' \circ i \ar[rd]^{\text{id}_f \star \alpha} \\
y \circ i \ar[ru]^{\beta \star \text{id}_i} \ar[rr]^\gamma & &
f \circ x
}
$$
soit commutatif. Un {\it morphisme de flèches pointillées}
$(x'_1, \alpha_1, \beta_1) \to (x'_2, \alpha_2, \beta_2)$ est une
$2$-flèche $\theta : x'_1 \to x'_2$ telle que
$\alpha_1 = \alpha_2 \circ (\theta \star \text{id}_i)$ et
$\beta_2 = (\text{id}_f \star \theta) \circ \beta_1$.

\medskip\noindent
La catégorie des flèches pointillées que l'on vient de décrire est un cas particulier de
Catégories, Définition \ref{categories-definition-dotted-arrows}.

\begin{lemma}
\label{lemma-reformulate-formal-smoothness}
Un morphisme $f : \mathcal{X} \to \mathcal{Y}$ de champs algébriques est
formellement lisse (Définition \ref{definition-formally-smooth})
si et seulement si, pour tout diagramme (\ref{equation-diagram}) et tout $\gamma$,
la catégorie des flèches pointillées est non vide.
\end{lemma}

\begin{proof}
La traduction entre les différents langages est omise.
\end{proof}

\begin{lemma}
\label{lemma-base-change-formal-smoothness}
Le changement de base d'un morphisme formellement lisse de champs algébriques
par un morphisme quelconque de champs algébriques est formellement lisse.
\end{lemma}

\begin{proof}
Cela résulte de
Catégories, Lemme \ref{categories-lemma-cat-dotted-arrows-base-change},
et de la définition.
\end{proof}

\begin{lemma}
\label{lemma-composition-formal-smoothness}
Le composé de morphismes formellement lisses de champs algébriques
est formellement lisse.
\end{lemma}

\begin{proof}
Cela résulte de
Catégories, Lemme \ref{categories-lemma-cat-dotted-arrows-composition},
et de la définition.
\end{proof}

\begin{lemma}
\label{lemma-formal-smoothness-representable}
Soit $f : \mathcal{X} \to \mathcal{Y}$ un morphisme de champs algébriques
représentable par des espaces algébriques. Les conditions suivantes sont équivalentes :
\begin{enumerate}
\item $f$ est formellement lisse,
\item pour tout schéma $T$ et tout morphisme $T \to \mathcal{Y}$,
le morphisme $\mathcal{X} \times_\mathcal{Y} T \to T$
est formellement lisse comme morphisme d'espaces algébriques.
\end{enumerate}
\end{lemma}

\begin{proof}
Cela résulte de
Catégories, Lemme \ref{categories-lemma-cat-dotted-arrows-base-change},
et de la définition.
\end{proof}

\begin{lemma}
\label{lemma-lift-to-smooth}
Soit $T \to T'$ un épaississement du premier ordre de schémas affines.
Soit $\mathcal{X}'$ un champ algébrique sur $T'$,
dont le morphisme structural $\mathcal{X}' \to T'$ est lisse.
Soit $x : T \to \mathcal{X}'$ un morphisme
sur $T'$. Alors il existe un morphisme $x' : T' \to \mathcal{X}'$
sur $T'$ tel que $x'|_T = x$.
\end{lemma}

\begin{proof}
On peut appliquer le résultat du Lemme \ref{lemma-local-lift-enough}.
Il suffit donc de construire un morphisme lisse surjectif
$W' \to T'$, avec $W'$ affine, tel que
$x|_{T \times_{T'} W'}$ se relève à $W'$.
(Nous invitons le lecteur à trouver sa propre démonstration de ce fait
en utilisant le résultat analogue déjà établi
pour les espaces algébriques.) Choisissons un
schéma $U'$ et un morphisme lisse surjectif $U' \to \mathcal{X}'$.
Observons que $U' \to T'$ est lisse et que la projection
$T \times_{\mathcal{X}'} U' \to T$ est lisse et surjective.
Choisissons un schéma affine $W$ et un morphisme \'etale
$W \to T \times_{\mathcal{X}'} U'$ tel que $W \to T$
soit surjectif. Alors $W \to T$ est un morphisme lisse de
schémas affines. Après avoir remplacé $W$ par une réunion disjointe d'ouverts
affines principaux, on peut supposer qu'il existe un
morphisme lisse de schémas affines $W' \to T'$ tel que
$W = T \times_{T'} W'$, voir Algèbre, Lemme \ref{algebra-lemma-lift-smooth}.
D'après Compléments sur les morphismes d'espaces, Lemme
\ref{spaces-more-morphisms-lemma-smooth-formally-smooth}
on peut trouver un morphisme $W' \to U'$ sur $T'$ relevant
le morphisme donné $W \to U'$. Cela achève la démonstration.
\end{proof}

\noindent
Le lemme suivant est le résultat principal de cette section.
Combiné à
Limites de champs, Proposition
\ref{stacks-limits-proposition-characterize-locally-finite-presentation},
il implique que l'on peut reconnaître si un morphisme de champs algébriques
$f : \mathcal{X} \to \mathcal{Y}$ est lisse en termes de
propriétés « simples » du $1$-morphisme de champs en groupoïdes
$\mathcal{X} \to \mathcal{Y}$.

\begin{lemma}[Critère infinitésimal de relèvement]
\label{lemma-smooth-formally-smooth}
Soit $f : \mathcal{X} \to \mathcal{Y}$ un morphisme de champs algébriques.
Les conditions suivantes sont équivalentes :
\begin{enumerate}
\item Le morphisme $f$ est lisse.
\item Le morphisme $f$ est localement de présentation finie et
formellement lisse.
\end{enumerate}
\end{lemma}

\begin{proof}
Supposons $f$ lisse. Alors $f$ est localement de présentation finie d'après
Morphismes de champs, Lemme
\ref{stacks-morphisms-lemma-smooth-locally-finite-presentation}.
Étant donné un diagramme (\ref{equation-diagram})
et un $\gamma : y \circ i \to f \circ x$, il suffit donc de trouver une
flèche pointillée (voir Lemme \ref{lemma-reformulate-formal-smoothness}).
En formant les produits fibrés, on obtient
$$
\xymatrix{
T \ar[d] \ar[r] &
T' \times_\mathcal{Y} \mathcal{X} \ar[d] \ar[r] &
\mathcal{X} \ar[d] \\
T' \ar[r] & T' \ar[r] & \mathcal{Y}
}
$$
Il suffit donc de trouver une flèche pointillée dans le
carré de gauche. Puisque $T' \times_\mathcal{Y} \mathcal{X} \to T'$
est lisse
(Morphismes de champs, Lemme \ref{stacks-morphisms-lemma-base-change-smooth}),
l'existence d'une flèche pointillée dans le carré de gauche est garantie par le
Lemme \ref{lemma-lift-to-smooth}.

\medskip\noindent
Réciproquement, supposons que $f$ soit localement de présentation finie et
formellement lisse. Choisissons un schéma $U$ et un morphisme lisse surjectif
$U \to \mathcal{X}$. Alors $a : U \to \mathcal{X}$ et $b : U \to \mathcal{Y}$
sont représentables par des espaces algébriques et localement de présentation finie
(utiliser Morphismes de champs, Lemme
\ref{stacks-morphisms-lemma-composition-finite-presentation}
et le fait constaté ci-dessus
qu'un morphisme lisse est localement de présentation finie).
Nous allons appliquer le principe général de
Champs algébriques, Lemme
\ref{algebraic-lemma-representable-transformations-property-implication}
en prenant comme donnée l'équivalence de Compléments sur les morphismes d'espaces,
Lemme \ref{spaces-more-morphisms-lemma-smooth-formally-smooth},
et utiliser simultanément la traduction du
Lemme de Critères de représentabilité
\ref{criteria-lemma-representable-by-spaces-formally-smooth}.
On applique d'abord ce principe à $a$ pour voir que $a$ est
formellement lisse sur les objets. Ensuite, on utilise que $f$ est
formellement lisse sur les objets par hypothèse
(voir Lemme \ref{lemma-reformulate-formal-smoothness})
ainsi que Critères de représentabilité, Lemme
\ref{criteria-lemma-composition-formally-smooth}
pour voir que $b = f \circ a$ est formellement lisse sur les objets.
On applique alors le principe une fois encore pour conclure que $b$ est lisse.
Cela signifie que $f$ est lisse par définition de la lissité
des morphismes de champs algébriques, ce qui achève la démonstration.
\end{proof}









\section{Éclatement et platitude}
\label{section-blowup-flat}

\noindent
Cette section expose rapidement ce que l'on peut déduire de
Compléments sur les morphismes d'espaces, Sections
\ref{spaces-more-morphisms-section-blowup-flat} et
\ref{spaces-more-morphisms-section-applications-flattening-by-blowing-up}
pour les champs algébriques sur des espaces algébriques.

\begin{lemma}
\label{lemma-flatten-stack}
Soit $f : \mathcal{X} \to Y$ un morphisme d'un champ algébrique
vers un espace algébrique. Soit $V \subset Y$ un sous-espace ouvert. Supposons que
\begin{enumerate}
\item $Y$ soit quasi-compact et quasi-séparé,
\item $f$ soit de type fini et quasi-séparé,
\item $V$ soit quasi-compact, et
\item $\mathcal{X}_V$ soit plat et localement de présentation finie sur $V$.
\end{enumerate}
Alors il existe un éclatement $V$-admissible $Y' \to Y$
et un sous-champ fermé $\mathcal{X}' \subset \mathcal{X}_{Y'}$
tel que $\mathcal{X}'_V = \mathcal{X}_V$ et que
$\mathcal{X}' \to Y'$ soit plat et de présentation finie.
\end{lemma}

\begin{proof}
Observons que $\mathcal{X}$ est quasi-compact.
Choisissons un schéma affine $U$ et un morphisme lisse
surjectif $U \to \mathcal{X}$.
Soit $R = U \times_\mathcal{X} U$; on obtient ainsi
un groupoïde $(U, R, s, t, c)$ en espaces algébriques sur $Y$, avec
$\mathcal{X} = [U/R]$
(Champs algébriques, Lemme \ref{algebraic-lemma-stack-presentation}).
On peut appliquer le
Lemme de Compléments sur les morphismes d'espaces
\ref{spaces-more-morphisms-lemma-flat-after-blowing-up}
à $U \to Y$ et à l'ouvert $V \subset Y$.
On obtient ainsi un éclatement $V$-admissible $Y' \to Y$
tel que la transformée stricte $U' \subset U_{Y'}$
soit plate et de présentation finie sur $Y'$.
Soit $R' \subset R_{Y'}$ la transformée stricte de $R$.
Puisque $s$ et $t$ sont lisses (et en particulier plats),
il résulte de
Diviseurs sur les espaces, Lemme
\ref{spaces-divisors-lemma-strict-transform-flat}
que l'on a des diagrammes cartésiens
$$
\vcenter{
\xymatrix{
R' \ar[r] \ar[d] & R_{Y'} \ar[d]^{s_{Y'}} \\
U' \ar[r] & U_{Y'}
}
}
\quad\text{et}\quad
\vcenter{
\xymatrix{
R' \ar[r] \ar[d] & R_{Y'} \ar[d]^{t_{Y'}} \\
U' \ar[r] & U_{Y'}
}
}
$$
Autrement dit, $U'$ est un sous-espace fermé invariant par $R_{Y'}$
de $U_{Y'}$. Ainsi, $U'$ définit un sous-champ fermé
$\mathcal{X}' \subset \mathcal{X}_{Y'}$ d'après
Propriétés des champs, Lemme
\ref{stacks-properties-lemma-substacks-presentation}.
Le morphisme $\mathcal{X}' \to Y'$ est plat et
localement de présentation finie, parce que c'est
le cas de $U' \to Y'$. D'autre part,
on sait déjà que $\mathcal{X}' \to Y'$ est quasi-compact et quasi-séparé
(par nos hypothèses sur $f$ et parce que c'est le cas des immersions fermées).
Cela achève la démonstration.
\end{proof}















\section{Lemme de Chow pour les champs algébriques}
\label{section-chow}

\noindent
Dans cette section, nous étudions le lemme de Chow pour les champs algébriques.

\begin{lemma}
\label{lemma-finite-cover-factor}
Soit $Y$ un espace algébrique quasi-compact et quasi-séparé.
Soit $V \subset Y$ un ouvert quasi-compact. Soit $f : \mathcal{X} \to V$
surjectif, plat et localement de présentation finie.
Alors il existe un morphisme fini surjectif $g : Y' \to Y$ tel que
$V' = g^{-1}(V) \to Y$ se factorise localement pour la topologie de Zariski par $f$.
\end{lemma}

\begin{proof}
Commençons par le démontrer lorsque $Y$ est un schéma.
On peut choisir un schéma $U$ et un morphisme lisse surjectif
$U \to \mathcal{X}$. Alors $\{U \to V\}$ est un recouvrement fppf de schémas.
D'après Compléments sur les morphismes, Lemme \ref{more-morphisms-lemma-dominate-fppf-finite},
il existe un morphisme fini surjectif
$V' \to V$ tel que $V' \to V$ se factorise localement pour la topologie de Zariski
par $U$. D'après
Compléments sur les morphismes, Lemme
\ref{more-morphisms-lemma-extend-finite-surjective-morphisms}
on peut trouver un morphisme fini surjectif $Y' \to Y$
dont la restriction à $V$ est $V' \to V$, comme souhaité.

\medskip\noindent
Si $Y$ est un espace algébrique, le lemme résulte d'abord d'un
changement de base fini par un morphisme fini surjectif
$Y' \to Y$, où $Y'$ est un schéma. Voir
Limites d'espaces, Proposition
\ref{spaces-limits-proposition-there-is-a-scheme-finite-over}.
\end{proof}

\begin{lemma}
\label{lemma-make-section}
Soit $f : \mathcal{X} \to Y$ un morphisme d'un champ algébrique
vers un espace algébrique. Soit $V \subset Y$ un sous-espace ouvert.
Supposons que
\begin{enumerate}
\item $f$ soit séparé et de type fini,
\item $Y$ soit quasi-compact et quasi-séparé,
\item $V$ soit quasi-compact, et
\item $\mathcal{X}_V$ soit une gerbe sur $V$.
\end{enumerate}
Alors il existe un diagramme commutatif
$$
\xymatrix{
\overline{Z} \ar[rd]_{\overline{g}} &
Z \ar[l]^j \ar[d]_g \ar[r]_h & \mathcal{X} \ar[ld]^f \\
& Y
}
$$
où $j$ est une immersion ouverte, $\overline{g}$ et $h$ sont propres,
et tel que $|V|$ soit contenu dans l'image de $|g|$.
\end{lemma}

\begin{proof}
Supposons donné un diagramme commutatif
$$
\xymatrix{
\mathcal{X}' \ar[d]_{f'} \ar[r] & \mathcal{X} \ar[d]^f \\
Y' \ar[r] & Y
}
$$
et un ouvert quasi-compact $V' \subset Y'$, tels que
$Y' \to Y$ soit un morphisme propre d'espaces algébriques,
$\mathcal{X}' \to \mathcal{X}$ soit un morphisme propre de champs algébriques,
$V' \subset Y'$ s'envoie surjectivement sur $V$, et
$\mathcal{X}'_{V'}$ soit une gerbe sur $V'$.
Il suffit alors de démontrer le lemme pour la paire
$(f' : \mathcal{X}' \to Y', V')$. Nous omettons certains détails.

\medskip\noindent
Stratégie générale de la démonstration. Nous nous ramènerons
au cas où l'image de $f$ est ouverte et où $f$
admet une section au-dessus de cet ouvert, en appliquant à plusieurs reprises
la remarque précédente. Chaque étape est directe, mais elles sont
assez nombreuses, ce qui rend la démonstration un peu longue.

\medskip\noindent
En utilisant Limites d'espaces, Proposition
\ref{spaces-limits-proposition-there-is-a-scheme-finite-over}
on se ramène au cas où $Y$ est un schéma.
(Soit $Y' \to Y$ un morphisme fini surjectif, où $Y'$ est
un schéma. Posons $\mathcal{X}' = \mathcal{X}_{Y'}$ et appliquons
la remarque initiale de la démonstration.)

\medskip\noindent
En utilisant le Lemme \ref{lemma-flatten-stack}
(et Morphismes de champs, Lemme \ref{stacks-morphisms-lemma-gerbe-fppf},
pour voir qu'une gerbe est plate et localement de présentation finie),
on se ramène au cas où $f$ est plat et de présentation finie.

\medskip\noindent
Puisque $f$ est plat et localement de
présentation finie, l'image de $|f|$ est un ouvert $W \subset Y$.
Puisque $\mathcal{X}$ est quasi-compact (car $f$ est de type fini
et $Y$ est quasi-compact), $W$ est quasi-compact.
D'après le Lemme \ref{lemma-finite-cover-factor},
on peut trouver un morphisme fini surjectif $g : Y' \to Y$
tel que $g^{-1}(W) \to Y$ se factorise localement pour la topologie de Zariski
par $\mathcal{X} \to Y$.
En remplaçant $Y$ par $Y'$ et $\mathcal{X}$ par
$\mathcal{X} \times_Y Y'$, on se ramène à la situation
décrite au paragraphe suivant.

\medskip\noindent
Supposons qu'il existe un entier $n \geq 0$, des ouverts quasi-compacts
$W_i \subset Y$, $i = 1, \ldots, n$, et des
morphismes $x_i : W_i \to \mathcal{X}$ tels que
(a) $f \circ x_i = \text{id}_{W_i}$,
(b) $W = \bigcup_{i = 1, \ldots, n} W_i$ contienne $V$, et
(c) $W$ soit l'image de $|f|$.
Nous raisonnons par récurrence sur $n$. Le cas initial est $n = 0$ : cela
implique $V = \emptyset$ et, dans ce cas, on peut prendre
$\overline{Z} = \emptyset$.
Si $n > 0$, considérons, pour $i = 1, \ldots, n$,
les sous-schémas fermés réduits
$Y_i$ dont l'espace topologique sous-jacent est
$Y \setminus W_i$. Considérons le morphisme fini
$$
Y' = Y \amalg \coprod\nolimits_{i = 1, \ldots, n} Y_i \longrightarrow Y
$$
et l'ouvert quasi-compact
$$
V' = (W_1 \cap \ldots \cap W_n \cap V) \amalg
\coprod_{i = 1, \ldots, n} (V \cap Y_i).
$$
D'après la remarque initiale de la démonstration, si l'on peut démontrer le lemme pour les paires
$$
(\mathcal{X} \to Y, W_1 \cap \ldots \cap W_n \cap V)
\quad\text{et}\quad
(\mathcal{X} \times_Y Y_i \to Y_i, V \cap Y_i),\quad
i = 1, \ldots, n
$$
alors le résultat est vrai. On utilise ici l'égalité ensembliste
$V = (W_1 \cap \ldots \cap W_n \cap V) \cup
\bigcup\nolimits_{i = 1, \ldots, n} (V \cap Y_i)$.
L'hypothèse de récurrence s'applique au second type de
paires ci-dessus. On se ramène donc à la situation décrite au
paragraphe suivant.

\medskip\noindent
Supposons qu'il existe un entier $n \geq 0$, des ouverts quasi-compacts
$W_i \subset Y$, $i = 1, \ldots, n$, et des
morphismes $x_i : W_i \to \mathcal{X}$ tels que
(a) $f \circ x_i = \text{id}_{W_i}$,
(b) $W = \bigcup_{i = 1, \ldots, n} W_i$ contienne $V$,
(c) $W$ soit l'image de $|f|$, et
(d) $V \subset W_1 \cap \ldots \cap W_n$.
Les morphismes
$$
T_{ij} = \mathit{Isom}_\mathcal{X}(x_i|_{W_i \cap W_j \cap V},
x_j|_{W_i \cap W_j \cap V}) \longrightarrow W_i \cap W_j \cap V
$$
sont surjectifs, plats et localement de présentation finie
(Morphismes de champs, Lemme \ref{stacks-morphisms-lemma-gerbe-isom-fppf}).
Appliquons le Lemme \ref{lemma-finite-cover-factor}
à chacun des ouverts quasi-compacts $W_i \cap W_j \cap V$ et aux morphismes
$T_{ij} \to W_i \cap W_j \cap V$; on obtient des morphismes finis surjectifs
$Y'_{ij} \to Y$. En remplaçant $Y$ par le produit fibré de tous
les $Y'_{ij}$ au-dessus de $Y$, on se ramène à la situation décrite
au paragraphe suivant.

\medskip\noindent
Supposons qu'il existe un entier $n \geq 0$, des ouverts quasi-compacts
$W_i \subset Y$, $i = 1, \ldots, n$, et des
morphismes $x_i : W_i \to \mathcal{X}$ tels que
(a) $f \circ x_i = \text{id}_{W_i}$,
(b) $W = \bigcup_{i = 1, \ldots, n} W_i$ contienne $V$,
(c) $W$ soit l'image de $|f|$,
(d) $V \subset W_1 \cap \ldots \cap W_n$, et
(e) $x_i$ et $x_j$ soient localement isomorphes pour la topologie de Zariski au-dessus de $W_i \cap W_j \cap V$.
Soit $y \in V$ arbitraire.
Supposons que l'on puisse trouver un voisinage ouvert quasi-compact
$y \in V_y \subset V$ tel que le lemme soit vrai pour
la paire $(\mathcal{X} \to Y, V_y)$, avec, disons, pour solution
$\overline{Z}_y, Z_y, \overline{g}_y, g_y, h_y$.
Comme $V$ est quasi-compact, on peut trouver un nombre fini de points
$y_1, \ldots, y_m$ tels que $V = V_{y_1} \cup \ldots \cup V_{y_m}$.
Il s'ensuit qu'en posant
$$
\overline{Z} = \coprod \overline{Z}_{y_j},\quad
Z = \coprod Z_{y_j},\quad
\overline{g} = \coprod \overline{g}_{y_j},\quad
g = \coprod g_{y_j},\quad
h = \coprod h_{y_j}
$$
on obtient une solution du lemme. Étant donné $y$, la condition (e)
permet de choisir un voisinage ouvert quasi-compact $y \in V_y \subset V$
et des isomorphismes $\varphi_i : x_1|_{V_y} \to x_i|_{V_y}$ pour
$i = 2, \ldots, n$. Posons $\varphi_{ij} = \varphi_j \circ \varphi_i^{-1}$.
Cela nous conduit à la situation décrite au paragraphe suivant.

\medskip\noindent
Supposons qu'il existe un entier $n \geq 0$, des ouverts quasi-compacts
$W_i \subset Y$, $i = 1, \ldots, n$, et des
morphismes $x_i : W_i \to \mathcal{X}$ tels que
(a) $f \circ x_i = \text{id}_{W_i}$,
(b) $W = \bigcup_{i = 1, \ldots, n} W_i$ contienne $V$,
(c) $W$ soit l'image de $|f|$,
(d) $V \subset W_1 \cap \ldots \cap W_n$, et
(f) il existe des isomorphismes $\varphi_{ij} : x_i|_V \to x_j|_V$
vérifiant $\varphi_{jk} \circ \varphi_{ij} = \varphi_{ik}$.
Les morphismes
$$
I_{ij} = \mathit{Isom}_\mathcal{X}(x_i|_{W_i \cap W_j},
x_j|_{W_i \cap W_j}) \longrightarrow W_i \cap W_j
$$
sont propres parce que $f$ est séparé
(Morphismes de champs, Lemme
\ref{stacks-morphisms-lemma-separated-implies-isom}).
Observons que $\varphi_{ij}$ définit une section $V \to I_{ij}$
de $I_{ij} \to W_i \cap W_j$ au-dessus de $V$.
D'après Compléments sur les morphismes d'espaces, Lemme
\ref{spaces-more-morphisms-lemma-get-section-after-blowup}
on peut trouver des éclatements $V$-admissibles
$p_{ij} : Y_{ij} \to Y$ tels que $\varphi_{ij}$
se prolonge à $p_{ij}^{-1}(W_i \cap W_j)$.
En remplaçant $Y$ par le produit fibré de tous les $Y_{ij}$
au-dessus de $Y$, on arrive à la situation décrite au paragraphe suivant.

\medskip\noindent
Supposons qu'il existe un entier $n \geq 0$, des ouverts quasi-compacts
$W_i \subset Y$, $i = 1, \ldots, n$, et des
morphismes $x_i : W_i \to \mathcal{X}$ tels que
(a) $f \circ x_i = \text{id}_{W_i}$,
(b) $W = \bigcup_{i = 1, \ldots, n} W_i$ contienne $V$,
(c) $W$ soit l'image de $|f|$,
(d) $V \subset W_1 \cap \ldots \cap W_n$, et
(g) il existe des isomorphismes
$\varphi_{ij} : x_i|_{W_i \cap W_j} \to x_j|_{W_i \cap W_j}$
vérifiant
$$
\varphi_{jk}|_V \circ \varphi_{ij}|_V = \varphi_{ik}|_V.
$$
Après avoir remplacé $Y$, si nécessaire, par un autre éclatement $V$-admissible,
on peut supposer que $V$ est dense et schématiquement dense
dans $Y$, et donc dans tout sous-espace ouvert de $Y$ contenant $V$.
Après un tel remplacement, on conclut que
$$
\varphi_{jk}|_{W_i \cap W_j \cap W_k} \circ
\varphi_{ij}|_{W_i \cap W_j \cap W_k} =
\varphi_{ik}|_{W_i \cap W_j \cap W_k}
$$
en invoquant Morphismes d'espaces, Lemme
\ref{spaces-morphisms-lemma-equality-of-morphisms}
et le fait que $I_{ik} \to W_i \cap W_k$ est propre
(donc séparé).
Bien entendu, cela signifie que $(x_i, \varphi_{ij})$
est une donnée de descente, et l'on obtient un morphisme
$x : W \to \mathcal{X}$ qui coïncide avec $x_i$ au-dessus de $W_i$,
car $\mathcal{X}$ est un champ.
Puisque $x$ est une section du morphisme séparé
$\mathcal{X} \to W$, $x$ est propre
(Morphismes de champs, Lemme \ref{stacks-morphisms-lemma-section-immersion}).
Le lemme est donc vérifié avec $\overline{Z} = Y$,
$Z = W$, $\overline{g} = \text{id}_Y$, $g = \text{id}_W$,
$h = x$.
\end{proof}

\begin{theorem}[Lemme de Chow]
\label{theorem-chow-finite-type}
\begin{reference}
Ce résultat est dû à Ofer Gabber; voir
\cite[Theorem 1.1]{olsson_proper}
\end{reference}
Soit $f : \mathcal{X} \to Y$ un morphisme d'un champ algébrique
vers un espace algébrique. Supposons que
\begin{enumerate}
\item $Y$ soit quasi-compact et quasi-séparé,
\item $f$ soit séparé de type fini.
\end{enumerate}
Alors il existe un diagramme commutatif
$$
\xymatrix{
\mathcal{X} \ar[rd] & X \ar[l] \ar[d] \ar[r] & \overline{X} \ar[ld] \\
& Y
}
$$
où $X \to \mathcal{X}$ est propre surjectif,
$X \to \overline{X}$ est une immersion ouverte, et
$\overline{X} \to Y$ est un morphisme propre d'espaces algébriques.
\end{theorem}

\begin{proof}
L'idée générale consiste à utiliser le fait que $\mathcal{X}$ possède un ouvert dense
qui est une gerbe (Morphismes de champs, Proposition
\ref{stacks-morphisms-proposition-open-stratum})
et à invoquer le Lemme \ref{lemma-make-section}.
Cela ne fonctionne pas directement, car l'ouvert peut ne pas être
quasi-compact et l'on se heurte alors à des difficultés techniques. Nous effectuons donc
d'abord une réduction (standard) au cas noethérien.

\medskip\noindent
Choisissons d'abord une immersion fermée
$\mathcal{X} \to \mathcal{X}'$, où $\mathcal{X}'$
est un champ algébrique séparé et de type fini sur $Y$.
Voir Limites de champs, Lemme
\ref{stacks-limits-lemma-separated-closed-in-finite-presentation}.
Il suffit manifestement de démontrer le théorème pour
$\mathcal{X}'$; on peut donc supposer que $\mathcal{X} \to Y$
est séparé et de présentation finie.

\medskip\noindent
Supposons que $\mathcal{X} \to Y$ soit séparé et de présentation finie.
D'après Limites d'espaces, Proposition \ref{spaces-limits-proposition-approximate},
on peut écrire $Y = \lim Y_i$ comme la limite projective d'un système filtrant
d'espaces algébriques noethériens à morphismes de transition affines.
D'après Limites de champs, Lemme \ref{stacks-limits-lemma-descend-a-stack-down},
il existe un $i$ et un morphisme $\mathcal{X}_i \to Y_i$ de présentation finie,
d'un champ algébrique vers $Y_i$, tels que
$\mathcal{X} = Y \times_{Y_i} \mathcal{X}_i$.
En augmentant $i$, on peut supposer que $\mathcal{X}_i \to Y_i$
est séparé; voir Limites de champs, Lemme
\ref{stacks-limits-lemma-eventually-separated}.
Il suffit alors de démontrer le théorème
pour $\mathcal{X}_i \to Y_i$. Cela nous ramène au cas considéré
au paragraphe suivant.

\medskip\noindent
Supposons $Y$ noethérien. On peut remplacer $\mathcal{X}$ par son
réduit (Propriétés des champs, Définition
\ref{stacks-properties-definition-reduced-induced-stack}).
Cela nous ramène au cas considéré
au paragraphe suivant.

\medskip\noindent
Supposons que $Y$ soit noethérien et que $\mathcal{X}$ soit réduit.
Puisque $\mathcal{X} \to Y$ est séparé et que $Y$ est quasi-séparé,
le champ algébrique $\mathcal{X}$ est quasi-séparé.
Ainsi, le morphisme d'inertie $\mathcal{I}_\mathcal{X} \to \mathcal{X}$
est quasi-compact. Par Morphismes de champs, Proposition
\ref{stacks-morphisms-proposition-open-stratum}
il existe un sous-champ ouvert dense $\mathcal{V} \subset \mathcal{X}$
qui est une gerbe. Soit $\mathcal{V} \to V$ le morphisme
qui présente $\mathcal{V}$ comme une gerbe sur l'espace algébrique $V$.
Voir
Morphismes de champs, Lemme \ref{stacks-morphisms-lemma-gerbe-over-iso-classes},
pour une construction de $\mathcal{V} \to V$.
Cette construction montre notamment que le morphisme
$\mathcal{V} \to Y$ se factorise en $\mathcal{V} \to V \to Y$.
On a le diagramme
$$
\xymatrix{
\mathcal{V} \ar[r] \ar[d] & \mathcal{X} \ar[d] \\
V \ar[r] & Y
}
$$
Puisque le morphisme $\mathcal{V} \to V$ est surjectif, plat et
de présentation finie
(Morphismes de champs, Lemme \ref{stacks-morphisms-lemma-gerbe-fppf})
et que $\mathcal{V} \to Y$ est localement de présentation finie,
il s'ensuit que $V \to Y$ est localement de présentation finie
(Morphismes de champs, Lemme
\ref{stacks-morphisms-lemma-flat-finite-presentation-permanence}).
Notons que $\mathcal{V} \to V$ est un homéomorphisme universel
(Morphismes de champs, Lemme
\ref{stacks-morphisms-lemma-gerbe-bijection-points}).
Puisque $\mathcal{V}$ est quasi-compact (voir
Morphismes de champs, Lemme
\ref{stacks-morphisms-lemma-locally-closed-in-noetherian}),
$V$ est quasi-compact.
Enfin, puisque $\mathcal{V} \to Y$ est séparé, il en va de même
de $V \to Y$ d'après
Morphismes de champs, Lemme
\ref{stacks-morphisms-lemma-check-separated-on-ui-cover}
appliqué à $\mathcal{V} \to V \to Y$
(dont les hypothèses sont satisfaites, comme nous l'avons déjà vu).

\medskip\noindent
Tout ce qui précède signifie que les hypothèses de
Limites d'espaces, Lemme \ref{spaces-limits-lemma-embedding-into-affine-over-qs},
s'appliquent au morphisme $V \to Y$. On peut donc trouver un sous-espace
ouvert dense $V' \subset V$ et une immersion $V' \to \mathbf{P}^n_Y$
sur $Y$. On peut manifestement remplacer $V$ par $V'$ et $\mathcal{V}$
par l'image inverse de $V'$ dans $\mathcal{V}$ (rappelons que
$|\mathcal{V}| = |V|$, comme nous l'avons vu plus haut).
On peut donc supposer donné un diagramme
$$
\xymatrix{
\mathcal{V} \ar[rr] \ar[d] & & \mathcal{X} \ar[d] \\
V \ar[r] & \mathbf{P}^n_Y \ar[r] & Y
}
$$
où la flèche $V \to \mathbf{P}^n_Y$ est une immersion.
Soit $\mathcal{X}'$ l'image schématique du morphisme
$$
j : \mathcal{V} \longrightarrow \mathbf{P}^n_Y \times_Y \mathcal{X}
$$
et soit $Y'$ l'image schématique du morphisme
$V \to \mathbf{P}^n_Y$. On obtient un diagramme commutatif
$$
\xymatrix{
\mathcal{V} \ar[r] \ar[d] &
\mathcal{X}' \ar[r] \ar[d] &
\mathbf{P}^n_Y \times_Y \mathcal{X} \ar[d] \ar[r] &
\mathcal{X} \ar[d] \\
V \ar[r] &
Y' \ar[r] &
\mathbf{P}^n_Y \ar[r] &
Y
}
$$
(Voir Morphismes de champs, Lemme \ref{stacks-morphisms-lemma-factor-factor}.)
Nous affirmons que $\mathcal{V} = V \times_{Y'} \mathcal{X}'$ et que le
Lemme \ref{lemma-make-section} s'applique au morphisme $\mathcal{X}' \to Y'$
et au sous-espace ouvert $V \subset Y'$. Si l'affirmation est vraie, on obtient
$$
\xymatrix{
\overline{X} \ar[rd]_{\overline{g}} &
X \ar[l] \ar[d]_g \ar[r]_h & \mathcal{X}' \ar[ld]^f \\
& Y'
}
$$
où $X \to \overline{X}$ est une immersion ouverte, $\overline{g}$ et $h$ sont propres,
et $|V|$ est contenu dans l'image de $|g|$.
La composée $X \to \mathcal{X}' \to \mathcal{X}$ est alors
propre (comme composée de morphismes propres) et son image
contient $|\mathcal{V}|$; cette composée est donc surjective.
De même, $\overline{X} \to Y' \to Y$ est propre en tant que composée
de morphismes propres.

\medskip\noindent
La dernière étape consiste à démontrer l'affirmation.
Observons que $\mathcal{X}' \to Y'$ est séparé et de type fini,
que $Y'$ est quasi-compact et quasi-séparé, et que $V$ est quasi-compact
(nous omettons la vérification complète de tous les détails).
Observons ensuite que
$b : \mathcal{X}' \to \mathcal{X}$ est un isomorphisme au-dessus de
$\mathcal{V}$ d'après Morphismes de champs, Lemme
\ref{stacks-morphisms-lemma-scheme-theoretic-image-of-partial-section}.
En particulier, $\mathcal{V}$ s'identifie à un sous-champ ouvert
de $\mathcal{X}'$.
Le morphisme $j$ est quasi-compact
(sa source est quasi-compacte et son but est quasi-séparé); la formation
de l'image schématique de $j$ commute donc au changement de base plat d'après
Morphismes de champs, Lemme
\ref{stacks-morphisms-lemma-existence-plus-flat-base-change}.
En particulier, $V \times_{Y'} \mathcal{X}'$ est précisément
l'image schématique de
$\mathcal{V} \to V \times_{Y'} \mathcal{X}'$.
Or, d'après Morphismes de champs, Lemme
\ref{stacks-morphisms-lemma-universally-closed-permanence}
l'image de $|\mathcal{V}| \to |V \times_{Y'} \mathcal{X}'|$
est fermée (on utilise que $\mathcal{V} \to V$ est un homéomorphisme universel,
comme nous l'avons vu plus haut, et qu'il est donc universellement fermé).
Cette image est aussi dense (combiner ce que nous venons de dire avec
Morphismes de champs, Lemme
\ref{stacks-morphisms-lemma-topology-scheme-theoretic-image})
et l'on conclut que $|\mathcal{V}| = |V \times_{Y'} \mathcal{X}'|$.
Ainsi, $\mathcal{V} \to V \times_{Y'} \mathcal{X}'$ est un
isomorphisme, ce qui achève la démonstration de l'affirmation.
\end{proof}




\section{Critère valuatif noethérien}
\label{section-Noetherian-valuative-criterion}

\noindent
Dans cette section, nous étudions des critères valuatifs (raffinés) pour les
morphismes de champs algébriques, en n'utilisant que des anneaux de valuation discrète
dans le cadre noethérien. Il en existe de nombreuses variantes,
et nous en ajouterons ici au fil du temps, selon les besoins.

\medskip\noindent
Soit $f : \mathcal{X} \to \mathcal{Y}$ un morphisme de champs algébriques
(ou d'espaces algébriques, ou encore de schémas). Un critère valuatif {\it raffiné}
est un critère où l'on se donne un morphisme $\mathcal{U} \to \mathcal{X}$
(possédant certaines propriétés) et où l'on ne considère que l'existence ou l'unicité
de flèches pointillées dans des diagrammes aux flèches pleines de la forme
$$
\xymatrix{
\Spec(K) \ar[d] \ar[r] & \mathcal{U} \ar[r] & \mathcal{X} \ar[d] \\
\Spec(A) \ar[rr] \ar@{..>}[rru] & & \mathcal{Y}
}
$$
Nous employons ci-dessous cette terminologie pour décrire les résultats obtenus jusqu'ici.

\medskip\noindent
Critères valuatifs non noethériens pour les morphismes de champs algébriques
\begin{enumerate}
\item Morphismes de champs, Section
\ref{stacks-morphisms-section-valuative-second}
(pour la séparation de la diagonale),
\item Morphismes de champs, Section
\ref{stacks-morphisms-section-valuative-diagonal}
(pour la séparation),
\item Morphismes de champs, Section
\ref{stacks-morphisms-section-valutive-criterion}
(pour le caractère universellement fermé),
\item Morphismes de champs, Section
\ref{stacks-morphisms-section-valutive-criterion-properness}
(pour la propreté).
\end{enumerate}
Pour les espaces algébriques, on dispose des critères valuatifs suivants
\begin{enumerate}
\item Morphismes d'espaces, Section
\ref{spaces-morphisms-section-valuative-criterion-universally-closed}
(pour le caractère universellement fermé),
\item Morphismes d'espaces, Lemme
\ref{spaces-morphisms-lemma-refined-valuative-criterion-universally-closed}
(forme raffinée pour le caractère universellement fermé),
\item Morphismes d'espaces, Section
\ref{spaces-morphisms-section-valuative-separatedness}
(pour la séparation),
\item Morphismes d'espaces, Section
\ref{spaces-morphisms-section-valuative-criterion-properness}
(pour la propreté),
\item Espaces décents, Section
\ref{decent-spaces-section-valuative-criterion-universally-closed}
(pour le caractère universellement fermé des espaces décents),
\item Espaces décents, Lemme
\ref{decent-spaces-lemma-re-characterize-universally-closed}
(pour le caractère universellement fermé des morphismes décents entre espaces algébriques),
\item Cohomologie des espaces, Section
\ref{spaces-cohomology-section-Noetherian-valuative-criterion}
contient des critères valuatifs noethériens
\begin{enumerate}
\item Cohomologie des espaces, Lemme
\ref{spaces-cohomology-lemma-check-separated-dvr}
(pour la séparation au moyen d'anneaux de valuation discrète),
\item Cohomologie des espaces, Lemme
\ref{spaces-cohomology-lemma-check-proper-dvr}
(pour la propreté au moyen d'anneaux de valuation discrète),
\item Cohomologie des espaces, Remarque
\ref{spaces-cohomology-remark-variant}
(explique comment se ramener à des anneaux de valuation discrète complets),
\end{enumerate}
\item Limites d'espaces, Section
\ref{spaces-limits-section-Noetherian-valuative-criterion}
consacrée aux critères valuatifs noethériens
\begin{enumerate}
\item Limites d'espaces, Lemme
\ref{spaces-limits-lemma-Noetherian-dvr-valuative-separation}
(pour la séparation au moyen d'anneaux de valuation discrète et de points génériques),
\item Limites d'espaces, Lemme
\ref{spaces-limits-lemma-Noetherian-dvr-valuative-proper}
(pour la propreté au moyen d'anneaux de valuation discrète et de points génériques),
\item Limites d'espaces, Lemme
\ref{spaces-limits-lemma-check-universally-closed-Noetherian}
(pour le caractère universellement fermé au moyen d'anneaux de valuation discrète).
\end{enumerate}
\item Limites d'espaces, Section
\ref{spaces-limits-section-refined-valuative-criteria}
consacrée aux critères valuatifs noethériens raffinés
\begin{enumerate}
\item Limites d'espaces, Lemmes
\ref{spaces-limits-lemma-refined-valuative-criterion-proper} et
\ref{spaces-limits-lemma-refined-valuative-criterion-universally-closed}
(forme raffinée pour la propreté au moyen d'anneaux de valuation discrète),
\item Limites d'espaces, Lemme
\ref{spaces-limits-lemma-refined-valuative-criterion-separated}
(forme raffinée pour la séparation au moyen d'anneaux de valuation discrète),
\end{enumerate}
\end{enumerate}
Pour les schémas, on dispose des critères valuatifs suivants
\begin{enumerate}
\item Schémas, Section
\ref{schemes-section-valuative-criterion-universal-closedness}
(pour le caractère universellement fermé),
\item Schémas, Section \ref{schemes-section-valuative-separatedness}
(pour la séparation),
\item Morphismes, Section \ref{morphisms-section-valuative-criteria}
(pour la propreté),
\item Morphismes, Lemme
\ref{morphisms-lemma-refined-valuative-criterion-universally-closed}
(forme raffinée pour le caractère universellement fermé),
\item Limites, Section \ref{limits-section-Noetherian-valuative-criterion}
consacrée aux critères valuatifs noethériens
\begin{enumerate}
\item Limites, Lemme \ref{limits-lemma-Noetherian-dvr-valuative-separation}
(pour la séparation au moyen d'anneaux de valuation discrète et de points génériques),
\item Limites, Lemme \ref{limits-lemma-Noetherian-dvr-valuative-proper}
(pour la propreté au moyen d'anneaux de valuation discrète et de points génériques),
\item Limites, Lemme \ref{limits-lemma-check-universally-closed-Noetherian}
(pour le caractère universellement fermé au moyen d'anneaux de valuation discrète).
\end{enumerate}
\item Limites, Section \ref{limits-section-refined-valuative-criteria}
consacrée aux critères valuatifs noethériens raffinés
\begin{enumerate}
\item Limites, Lemmes
\ref{limits-lemma-refined-valuative-criterion-proper} et
\ref{limits-lemma-refined-valuative-criterion-universally-closed}
(forme raffinée pour la propreté au moyen d'anneaux de valuation discrète),
\item Limites, Lemme \ref{limits-lemma-refined-valuative-criterion-separated}
(forme raffinée pour la séparation au moyen d'anneaux de valuation discrète),
\end{enumerate}
\item Limites, Section \ref{limits-section-nagata-valuative}
consacrée aux critères valuatifs sur une base noethérienne, pour lesquels
on peut obtenir des anneaux de valuation discrète essentiellement de type fini
sur la base.
\end{enumerate}
Cela termine notre liste de résultats antérieurs.

\medskip\noindent
Nombre des résultats de cette section peuvent (et devraient peut-être)
être démontrés en invoquant le lemme suivant, même si nous ne l'avons pas
toujours fait.

\begin{lemma}
\label{lemma-reach-point-closure-Noetherian}
Soit $f : \mathcal{X} \to \mathcal{Y}$ un morphisme de champs algébriques.
Supposons $f$ de type fini et $\mathcal{Y}$ localement noethérien.
Soit $y \in |\mathcal{Y}|$ un point de l'adhérence de l'image de $|f|$.
Alors il existe un diagramme commutatif
$$
\xymatrix{
\Spec(K) \ar[r] \ar[d] & \mathcal{X} \ar[d]^f \\
\Spec(A) \ar[r] & \mathcal{Y}
}
$$
de champs algébriques, où $A$ est un anneau de valuation discrète et $K$
son corps des fractions, le point fermé de $\Spec(A)$ étant envoyé sur $y$.
\end{lemma}

\begin{proof}
Choisissons un schéma affine $V$, un point $v \in V$ et un morphisme lisse
$V \to \mathcal{Y}$ qui envoie $v$ sur $y$. L'application $|V| \to |\mathcal{Y}|$
est ouverte et, d'après
Propriétés des champs, Lemme \ref{stacks-properties-lemma-points-cartesian},
l'image de $|\mathcal{X} \times_\mathcal{Y} V| \to |V|$
est l'image inverse de l'image de $|f|$.
On en déduit que le point $v$ appartient à l'adhérence de
l'image de $|\mathcal{X} \times_\mathcal{Y} V| \to |V|$.
Si l'on démontre le lemme pour
$\mathcal{X} \times_\mathcal{Y} V \to V$ et le point $v$, alors
le lemme en résulte pour $f$ et $y$.
On se ramène ainsi à la situation décrite au
paragraphe suivant.

\medskip\noindent
Supposons donnés $f : \mathcal{X} \to Y$ et $y \in |Y|$ comme dans le lemme, où
$Y$ est un schéma affine noethérien. Puisque $f$ est quasi-compact, on en déduit que
$\mathcal{X}$ est quasi-compact. On peut donc choisir un schéma affine $W$ et
un morphisme lisse surjectif $W \to \mathcal{X}$. L'image de
$|f|$ est alors la même que celle de $|W| \to |Y|$. On se ramène ainsi
au cas des schémas, qui est
Limites, Lemme \ref{limits-lemma-reach-point-closure-Noetherian}.
\end{proof}

\begin{lemma}
\label{lemma-check-separated-dvr}
Soit $f : \mathcal{X} \to \mathcal{Y}$ un morphisme de champs algébriques. Supposons que
\begin{enumerate}
\item $\mathcal{Y}$ soit localement noethérien,
\item $f$ soit localement de type fini et quasi-séparé,
\item pour tout diagramme commutatif
$$
\xymatrix{
\Spec(K) \ar[r]_x \ar[d]_j & \mathcal{X} \ar[d]^f \\
\Spec(A) \ar[r]^y \ar@{-->}[ru] & \mathcal{Y}
}
$$
où $A$ est un anneau de valuation discrète, $K$ son corps des fractions,
et pour toute $2$-flèche $\gamma : y \circ j \to f \circ x$, la catégorie
des flèches pointillées (Morphismes de champs, Définition
\ref{stacks-morphisms-definition-fill-in-diagram})
soit vide, soit un sétoïde ayant exactement une classe d'isomorphisme.
\end{enumerate}
Alors $f$ est séparé.
\end{lemma}

\begin{proof}
Pour démontrer que $f$ est séparé, il faut montrer que
$\Delta : \mathcal{X} \to \mathcal{X} \times_\mathcal{Y} \mathcal{X}$
est propre. On sait déjà que $\Delta$ est représentable par
des espaces algébriques, localement de type fini (Morphismes de champs,
Lemme \ref{stacks-morphisms-lemma-properties-diagonal}) et
quasi-compact et quasi-séparé (par définition de la quasi-séparation de $f$).
Choisissons un schéma $U$ et un morphisme lisse surjectif
$U \to \mathcal{X} \times_\mathcal{Y} \mathcal{X}$.
Posons
$$
V = \mathcal{X} \times_{\Delta, \mathcal{X} \times_\mathcal{Y} \mathcal{X}} U
$$
Il suffit de montrer que le morphisme d'espaces algébriques $V \to U$
est propre (Propriétés des champs, Lemme
\ref{stacks-properties-lemma-check-property-covering}).
Observons que $U$ est localement noethérien
(utiliser Morphismes de champs, Lemme
\ref{stacks-morphisms-lemma-locally-finite-type-locally-noetherian}
et le fait que $U \to \mathcal{Y}$ est localement de type fini),
et que $V \to U$
est de type fini et quasi-séparé (comme changement de base
de $\Delta$, compte tenu des propriétés de $\Delta$ énumérées ci-dessus). D'après
Cohomologie des espaces, Lemme \ref{spaces-cohomology-lemma-check-proper-dvr},
il suffit de montrer ceci : étant donné un diagramme commutatif
$$
\xymatrix{
\Spec(K) \ar[r]_v \ar[d]_j &
V \ar[d]^g \ar[r] &
\mathcal{X} \ar[d]^\Delta \\
\Spec(A) \ar[r]^u \ar@{-->}[ru] \ar@{..>}[rru] &
U \ar[r] &
\mathcal{X} \times_\mathcal{Y} \mathcal{X}
}
$$
où $A$ est un anneau de valuation discrète et $K$ son corps des fractions,
il existe une unique flèche en tirets qui rend le diagramme commutatif.
D'après Morphismes de champs, Lemme
\ref{stacks-morphisms-lemma-cat-dotted-arrows-base-change}
les catégories des flèches en tirets et des flèches pointillées sont équivalentes.
L'hypothèse (3) entraîne l'existence d'une unique flèche pointillée à
isomorphisme près; voir Morphismes de champs, Lemme
\ref{stacks-morphisms-lemma-helper-diagonal}. On en déduit
l'existence de l'unique flèche en tirets voulue.
\end{proof}

\begin{lemma}
\label{lemma-refined-valuative-criterion-proper}
Soient $f : \mathcal{X} \to \mathcal{Y}$ et $h : \mathcal{U} \to \mathcal{X}$
des morphismes de champs algébriques. Supposons que $\mathcal{Y}$ soit
localement noethérien, que $f$ et $h$ soient de type fini,
que $f$ soit séparé, et que l'image de
$|h| : |\mathcal{U}| \to |\mathcal{X}|$ soit dense dans $|\mathcal{X}|$.
Si, pour tout diagramme $2$-commutatif
$$
\xymatrix{
\Spec(K) \ar[r]_-u \ar[d]_j & \mathcal{U} \ar[r]_h & \mathcal{X} \ar[d]^f \\
\Spec(A) \ar[rr]^-y & & \mathcal{Y}
}
$$
où $A$ est un anneau de valuation discrète de corps des fractions $K$
et $\gamma : y \circ j \to f \circ h \circ u$, il
existe une extension de corps $K'/K$ et un anneau de valuation $A' \subset K'$
dominant $A$, tels que la catégorie des flèches pointillées du
diagramme induit
$$
\xymatrix{
\Spec(K') \ar[r]_-{x'} \ar[d]_{j'} & \mathcal{X} \ar[d]^f \\
\Spec(A') \ar[r]^-{y'} \ar@{..>}[ru] & \mathcal{Y}
}
$$
muni de la $2$-flèche induite $\gamma' : y' \circ j' \to f \circ x'$
soit non vide (Morphismes de champs, Définition
\ref{stacks-morphisms-definition-fill-in-diagram}), alors $f$ est propre.
\end{lemma}

\begin{proof}
Il suffit de démontrer que $f$ est universellement fermé.
Soit $V \to \mathcal{Y}$ un morphisme lisse, où $V$ est un schéma affine.
D'après Propriétés des champs, Lemme
\ref{stacks-properties-lemma-points-cartesian}
l'image $I$ de
$|\mathcal{U} \times_\mathcal{Y} V| \to |\mathcal{X} \times_\mathcal{Y} V|$
est l'image inverse de l'image de $|h|$. Puisque
$|\mathcal{X} \times_\mathcal{Y} V| \to |\mathcal{X}|$ est ouverte
(Morphismes de champs, Lemme \ref{stacks-morphisms-lemma-fppf-open}),
on en déduit que $I$ est dense dans $|\mathcal{X} \times_\mathcal{Y} V|$.
En outre, comme la catégorie des flèches pointillées se comporte bien par
changement de base (Morphismes de champs, Lemme
\ref{stacks-morphisms-lemma-cat-dotted-arrows-base-change})
les morphismes suivants héritent de l'hypothèse d'existence de flèches pointillées
(après extension) :
$\mathcal{U} \times_\mathcal{Y} V \to \mathcal{X} \times_\mathcal{Y} V \to V$.
Les hypothèses du lemme sont donc
satisfaites pour les morphismes
$\mathcal{U} \times_\mathcal{Y} V \to \mathcal{X} \times_\mathcal{Y} V \to V$.
On peut par conséquent supposer que $\mathcal{Y}$ est un schéma affine.

\medskip\noindent
Supposons que $\mathcal{Y} = Y$ soit un schéma affine.
(Désormais, nous n'avons plus à nous soucier des
$2$-flèches $\gamma$ et $\gamma'$; voir Morphismes de champs, Lemme
\ref{stacks-morphisms-lemma-cat-dotted-arrows-independent}.)
Alors $\mathcal{U}$ est quasi-compact. Choisissons un schéma affine $U$ et un
morphisme lisse surjectif $U \to \mathcal{U}$.
Nous remplaçons alors $\mathcal{U}$ par $U$.
On peut donc supposer que $\mathcal{U}$ est un schéma affine.

\medskip\noindent
Supposons que $\mathcal{Y} = Y$ et $\mathcal{U} = U$ soient des schémas affines.
D'après le lemme de Chow (Théorème \ref{theorem-chow-finite-type}),
on peut choisir un morphisme propre surjectif $X \to \mathcal{X}$,
où $X$ est un espace algébrique. Nous utiliserons plus bas que $X \to Y$ est séparé
comme composé de morphismes séparés. Considérons
l'espace algébrique $W = X \times_\mathcal{X} U$. Le morphisme de projection
$W \to X$ est de type fini.
On peut remplacer $X$ par l'image schématique de
$W \to X$ et supposer ainsi que l'image de $|W|$ dans $|X|$
est dense dans $|X|$ (on utilise ici que l'image de $|h|$
est dense dans $|\mathcal{X}|$, de sorte qu'après ce remplacement, le
morphisme $X \to \mathcal{X}$ reste surjectif).
Nous affirmons que, pour tout diagramme commutatif aux flèches pleines
$$
\xymatrix{
\Spec(K) \ar[r] \ar[d] & W \ar[r] & X \ar[d] \\
\Spec(A) \ar[rr] \ar@{..>}[rru] & & Y
}
$$
où $A$ est un anneau de valuation discrète de corps des fractions $K$, il
existe une flèche pointillée qui rend le diagramme commutatif. Tout d'abord, il suffit de
montrer qu'une telle flèche existe après avoir remplacé $K$ par une extension
et $A$ par un anneau de valuation de cette extension dominant $A$; voir
Morphismes d'espaces, Lemme \ref{spaces-morphisms-lemma-push-down-solution}.
L'hypothèse du lemme fournit une extension $K'/K$, un anneau de
valuation $A' \subset K'$ dominant $A$, ainsi
qu'une flèche $\Spec(A') \to \mathcal{X}$ qui relève la composée
$\Spec(A') \to \Spec(A) \to Y$ et qui est compatible avec la composée
$\Spec(K') \to \Spec(K) \to W \to X$. Puisque $X \to \mathcal{X}$
est propre, le critère valuatif de propreté
(Morphismes de champs, Lemme \ref{stacks-morphisms-lemma-criterion-proper})
donne une extension $K''/K'$, un anneau de valuation $A'' \subset K''$
dominant $A'$, et un morphisme $\Spec(A'') \to X$ qui relève la composée
$\Spec(A'') \to \Spec(A') \to \mathcal{X}$ et qui est compatible avec la composée
$\Spec(K'') \to \Spec(K') \to \Spec(K) \to X$.
L'extension $K''/K$, l'anneau $A'' \subset K''$ et le morphisme $\Spec(A'') \to X$
fournissent alors une solution au problème posé ci-dessus, ce qui démontre l'affirmation.

\medskip\noindent
L'affirmation entraîne que le morphisme $X \to Y$ est propre, par le
cas du lemme relatif aux espaces algébriques
(Limites d'espaces, Lemme
\ref{spaces-limits-lemma-refined-valuative-criterion-proper}).
D'après Morphismes de champs, Lemme
\ref{stacks-morphisms-lemma-image-proper-is-proper}
on en déduit que $\mathcal{X} \to Y$ est propre, comme souhaité.
\end{proof}

\begin{lemma}
\label{lemma-refined-valuative-criterion-separated}
Soient $f : \mathcal{X} \to \mathcal{Y}$ et $h : \mathcal{U} \to \mathcal{X}$
des morphismes de champs algébriques. Supposons que $\mathcal{Y}$ soit
localement noethérien, que $f$ soit localement de type fini et quasi-séparé,
que $h$ soit de type fini, et que l'image de
$|h| : |\mathcal{U}| \to |\mathcal{X}|$ soit dense dans $|\mathcal{X}|$.
Si, pour tout diagramme $2$-commutatif
$$
\xymatrix{
\Spec(K) \ar[r]_-u \ar[d]_j & \mathcal{U} \ar[r]_h & \mathcal{X} \ar[d]^f \\
\Spec(A) \ar[rr]^-y \ar@{..>}[rru] & & \mathcal{Y}
}
$$
où $A$ est un anneau de valuation discrète de corps des fractions $K$
et $\gamma : y \circ j \to f \circ h \circ u$, la catégorie
des flèches pointillées est soit vide, soit un sétoïde ayant exactement
une classe d'isomorphisme, alors $f$ est séparé.
\end{lemma}

\begin{proof}
Il faut démontrer que $\Delta$ est un morphisme propre.
Supposons d'abord que $\Delta$ soit séparé. On peut alors appliquer le
Lemme \ref{lemma-refined-valuative-criterion-proper}
aux morphismes $\mathcal{U} \to \mathcal{X}$ et
$\Delta : \mathcal{X} \to \mathcal{X} \times_\mathcal{Y} \mathcal{X}$.
Observons que $\Delta$ est quasi-compact puisque $f$ est quasi-séparé.
Bien entendu, $\Delta$ est localement de type fini (cela vaut pour tout
morphisme diagonal; voir Morphismes de champs, Lemme
\ref{stacks-morphisms-lemma-properties-diagonal}).
Enfin, supposons donné un diagramme $2$-commutatif
$$
\xymatrix{
\Spec(K) \ar[r]_-u \ar[d]_j &
\mathcal{U} \ar[r]_h &
\mathcal{X} \ar[d]^\Delta \\
\Spec(A) \ar[rr]^-y \ar@{..>}[rru] & &
\mathcal{X} \times_\mathcal{Y} \mathcal{X}
}
$$
où $A$ est un anneau de valuation discrète de corps des fractions $K$
et $\gamma : y \circ j \to \Delta \circ h \circ u$.
D'après Morphismes de champs, Lemme \ref{stacks-morphisms-lemma-helper-diagonal},
et l'hypothèse du lemme,
il existe une unique flèche pointillée.
Cela établit la dernière hypothèse du
Lemme \ref{lemma-refined-valuative-criterion-proper},
et le résultat en découle.

\medskip\noindent
Dans le cas général, il suffit de démontrer que $\Delta$ est séparé,
car on sera alors ramené au cas précédent. Nous affirmons en fait
que les hypothèses du lemme sont satisfaites pour
$$
\mathcal{U} \to \mathcal{X}
\quad\text{et}\quad
\Delta :
\mathcal{X} \to
\mathcal{X} \times_\mathcal{Y} \mathcal{X}
$$
En effet, puisque $\Delta$ est représentable par des espaces algébriques, la
catégorie des flèches pointillées d'un diagramme comme au paragraphe précédent
est un sétoïde (voir par exemple
Morphismes de champs, Lemme \ref{stacks-morphisms-lemma-cat-dotted-arrows}).
L'argument du paragraphe précédent montre que ces catégories
sont soit vides, soit munies d'une unique classe d'isomorphisme.
Ainsi, $\Delta$ est séparé.
\end{proof}

\begin{lemma}
\label{lemma-valuative-criterion-universally-closed}
Soit $f : \mathcal{X} \to \mathcal{Y}$ un morphisme de champs algébriques.
Supposons que $\mathcal{Y}$ soit localement noethérien et que $f$ soit de type fini.
Si, pour tout diagramme $2$-commutatif
$$
\xymatrix{
\Spec(K) \ar[r]_-x \ar[d]_j & \mathcal{X} \ar[d]^f \\
\Spec(A) \ar[r]^-y & \mathcal{Y}
}
$$
où $A$ est un anneau de valuation discrète de corps des fractions $K$
et $\gamma : y \circ j \to f \circ x$, il existe une extension de corps $K'/K$
et un anneau de valuation $A' \subset K'$ dominant $A$, tels que
la catégorie des flèches pointillées du diagramme induit
$$
\xymatrix{
\Spec(K') \ar[r]_-{x'} \ar[d]_{j'} & \mathcal{X} \ar[d]^f \\
\Spec(A') \ar[r]^-{y'} \ar@{..>}[ru] & \mathcal{Y}
}
$$
muni de la $2$-flèche induite $\gamma' : y' \circ j' \to f \circ x'$
soit non vide (Morphismes de champs, Définition
\ref{stacks-morphisms-definition-fill-in-diagram}), alors $f$ est
universellement fermé.
\end{lemma}

\begin{proof}
Soit $V \to \mathcal{Y}$ un morphisme lisse, où $V$ est un schéma affine.
La catégorie des flèches pointillées se comporte bien par changement de base
(Morphismes de champs, Lemme
\ref{stacks-morphisms-lemma-cat-dotted-arrows-base-change}).
L'hypothèse d'existence de flèches pointillées (après extension)
est donc héritée par le morphisme $\mathcal{X} \times_\mathcal{Y} V \to V$.
Les hypothèses du lemme sont par conséquent satisfaites pour le morphisme
$\mathcal{X} \times_\mathcal{Y} V \to V$.
On peut donc supposer que $\mathcal{Y}$ est un schéma affine.

\medskip\noindent
Supposons que $\mathcal{Y} = Y$ soit un schéma affine noethérien.
(Désormais, nous n'avons plus à nous soucier des
$2$-flèches $\gamma$ et $\gamma'$; voir Morphismes de champs, Lemme
\ref{stacks-morphisms-lemma-cat-dotted-arrows-independent}.)
Pour démontrer que $f$ est universellement fermé, il suffit de montrer que
$|\mathcal{X} \times \mathbf{A}^n| \to |Y \times \mathbf{A}^n|$ est fermée
pour tout $n$, d'après Limites de champs, Lemme
\ref{stacks-limits-lemma-test-universally-closed}. Comme l'hypothèse
du lemme est héritée par le morphisme produit
$\mathcal{X} \times \mathbf{A}^n \to Y \times \mathbf{A}^n$ (nous omettons les détails),
on se ramène à démontrer que $|\mathcal{X}| \to |Y|$ est fermée.

\medskip\noindent
Supposons que $Y$ soit un schéma affine noethérien.
Soit $T \subset |\mathcal{X}|$ une partie fermée. Il faut montrer que
l'image de $T$ dans $|Y|$ est fermée. On peut remplacer $\mathcal{X}$
par la structure réduite induite de sous-champ fermé sur $T$; nous
omettons de vérifier que la propriété d'existence de
flèches pointillées est préservée par ce remplacement.
On se ramène ainsi à démontrer que l'image
de $|\mathcal{X}| \to |Y|$ est fermée.

\medskip\noindent
Soit $y \in |Y|$ un point de l'adhérence de l'image de
$|\mathcal{X}| \to |Y|$. D'après le Lemme \ref{lemma-reach-point-closure-Noetherian},
on peut choisir un diagramme commutatif
$$
\xymatrix{
\Spec(K) \ar[r] \ar[d] & \mathcal{X} \ar[d]^f \\
\Spec(A) \ar[r] & Y
}
$$
où $A$ est un anneau de valuation discrète et $K$ son corps des fractions,
le point fermé de $\Spec(A)$ étant envoyé sur $y$. Il résulte immédiatement
de l'hypothèse du lemme que $y$ appartient à l'image de
$|\mathcal{X}| \to |Y|$, ce qui achève la démonstration.
\end{proof}










\section{Espaces de modules}
\label{section-moduli-spaces}

\noindent
Cette section porte sur les morphismes $f : \mathcal{X} \to Y$ de
champs algébriques vers des espaces algébriques. Sous des hypothèses convenables,
$Y$ est appelé un {\it espace de modules} de $\mathcal{X}$. Si
$\mathcal{X} = [U/R]$ est une présentation, on obtient alors un
morphisme $R$-invariant $U \to Y$ et, sous des hypothèses convenables,
$Y$ est un {\it quotient} du groupoïde $(U, R, s, t, c)$.
On trouvera une étude des différents types de quotients
à partir de
Quotients de groupoïdes, Section \ref{groupoids-quotients-section-introduction}.

\begin{definition}
\label{definition-categorical-quotient}
Soit $\mathcal{X}$ un champ algébrique. Soit
$f : \mathcal{X} \to Y$ un morphisme vers un espace algébrique $Y$.
\begin{enumerate}
\item On dit que $f$ est un {\it espace de modules catégorique} si tout morphisme
$\mathcal{X} \to W$ vers un espace algébrique $W$ se factorise de manière unique par $f$.
\item On dit que $f$ est un {\it espace de modules catégorique uniforme}
si, pour tout morphisme plat $Y' \to Y$ d'espaces algébriques, le changement de base
$f' : Y' \times_Y \mathcal{X} \to Y'$ est un espace de modules catégorique.
\end{enumerate}
Soit $\mathcal{C}$ une sous-catégorie pleine de la catégorie des espaces
algébriques.
\begin{enumerate}
\item[(3)] On dit que $f$ est un {\it espace de modules catégorique dans $\mathcal{C}$}
si $Y \in \Ob(\mathcal{C})$ et si tout morphisme $\mathcal{X} \to W$ avec
$W \in \Ob(\mathcal{C})$ se factorise de manière unique par $f$.
\item[(4)] On dit que $f$ est un {\it espace de modules catégorique uniforme dans $\mathcal{C}$}
si $Y \in \Ob(\mathcal{C})$ et si, pour tout morphisme plat $Y' \to Y$ dans
$\mathcal{C}$, le changement de base $f' : Y' \times_Y \mathcal{X} \to Y'$ est un
espace de modules catégorique dans $\mathcal{C}$.
\end{enumerate}
\end{definition}

\noindent
D'après le lemme de Yoneda, un espace de modules catégorique, s'il existe, est unique.
Rapprochons cette notion du langage introduit pour les quotients.

\begin{lemma}
\label{lemma-quotient-compare}
Soit $(U, R, s, t, c)$ un groupoïde en espaces algébriques tel que
$s, t : R \to U$ soient plats et localement de présentation finie.
Considérons le champ algébrique $\mathcal{X} = [U/R]$.
Pour tout espace algébrique $Y$, il existe une correspondance biunivoque ($1$ à $1$) entre les
morphismes $f : \mathcal{X} \to Y$ et les morphismes $R$-invariants
$\phi : U \to Y$.
\end{lemma}

\begin{proof}
Critères de représentabilité, Théorème
\ref{criteria-theorem-flat-groupoid-gives-algebraic-stack}
nous dit que $\mathcal{X}$ est un champ algébrique.
Étant donné un morphisme $f : \mathcal{X} \to Y$, soit $\phi : U \to Y$
la composée $U \to \mathcal{X} \to Y$. Puisque $R = U \times_\mathcal{X} U$
(Groupoïdes en espaces, Lemme
\ref{spaces-groupoids-lemma-quotient-stack-2-cartesian})
il est immédiat que $\phi$ est $R$-invariant.
Réciproquement, si $\phi : U \to Y$ est un morphisme $R$-invariant vers
un espace algébrique, on obtient un morphisme
$f : \mathcal{X} \to Y$ d'après
Groupoïdes en espaces, Lemme
\ref{spaces-groupoids-lemma-quotient-stack-2-coequalizer}.
On peut aussi construire $f$ à partir de $\phi$ en utilisant la description explicite du
champ quotient donnée dans
Groupoïdes en espaces, Lemme
\ref{spaces-groupoids-lemma-quotient-stack-objects}.
\end{proof}

\begin{lemma}
\label{lemma-categorical-quotient-compare}
Sous les hypothèses et avec les notations du Lemme \ref{lemma-quotient-compare},
$f$ est un espace de modules catégorique (uniforme)
si et seulement si $\phi$ est un quotient catégorique (uniforme).
Il en va de même des espaces de modules dans une sous-catégorie pleine.
\end{lemma}

\begin{proof}
La correspondance biunivoque ($1$ à $1$) établie dans le
Lemme \ref{lemma-quotient-compare} montre immédiatement que $f$ est un espace de modules catégorique
si et seulement si $\phi$ est un quotient catégorique
(Quotients de groupoïdes, Définition
\ref{groupoids-quotients-definition-categorical}).
Si $Y' \to Y$ est un morphisme, alors
$U' = Y' \times_Y U \to Y' \times_Y \mathcal{X} = \mathcal{X}'$
est un morphisme surjectif, plat et localement de présentation finie
comme changement de base de $U \to \mathcal{X}$
(Critères de représentabilité, Lemme
\ref{criteria-lemma-flat-quotient-flat-presentation}).
De plus, $R' = Y' \times_Y R$ est égal à $U' \times_{\mathcal{X}'} U'$
par transitivité des produits fibrés.
Ainsi, $\mathcal{X}' = [U'/R']$; voir
Champs algébriques, Remarque \ref{algebraic-remark-flat-fp-presentation}.
Le changement de base de notre situation à $Y'$ est donc une nouvelle situation
du type décrit dans l'énoncé du lemme. Il en
résulte immédiatement que $f$ est un espace de modules catégorique uniforme
si et seulement si $\phi$ est un quotient catégorique uniforme.
\end{proof}

\begin{lemma}
\label{lemma-check-uniform-categorical-quotient-on-affines}
Soit $f : \mathcal{X} \to Y$ un morphisme d'un champ algébrique
vers un espace algébrique. Si, pour tout schéma affine $Y'$ et tout
morphisme plat $Y' \to Y$, le changement de base
$f' : Y' \times_Y \mathcal{X} \to Y'$ est un espace de modules catégorique,
alors $f$ est un espace de modules catégorique uniforme.
\end{lemma}

\begin{proof}
Choisissons un recouvrement \'etale $\{Y_i \to Y\}$ où $Y_i$ est un schéma affine.
Pour chaque $i$ et $j$, choisissons un recouvrement par des ouverts affines
$Y_i \times_Y Y_j = \bigcup Y_{ijk}$.
Posons $\mathcal{X}_i = Y_i \times_Y \mathcal{X}$ et
$\mathcal{X}_{ijk} = Y_{ijk} \times_Y \mathcal{X}$.
Soit $g : \mathcal{X} \to W$ un morphisme vers
un espace algébrique. Considérons alors le diagramme
$$
\xymatrix{
\mathcal{X}_i \ar[r] \ar[d] & \mathcal{X} \ar[d] \ar[r]_g & W \\
Y_i \ar[r] \ar@{..>}[rru] & Y
}
$$
L'hypothèse que $\mathcal{X}_i \to Y_i$ est un espace de modules
catégorique fournit une unique flèche pointillée $h_i : Y_i \to W$.
L'hypothèse que $\mathcal{X}_{ijk} \to Y_{ijk}$ est un espace de modules
catégorique entraîne que les restrictions de $h_i$ et $h_j$ à
$Y_{ijk}$ sont égales. Ainsi, $h_i$ et $h_j$ coïncident sur $Y_i \times_Y Y_j$.
Puisque $Y = \coprod Y_i / \coprod Y_i \times_Y Y_j$
(d'après Espaces, Section \ref{spaces-section-presentations}), on en déduit
qu'il existe un unique morphisme $Y \to W$ par lequel $g$ se factorise.
Ainsi, $f$ est un espace de modules catégorique. Le même argument s'applique
après un changement de base plat; $f$ est donc un espace de modules catégorique uniforme.
\end{proof}







\section{Le théorème de Keel--Mori}
\label{section-Keel-Mori}

\noindent
Dans cette section, nous commençons l'étude du théorème de Keel et Mori
dans le cadre des champs algébriques. Pour une présentation de la
littérature, voir Guide bibliographique, Sous-section
\ref{guide-subsection-coarse-moduli-spaces}.

\begin{definition}
\label{definition-well-nigh-affine}
Soit $\mathcal{X}$ un champ algébrique. On dit que $\mathcal{X}$
est {\it presque affine} s'il existe un schéma affine $U$
et un morphisme surjectif, plat, fini et de présentation finie
$U \to \mathcal{X}$.
\end{definition}

\noindent
Nous donnons à cette propriété un nom quelque peu ridicule, car nous ne comptons pas
l'utiliser très souvent.

\begin{lemma}
\label{lemma-well-nigh-affine}
Soit $\mathcal{X}$ un champ algébrique. Les assertions suivantes sont équivalentes :
\begin{enumerate}
\item $\mathcal{X}$ est presque affine, et
\item il existe un schéma en groupoïdes $(U, R, s, t, c)$ avec $U$ et
$R$ affines et $s, t : R \to U$ finis localement libres, tel que
$\mathcal{X} = [U/R]$.
\end{enumerate}
Lorsque ces assertions sont vraies, $\mathcal{X}$ est quasi-compact, quasi-DM et séparé.
\end{lemma}

\begin{proof}
Supposons $\mathcal{X}$ presque affine. Choisissons un schéma affine $U$
et un morphisme surjectif, plat, fini et de présentation finie
$U \to \mathcal{X}$. Posons $R = U \times_\mathcal{X} U$. On
obtient alors un groupoïde $(U, R, s, t, c)$ en espaces algébriques et un
isomorphisme $[U/R] \to \mathcal{X}$; voir
Champs algébriques, Lemme \ref{algebraic-lemma-map-space-into-stack}
et Remarque \ref{algebraic-remark-flat-fp-presentation}.
Puisque $s, t : R \to U$ sont des morphismes
plats, finis et de présentation finie
(comme changements de base de $U \to \mathcal{X})$, on voit que
$s, t$ sont finis localement libres
(Morphismes, Lemme \ref{morphisms-lemma-finite-flat}).
Il en résulte que $R$ est affine (car les morphismes finis sont affines)
et donc que (2) est vérifiée.

\medskip\noindent
Supposons donné un schéma en groupoïdes $(U, R, s, t, c)$ avec $U$ et
$R$ affines et $s, t : R \to U$ finis localement libres.
Posons $\mathcal{X} = [U/R]$. Alors $\mathcal{X}$ est un champ algébrique
d'après Critères de représentabilité, Théorème
\ref{criteria-theorem-flat-groupoid-gives-algebraic-stack} (à strictement parler,
nous n'en avons pas besoin ici, mais on ne saurait trop insister sur ce fait).
Le morphisme $U \to \mathcal{X}$ est surjectif, plat et localement
de présentation finie d'après
Critères de représentabilité, Lemme
\ref{criteria-lemma-flat-quotient-flat-presentation}.
On peut donc vérifier que $U \to \mathcal{X}$ est fini en
vérifiant que la projection $U \times_\mathcal{X} U \to U$
possède cette propriété; voir Propriétés des champs, Lemme
\ref{stacks-properties-lemma-check-property-covering}.
Puisque $U \times_\mathcal{X} U = R$ d'après
Groupoïdes en espaces, Lemme
\ref{spaces-groupoids-lemma-quotient-stack-2-cartesian}
on voit que tel est le cas. Ainsi, $\mathcal{X}$ est presque affine.

\medskip\noindent
Démontrons l'assertion finale. On voit que $\mathcal{X}$
est quasi-compact d'après Propriétés des champs, Lemme
\ref{stacks-properties-lemma-quasi-compact-stack}.
On voit que $\mathcal{X} = [U/R]$ est quasi-DM et séparé d'après
Morphismes de champs, Lemme
\ref{stacks-morphisms-lemma-properties-diagonal-from-presentation}.
\end{proof}

\begin{lemma}
\label{lemma-affine-over-well-nigh-affine}
Soit $\mathcal{X}$ un champ algébrique presque affine.
\begin{enumerate}
\item Si $\mathcal{X}$ est un espace algébrique, alors il est affine.
\item Si $\mathcal{X}' \to \mathcal{X}$ est un morphisme affine
de champs algébriques, alors $\mathcal{X}'$ est presque affine.
\end{enumerate}
\end{lemma}

\begin{proof}
La partie (1) résulte immédiatement de
Limites d'espaces, Lemme \ref{spaces-limits-lemma-affine}.
Ce résultat est toutefois surdimensionné, car (1) résulte aussi du
Lemme \ref{lemma-well-nigh-affine} combiné avec
Groupoïdes, Proposition
\ref{groupoids-proposition-finite-flat-equivalence}.

\medskip\noindent
Pour démontrer (2), choisissons un schéma affine $U$ et un
morphisme surjectif, plat, fini et de présentation finie $U \to \mathcal{X}$.
Alors $U' = \mathcal{X}' \times_\mathcal{X} U$ admet un morphisme
affine vers $U$ (Morphismes de champs, Lemme
\ref{stacks-morphisms-lemma-base-change-affine}).
Ainsi, $U'$ est un schéma affine. Bien entendu,
$U' \to \mathcal{X}'$ est surjectif, plat, fini et de présentation finie
comme changement de base de $U \to \mathcal{X}$.
\end{proof}

\begin{lemma}
\label{lemma-well-nigh-affine-moduli-space}
Soit $\mathcal{X}$ un champ algébrique presque affine. Il existe
un espace de modules catégorique uniforme
$$
f : \mathcal{X} \longrightarrow M
$$
dans la catégorie des schémas affines. De plus,
$f$ est séparé, quasi-compact et un homéomorphisme universel.
\end{lemma}

\begin{proof}
Écrivons $\mathcal{X} = [U/R]$, avec $(U, R, s, t, c)$ comme dans le
Lemme \ref{lemma-well-nigh-affine}. Soit $C$ l'anneau des
fonctions $R$-invariantes sur $U$; voir
Groupoïdes, Section \ref{groupoids-section-finite-flat}.
Posons $M = \Spec(C)$. Le morphisme $R$-invariant
$U \to M$ correspond à un morphisme $f : \mathcal{X} \to M$ d'après le
Lemme \ref{lemma-quotient-compare}.
La caractérisation des morphismes vers les schémas affines donnée dans
Schémas, Lemme \ref{schemes-lemma-morphism-into-affine},
garantit immédiatement que $\phi : U \to M$ est un quotient
catégorique dans la catégorie des schémas affines. Ainsi, $f$ est un
espace de modules catégorique dans la catégorie des schémas affines
(Lemme \ref{lemma-categorical-quotient-compare}).

\medskip\noindent
Puisque $\mathcal{X}$ est séparé d'après le Lemme \ref{lemma-well-nigh-affine},
on voit que $f$ est séparé d'après Morphismes de champs, Lemme
\ref{stacks-morphisms-lemma-compose-after-separated}.

\medskip\noindent
Puisque $U \to \mathcal{X}$ est surjectif et que $U \to M$ est quasi-compact,
on voit que $f$ est quasi-compact d'après Morphismes de champs, Lemme
\ref{stacks-morphisms-lemma-surjection-from-quasi-compact}.

\medskip\noindent
D'après Groupoïdes, Lemme \ref{groupoids-lemma-integral-over-invariants},
la composée
$$
U \to \mathcal{X} \to M
$$
est un morphisme intégral de schémas affines. En particulier, elle est
universellement fermée
(Morphismes, Lemme \ref{morphisms-lemma-integral-universally-closed}).
Puisque $U \to \mathcal{X}$ est surjectif, il s'ensuit que $\mathcal{X} \to M$
est universellement fermé (Morphismes de champs, Lemme
\ref{stacks-morphisms-lemma-image-proper-is-proper}).
Pour conclure que $\mathcal{X} \to M$ est un homéomorphisme universel,
il suffit de montrer qu'il est universellement bijectif, c'est-à-dire
surjectif et universellement injectif.

\medskip\noindent
On a $|\mathcal{X}| = |U|/|R|$ d'après
Morphismes de champs, Lemme \ref{stacks-morphisms-lemma-points-presentation}.
Ainsi, $|f|$ est surjectif, et même bijectif,
d'après Groupoïdes, Lemme \ref{groupoids-lemma-points}.

\medskip\noindent
Soit $C \to C'$ un homomorphisme d'anneaux. Soit $(U', R', s', t', c')$
le changement de base de $(U, R, s, t, c)$ par $M' = \Spec(C') \to M$.
En posant $\mathcal{X}' = [U'/R']$, on observe que
$M' \times_M \mathcal{X} = \mathcal{X}'$ d'après
Quotients de groupoïdes, Lemme
\ref{groupoids-quotients-lemma-base-change-quotient-stack}.
Soit $C^1$ l'anneau des fonctions $R'$-invariantes sur $U'$.
Posons $M^1 = \Spec(C^1)$ et considérons le diagramme
$$
\xymatrix{
\mathcal{X}' \ar[d]^{f'} \ar[r] & \mathcal{X} \ar[dd]^f \\
M^1 \ar[d] \\
M' \ar[r] & M
}
$$
D'après Groupoïdes, Lemme \ref{groupoids-lemma-invariants-base-change}, et
Algèbre, Lemme \ref{algebra-lemma-universally-bijective},
le morphisme $M^1 \to M'$ est un homéomorphisme.
D'autre part, le paragraphe précédent appliqué à
$(U', R', s', t', c')$ montre que $|f'|$ est bijectif.
On en déduit que $f$ induit une bijection sur les points après tout
changement de base par un schéma affine. Ainsi, $f$ est universellement injectif
d'après Morphismes de champs, Lemme
\ref{stacks-morphisms-lemma-universally-injective-local}.

\medskip\noindent
Enfin, il reste à montrer que $f$ est un espace de modules catégorique uniforme
dans la catégorie des schémas affines. Cela résulte de la discussion
précédente et du fait que, si
l'homomorphisme d'anneaux $C \to C'$ est plat, alors $C' \to C^1$ est un isomorphisme
d'après Groupoïdes, Lemme \ref{groupoids-lemma-invariants-base-change}.
\end{proof}

\begin{lemma}
\label{lemma-well-nigh-affine-moduli-space-etale}
Soit $h : \mathcal{X}' \to \mathcal{X}$ un morphisme de champs algébriques.
Supposons que $\mathcal{X}'$ et $\mathcal{X}$ soient presque affines,
que $h$ est \'etale et que $h$ induit des isomorphismes entre les groupes d'automorphismes
(Morphismes de champs, remarque
\ref{stacks-morphisms-remark-identify-automorphism-groups}).
Il existe alors un diagramme cartésien
$$
\xymatrix{
\mathcal{X}' \ar[d] \ar[r] & \mathcal{X} \ar[d] \\
M' \ar[r] & M
}
$$
où $M' \to M$ est \'etale et
les flèches verticales sont les espaces de modules construits dans le
lemme \ref{lemma-well-nigh-affine-moduli-space}.
\end{lemma}

\begin{proof}
Remarquons que $h$ est représentable par des espaces algébriques, d'après
Morphismes de champs, lemmes
\ref{stacks-morphisms-lemma-aut-iso-unramified} et
\ref{stacks-morphisms-lemma-stabilizer-preserving}.
Choisissons un schéma affine $U$ et un
morphisme $U \to \mathcal{X}$ surjectif, plat, fini et de présentation finie.
Alors $U' = \mathcal{X}' \times_\mathcal{X} U$ est un espace
algébrique muni d'un morphisme fini (en particulier affine) $U' \to \mathcal{X}'$.
D'après le lemme \ref{lemma-affine-over-well-nigh-affine},
nous en déduisons que $U'$ est affine.
En posant $R = U \times_\mathcal{X} U$ et $R' = U' \times_{\mathcal{X}'} U'$,
nous obtenons des groupoïdes $(U, R, s, t, c)$ et $(U', R', s', t', c')$
tels que $\mathcal{X} = [U/R]$ et $\mathcal{X}' = [U'/R']$ ;
voir la démonstration du lemme \ref{lemma-well-nigh-affine}.
Nous voyons que les diagrammes
$$
\xymatrix{
R' \ar[d]_{s'} \ar[r]_f & R \ar[d]^s \\
U' \ar[r]^f & U
}
\quad
\quad
\xymatrix{
R' \ar[d]_{t'} \ar[r]_f & R \ar[d]^t \\
U' \ar[r]^f & U
}
\quad
\quad
\xymatrix{
G' \ar[d] \ar[r]_f & G \ar[d] \\
U' \ar[r]^f & U
}
$$
sont cartésiens, où $G$ et $G'$ sont les schémas en groupes stabilisateurs.
Pour les deux premiers, cela résulte de la transitivité des produits fibrés,
et, pour le dernier, du fait qu'il est le changement de base de
l'isomorphisme $\mathcal{I}_{\mathcal{X}'} \to
\mathcal{X}' \times_\mathcal{X} \mathcal{I}_\mathcal{X}$
(d'après le lemme déjà utilisé de Morphismes de champs
\ref{stacks-morphisms-lemma-aut-iso-unramified}).
Rappelons que $M$, resp.\ $M'$, a été construit dans le
lemme \ref{lemma-well-nigh-affine-moduli-space}
comme le spectre de l'anneau des fonctions $R$-invariantes sur $U$,
resp.\ de l'anneau des fonctions $R'$-invariantes sur $U'$.
Ainsi, $M' \to M$ est \'etale et $U' = M' \times_M U$
d'après Groupoïdes, lemme \ref{groupoids-lemma-etale}.
Il s'ensuit que $R' = M' \times_M R$ ; autrement dit,
le groupoïde $(U', R', s', t', c')$ est le changement de base de
$(U, R, s, t, c)$ par $M' \to M$.
Cela implique que le diagramme de l'énoncé est
cartésien d'après
Quotients de groupoïdes, lemme
\ref{groupoids-quotients-lemma-base-change-quotient-stack}.
\end{proof}

\begin{lemma}
\label{lemma-moduli-space-finite-affine}
Soit $\mathcal{X}$ un champ algébrique presque affine. Le morphisme
$$
f : \mathcal{X} \longrightarrow M
$$
du lemme \ref{lemma-well-nigh-affine-moduli-space}
est un espace de modules catégorique uniforme.
\end{lemma}

\begin{proof}
Nous savons déjà que $M$ est un espace de modules catégorique uniforme
dans la catégorie des schémas affines. D'après le
lemme \ref{lemma-check-uniform-categorical-quotient-on-affines},
il suffit de montrer que le changement de base
$f' : M' \times_M \mathcal{X} \to M'$
est un espace de modules catégorique pour tout morphisme plat
$M' \to M$ de schémas affines.
Notons que $\mathcal{X}' = M' \times_M \mathcal{X}$ est presque affine d'après le
lemme \ref{lemma-affine-over-well-nigh-affine}.
Ainsi, en remplaçant $\mathcal{X}$ par $\mathcal{X}'$
et $M$ par $M'$, nous nous ramenons à prouver que $f$ est un espace de modules
catégorique.

\medskip\noindent
Soit $g : \mathcal{X} \to Y$ un morphisme, où $Y$ est un espace algébrique.
Nous devons montrer que $g = h \circ f$ pour un unique morphisme $h : M \to Y$.

\medskip\noindent
Unicité. Supposons donnés deux morphismes $h_i : M \to Y$ tels que
$g = h_1 \circ f = h_2 \circ f$. Soit $M' \subset M$ l'égalisateur
de $h_1$ et $h_2$. Alors $M' \to M$ est un monomorphisme et
$f : \mathcal{X} \to M$ se factorise par $M'$. Ainsi, $M' \to M$
est un homéomorphisme universel. Nous en concluons que $M'$ est affine
(Morphismes, lemme \ref{morphisms-lemma-universal-homeomorphism}).
Or, puisque $f : \mathcal{X} \to M$
est un espace de modules catégorique dans la catégorie des schémas affines,
nous avons $M' = M$.

\medskip\noindent
Existence. Nous allons montrer ci-dessous que, pour tout $p \in M$, il existe
un carré cartésien
$$
\xymatrix{
\mathcal{X}' \ar[r] \ar[d] & \mathcal{X} \ar[d] \\
M' \ar[r] & M
}
$$
où $M' \to M$ est un morphisme \'etale de schémas affines dont l'image contient $p$, tel que
la composée $\mathcal{X}' \to \mathcal{X} \to Y$ se factorise par $M'$.
Cela signifie que nous pouvons construire le morphisme $h : M \to Y$ localement pour la topologie \'etale sur $M$.
Puisque $Y$ est un faisceau pour la topologie \'etale et en vertu de l'unicité établie
ci-dessus, cela suffit (nous omettons un petit détail).

\medskip\noindent
Soit $y \in |Y|$ l'image de $p$.
Soit $(V, v) \to (Y, y)$ un morphisme \'etale avec $V$ affine.
Considérons $\mathcal{X}' = V \times_Y \mathcal{X}$.
Remarquons que $\mathcal{X}' \to \mathcal{X}$ est séparé et \'etale
comme changement de base de $V \to Y$. De plus, $\mathcal{X}' \to \mathcal{X}$
induit des isomorphismes entre les groupes d'automorphismes
(Morphismes de champs, remarque
\ref{stacks-morphisms-remark-identify-automorphism-groups})
car cela vaut pour
$V \to Y$ ; voir Morphismes de champs, lemme
\ref{stacks-morphisms-lemma-base-change-stabilizer-preserving}.
Choisissons une présentation $\mathcal{X} = [U/R]$
comme dans le lemme \ref{lemma-well-nigh-affine}.
Posons $U' = \mathcal{X}' \times_\mathcal{X} U = V \times_Y U$
et choisissons $u' \in U'$ se projetant sur $p$ et $v$ (ce qui est possible
d'après Propriétés des espaces, lemme \ref{spaces-properties-lemma-points-cartesian}).
Puisque $U' \to U$ est séparé et \'etale, nous voyons que
tout ensemble fini de points de $U'$ est contenu dans un ouvert affine ; voir
Compléments sur les morphismes, lemme
\ref{more-morphisms-lemma-separated-locally-quasi-finite-over-affine}.
D'autre part, le morphisme $U' \to \mathcal{X}'$ est
surjectif, fini, plat et localement de présentation finie.
En posant $R' = U' \times_{\mathcal{X}'} U'$, nous voyons
que $s', t' : R' \to U'$ sont finis localement libres.
D'après Groupoïdes, lemme \ref{groupoids-lemma-find-invariant-affine},
il existe un sous-schéma ouvert affine invariant par $R'$, noté $U'' \subset U'$,
contenant $u'$.
Soit $\mathcal{X}'' \subset \mathcal{X}'$ le
sous-champ ouvert correspondant. Alors $\mathcal{X}''$ est
presque affine. D'après le lemme \ref{lemma-well-nigh-affine-moduli-space-etale},
nous obtenons un carré cartésien
$$
\xymatrix{
\mathcal{X}'' \ar[r] \ar[d] & \mathcal{X} \ar[d] \\
M'' \ar[r] & M
}
$$
où $M'' \to M$ est \'etale. Puisque $\mathcal{X}'' \to M''$ est
un espace de modules catégorique dans la catégorie des schémas affines,
nous obtenons un morphisme $M'' \to V$ tel que la composée
$\mathcal{X}'' \to \mathcal{X}' \to V$ est égale à la composée
$\mathcal{X}'' \to M'' \to V$. Cela prouve notre assertion et achève
la démonstration.
\end{proof}

\begin{lemma}
\label{lemma-etale-separated-over-well-nigh-affine}
Soit $h : \mathcal{X}' \to \mathcal{X}$ un morphisme de champs algébriques.
Supposons que $\mathcal{X}$ soit presque affine, que $h$ soit \'etale, que $h$ soit séparé
et que $h$ induit des isomorphismes entre les groupes d'automorphismes
(Morphismes de champs, remarque
\ref{stacks-morphisms-remark-identify-automorphism-groups}).
Il existe alors un diagramme cartésien
$$
\xymatrix{
\mathcal{X}' \ar[d] \ar[r] & \mathcal{X} \ar[d] \\
M' \ar[r] & M
}
$$
où $M' \to M$ est un morphisme de schémas séparé et \'etale, et où
$\mathcal{X} \to M$ est l'espace de modules construit dans le
lemme \ref{lemma-well-nigh-affine-moduli-space}.
\end{lemma}

\begin{proof}
Choisissons un schéma affine $U$ et un morphisme surjectif, plat, fini et
localement de présentation finie $U \to \mathcal{X}$.
Puisque $h$ est représentable par des espaces algébriques
(Morphismes de champs, lemmes
\ref{stacks-morphisms-lemma-aut-iso-unramified} et
\ref{stacks-morphisms-lemma-stabilizer-preserving})
nous voyons que $U' = \mathcal{X}' \times_\mathcal{X} U$ est
un espace algébrique. Puisque $U' \to U$ est séparé et \'etale,
nous voyons que $U'$ est un schéma et que tout ensemble fini de points
de $U'$ est contenu dans un ouvert affine ; voir
Morphismes d'espaces, lemme
\ref{spaces-morphisms-lemma-locally-quasi-finite-separated-representable}
et
Compléments sur les morphismes, lemme
\ref{more-morphisms-lemma-separated-locally-quasi-finite-over-affine}.
En posant $R' = U' \times_{\mathcal{X}'} U'$, nous voyons
que $s', t' : R' \to U'$ sont finis localement libres.
D'après Groupoïdes, lemme \ref{groupoids-lemma-find-invariant-affine},
il existe un recouvrement ouvert $U' = \bigcup U'_i$ par
des sous-schémas ouverts affines invariants par $R'$, $U'_i \subset U'$.
Soient $\mathcal{X}'_i \subset \mathcal{X}'$ les
sous-champs ouverts correspondants. Ils sont presque affines, car $U'_i \to \mathcal{X}'_i$
est surjectif, plat, fini et de présentation finie. D'après le
lemme \ref{lemma-well-nigh-affine-moduli-space-etale},
nous obtenons des diagrammes cartésiens
$$
\xymatrix{
\mathcal{X}'_i \ar[r] \ar[d] & \mathcal{X} \ar[d] \\
M'_i \ar[r] & M
}
$$
où $M'_i \to M$ est un morphisme \'etale de schémas affines,
les flèches verticales étant comme dans le
lemme \ref{lemma-well-nigh-affine-moduli-space}.
Remarquons que
$\mathcal{X}'_{ij} = \mathcal{X}'_i \cap \mathcal{X}'_j$
est un sous-champ ouvert de $\mathcal{X}'_i$ et de $\mathcal{X}'_j$.
Nous obtenons donc les sous-schémas ouverts correspondants
$V_{ij} \subset M'_i$ et $V_{ji} \subset M'_j$.
En vertu du
lemme \ref{lemma-moduli-space-finite-affine},
nous voyons que
$\mathcal{X}'_{ij} \to V_{ij}$ et
$\mathcal{X}'_{ji} \to V_{ji}$ sont tous deux des espaces de modules catégoriques !
Nous obtenons ainsi un unique isomorphisme $\varphi_{ij} : V_{ij} \to V_{ji}$
tel que le diagramme
$$
\xymatrix{
\mathcal{X}'_i \ar[d] & &
\mathcal{X}'_i \cap \mathcal{X}'_j \ar[rr] \ar[ll] \ar[ld] \ar[rd] & &
\mathcal{X}'_j \ar[d] \\
M'_i &
V_{ij} \ar[l] \ar[rr]^{\varphi_{ij}} & &
V_{ji} \ar[r] &
M'_j
}
$$
est commutatif. Ces isomorphismes satisfont à la condition de cocycle de
Schémas, section \ref{schemes-section-glueing-schemes}, comme le montre un calcul
(et une nouvelle application du lemme précédent) que nous omettons.
Nous pouvons donc recoller les schémas affines pour former un schéma $M'$ ; voir
Schémas, lemme \ref{schemes-lemma-glue}.
Identifions les $M'_i$ à leur image dans $M'$.
Nous affirmons qu'il existe un morphisme $\mathcal{X}' \to M'$ s'inscrivant dans des
diagrammes cartésiens
$$
\xymatrix{
\mathcal{X}'_i \ar[r] \ar[d] & \mathcal{X}' \ar[d] \\
M'_i \ar[r] & M'
}
$$
Cela ressort de la description des morphismes vers le schéma recollé $M'$
dans Schémas, lemme \ref{schemes-lemma-glue}, et du fait que donner un morphisme
$\mathcal{X}' \to M'$ revient à donner un morphisme $T \to M'$
pour tout morphisme $T \to \mathcal{X}'$.
De même, il existe un morphisme $M' \to M$ dont les restrictions sont les
morphismes donnés $M'_i \to M$ sur $M'_i$.
Le morphisme $M' \to M$ est \'etale (car il l'est sur les membres d'un
recouvrement \'etale), et la propriété de produit fibré vaut puisqu'elle peut
se vérifier sur les membres du recouvrement ouvert (affine) $M' = \bigcup M'_i$.
Enfin, $M' \to M$ est séparé, car la composée
$U' \to \mathcal{X}' \to M'$ est surjective et universellement fermée,
et nous pouvons appliquer Morphismes, lemme
\ref{morphisms-lemma-image-universally-closed-separated}.
\end{proof}

\begin{lemma}
\label{lemma-etale-local-finite-inertia}
Soit $\mathcal{X}$ un champ algébrique. Supposons que
$\mathcal{I}_\mathcal{X} \to \mathcal{X}$ est fini.
Il existe alors un ensemble $I$ et, pour $i \in I$, un morphisme de champs algébriques
$$
g_i : \mathcal{X}_i \longrightarrow \mathcal{X}
$$
ayant les propriétés suivantes :
\begin{enumerate}
\item $|\mathcal{X}| = \bigcup |g_i|(|\mathcal{X}_i|)$,
\item $\mathcal{X}_i$ est presque affine,
\item $\mathcal{I}_{\mathcal{X}_i} \to
\mathcal{X}_i \times_\mathcal{X} \mathcal{I}_\mathcal{X}$
est un isomorphisme, et
\item $g_i : \mathcal{X}_i \to \mathcal{X}$ est représentable
par des espaces algébriques, séparé et \'etale,
\end{enumerate}
\end{lemma}

\begin{proof}
Pour tout $x \in |\mathcal{X}|$, nous pouvons choisir
$g : \mathcal{U} \to \mathcal{X}$, $\mathcal{U} = [U/R]$ et $u$ comme dans
Morphismes de champs,
lemme \ref{stacks-morphisms-lemma-etale-local-quasi-DM-at-x}.
Puis, d'après
Morphismes de champs, lemme
\ref{stacks-morphisms-lemma-stabilizer-preserving-unramified}
nous voyons qu'il existe un sous-champ ouvert
$\mathcal{U}' \subset \mathcal{U}$ contenant $u$
tel que $\mathcal{I}_{\mathcal{U}'} \to
\mathcal{U}' \times_\mathcal{X} \mathcal{I}_\mathcal{X}$
est un isomorphisme.
Soit $U' \subset U$ l'ouvert invariant par $R$ correspondant au
sous-champ ouvert $\mathcal{U}'$.
Soit $u' \in U'$ un point de $U'$ s'envoyant sur $u$.
Remarquons que $t(s^{-1}(\{u'\}))$ est fini, puisque $s : R \to U$ est fini.
D'après Propriétés, lemme \ref{properties-lemma-ample-finite-set-in-affine}
et Groupoïdes, lemme \ref{groupoids-lemma-find-invariant-affine},
nous pouvons trouver un ouvert affine invariant par $R$, noté $U'' \subset U'$,
contenant $u'$. Soit $R''$ la restriction de $R$ à $U''$.
Alors $\mathcal{U}'' = [U''/R'']$ est un sous-champ ouvert de
$\mathcal{U}'$ contenant $u$, est presque affine,
$\mathcal{I}_{\mathcal{U}''} \to
\mathcal{U}'' \times_\mathcal{X} \mathcal{I}_\mathcal{X}$
est un isomorphisme, et $\mathcal{U}'' \to \mathcal{X}$
est représentable par des espaces algébriques et \'etale.
Enfin, $\mathcal{U}'' \to \mathcal{X}$ est séparé, car
$\mathcal{U}''$ est séparé (lemme \ref{lemma-well-nigh-affine}),
la diagonale de $\mathcal{X}$ est séparée
(Morphismes de champs, lemme \ref{stacks-morphisms-lemma-diagonal-diagonal}),
et la séparation résulte de Morphismes de champs, lemme
\ref{stacks-morphisms-lemma-compose-after-separated}.
Puisque le point $x \in |\mathcal{X}|$ est arbitraire, la démonstration est achevée.
\end{proof}

\begin{theorem}[Keel-Mori]
\label{theorem-keel-mori}
Soit $\mathcal{X}$ un champ algébrique. Supposons que
$\mathcal{I}_\mathcal{X} \to \mathcal{X}$ est fini.
Il existe alors un espace de modules catégorique uniforme
$$
f : \mathcal{X} \longrightarrow M
$$
et $f$ est séparé, quasi-compact et un homéomorphisme universel.
\end{theorem}

\begin{proof}
Choisissons un ensemble $I$\footnote{Le lecteur qui tient encore
compte des questions ensemblistes devra s'assurer que $I$ n'est pas trop grand.}
et, pour $i \in I$, un morphisme de champs algébriques
$g_i : \mathcal{X}_i \to \mathcal{X}$ comme dans
le lemme \ref{lemma-etale-local-finite-inertia} ; nous utiliserons sans autre mention toutes
les propriétés énumérées dans ce lemme.
Soit
$$
f_i : \mathcal{X}_i \to M_i
$$
le morphisme donné par le lemme \ref{lemma-well-nigh-affine-moduli-space}.
Considérons les champs
$$
\mathcal{X}_{ij} = \mathcal{X}_i \times_{g_i, \mathcal{X}, g_j} \mathcal{X}_j
$$
pour $i, j \in I$. Les projections $\mathcal{X}_{ij} \to \mathcal{X}_i$
et $\mathcal{X}_{ij} \to \mathcal{X}_j$ sont séparées
d'après Morphismes de champs, lemme
\ref{stacks-morphisms-lemma-base-change-separated},
\'etales d'après Morphismes de champs, lemme
\ref{stacks-morphisms-lemma-base-change-etale},
et induisent des isomorphismes sur les groupes d'automorphismes
(comme dans Morphismes de champs, remarque
\ref{stacks-morphisms-remark-identify-automorphism-groups}), d'après
Morphismes de champs, lemme
\ref{stacks-morphisms-lemma-base-change-stabilizer-preserving}.
Nous pouvons donc appliquer le lemme \ref{lemma-etale-separated-over-well-nigh-affine}
afin d'obtenir un diagramme commutatif
$$
\xymatrix{
\mathcal{X}_i \ar[d]_{f_i} &
\mathcal{X}_{ij} \ar[d]_{f_{ij}} \ar[l] \ar[r] &
\mathcal{X}_j \ar[d]_{f_j} \\
M_i &
M_{ij} \ar[l] \ar[r] &
M_j
}
$$
dont les carrés sont cartésiens et où $M_{ij} \to M_i$ et $M_{ij} \to M_j$
sont des morphismes de schémas séparés et \'etales ; nous utilisons ici aussi que $f_i$
est un quotient catégorique uniforme d'après le
lemme \ref{lemma-moduli-space-finite-affine}.
Affirmation :
$$
\coprod M_{ij} \longrightarrow \coprod M_i \times \coprod M_i
$$
est une relation d'équivalence \'etale.

\medskip\noindent
Démonstration de l'affirmation. Posons $R = \coprod M_{ij}$ et $U = \coprod M_i$.
Nous avons déjà vu que $t : R \to U$ et $s : R \to U$ sont \'etales.
Construisons un morphisme $c : R \times_{s, U, t} R \to R$
compatible avec $\text{pr}_{13} : U \times U \times U \to U \times U$.
Plus précisément, pour $i, j, k \in I$, considérons
$$
\mathcal{X}_{ijk} =
\mathcal{X}_i \times_{g_i, \mathcal{X}, g_j} \mathcal{X}_j
\times_{g_j, \mathcal{X}, g_k} \mathcal{X}_k =
\mathcal{X}_{ij} \times_{\mathcal{X}_j} \mathcal{X}_{jk}
$$
En raisonnant exactement comme au paragraphe précédent,
nous voyons que $M_{ijk} = M_{ij} \times_{M_j} M_{jk}$
est un espace de modules catégorique pour $\mathcal{X}_{ijk}$.
En particulier, il existe un morphisme canonique
$M_{ijk} = M_{ij} \times_{M_j} M_{jk} \to M_{ik}$
provenant de la projection $\mathcal{X}_{ijk} \to \mathcal{X}_{ik}$.
En réunissant ces morphismes, nous obtenons le morphisme $c$.
De même, nous construisons un morphisme $e : U \to R$
compatible avec $\Delta : U \to U \times U$ et
$i : R \to R$ compatible avec la permutation des facteurs $U \times U \to U \times U$.
Soit $k$ un corps algébriquement clos. Alors
$$
\Mor(\Spec(k), \mathcal{X}_i) \to \Mor(\Spec(k), M_i) = M_i(k)
$$
est bijective sur les classes d'isomorphie, et il en reste de même après tout
changement de base par un morphisme $M' \to M_i$. Cela résulte du choix
de $f_i$ et de Morphismes de champs, lemmes
\ref{stacks-morphisms-lemma-universally-injective} et
\ref{stacks-morphisms-lemma-universally-injective-point}.
Par construction des $2$-produits fibrés, le diagramme
$$
\xymatrix{
\Mor(\Spec(k), \mathcal{X}_{ij}) \ar[d] \ar[r] &
\Mor(\Spec(k), \mathcal{X}_j) \ar[d] \\
\Mor(\Spec(k), \mathcal{X}_i) \ar[r] &
\Mor(\Spec(k), \mathcal{X})
}
$$
est un 2-produit fibré de catégories. D'après notre choix de $g_i$, les
foncteurs de ce diagramme induisent des bijections sur les groupes d'automorphismes.
Il s'ensuit que ce diagramme induit un diagramme cartésien
sur les ensembles de classes d'isomorphie ! Nous voyons donc que
$$
R(k) = U(k) \times_{|\Mor(\Spec(k), \mathcal{X})|} U(k)
$$
où $|\Mor(\Spec(k), \mathcal{X})|$ désigne l'ensemble
des classes d'isomorphie.
En particulier, pour tout corps algébriquement clos $k$,
l'application sur les points à valeurs dans $k$ définit une relation d'équivalence.
L'affirmation en résulte par
Groupoïdes, lemme \ref{groupoids-lemma-etale-equivalence-relation}.

\medskip\noindent
Soit $M = U/R$ l'espace algébrique quotient de la
relation d'équivalence \'etale ci-dessus ; voir
Espaces, théorème \ref{spaces-theorem-presentation}.
Il existe un morphisme canonique $f : \mathcal{X} \to M$
s'insérant dans les diagrammes commutatifs
\begin{equation}
\label{equation-fundamental-diagram}
\xymatrix{
\mathcal{X}_i \ar[r]_{g_i} \ar[d]_{f_i} & \mathcal{X} \ar[d]^f \\
M_i \ar[r] & M
}
\end{equation}
En effet, un tel morphisme $f$ est défini par un foncteur
$$
f : \Mor(T, \mathcal{X}) \longrightarrow \Mor(T, M)
$$
pour tout schéma $T$, la famille ainsi obtenue étant compatible avec les changements de base. Soit $a : T \to \mathcal{X}$
un objet du membre de gauche. Nous obtenons un recouvrement \'etale
$\{T_i \to T\}$, avec $T_i = \mathcal{X}_i \times_\mathcal{X} T$,
et des morphismes $a_i : T_i \to \mathcal{X}_i$. Nous obtenons alors
$b_i = f_i \circ a_i : T_i \to M_i$. Puisque
$T_i \times_T T_j = \mathcal{X}_{ij} \times_\mathcal{X} T$,
nous obtenons de plus un morphisme $a_{ij} : T_i \times_T T_j \to \mathcal{X}_{ij}$.
En posant $b_{ij} = f_{ij} \circ a_{ij}$, nous constatons que
$b_i \times b_j$ se factorise par le monomorphisme
$M_{ij} \to M_i \times M_j$. Par conséquent, les morphismes
$$
T_i \xrightarrow{b_i} M_i \to M
$$
coïncident sur $T_i \times_T T_j$. Comme $M$ est un faisceau pour la
topologie \'etale, nous voyons que ces morphismes se recollent en un unique
morphisme $b = f(a) : T \to M$. Nous omettons de vérifier que
cette construction est compatible avec les changements de base, ainsi que de vérifier
que les diagrammes
(\ref{equation-fundamental-diagram}) sont commutatifs.

\medskip\noindent
Affirmation : les diagrammes (\ref{equation-fundamental-diagram}) sont cartésiens.
Pour le voir, étudions le morphisme induit
$$
h_i : \mathcal{X}_i \longrightarrow M_i \times_M \mathcal{X}
$$
C'est un morphisme de champs \'etale sur $\mathcal{X}$,
donc $h_i$ est \'etale (Morphismes de champs, lemme
\ref{stacks-morphisms-lemma-etale-permanence}).
Comme $g_i$ est séparé, nous voyons que $h_i$ est séparé
(utiliser Morphismes de champs, lemme
\ref{stacks-morphisms-lemma-compose-after-separated}, et le fait,
vu plus haut, que la diagonale de $\mathcal{X}$ est séparée).
Le morphisme $h_i$ induit des isomorphismes sur les groupes d'automorphismes
(Morphismes de champs, remarque
\ref{stacks-morphisms-remark-identify-automorphism-groups})
puisqu'il en est ainsi pour $g_i$. Pour un corps algébriquement clos $k$,
le diagramme
$$
\xymatrix{
\Mor(\Spec(k), M_i \times_M \mathcal{X}) \ar[r] \ar[d] &
\Mor(\Spec(k), \mathcal{X}) \ar[d] \\
M_i(k) \ar[r] &
M(k)
}
$$
est un diagramme cartésien de catégories, et la flèche du haut
induit des bijections sur les groupes d'automorphismes.
D'autre part, nous avons
$$
M(k) = U(k)/R(k) = U(k)/
U(k) \times_{|\Mor(\Spec(k), \mathcal{X})|} U(k) =
|\Mor(\Spec(k), \mathcal{X})|
$$
d'après ce qui précède. Ainsi, la flèche verticale de droite du
diagramme cartésien ci-dessus est bijective sur les classes d'isomorphie.
Nous concluons que
$|\Mor(\Spec(k), M_i \times_M \mathcal{X})| \to M_i(k)$ est bijective.
Récapitulons : $h_i$ est séparé et \'etale, il induit des isomorphismes sur les
groupes d'automorphismes (comme dans Morphismes de champs, remarque
\ref{stacks-morphisms-remark-identify-automorphism-groups}),
et une équivalence sur toute catégorie fibre au-dessus d'un corps algébriquement clos.
Il s'agit donc d'un isomorphisme d'après Morphismes de champs, lemme
\ref{stacks-morphisms-lemma-etale-iso}.

\medskip\noindent
L'affirmation donne en particulier ceci :
il existe un morphisme \'etale surjectif $U \to M$
tel que le changement de base de $f$ est séparé, quasi-compact
et un homéomorphisme universel. Il s'ensuit que $f$ est séparé,
quasi-compact et un homéomorphisme universel.
Voir Morphismes de champs, lemme
\ref{stacks-morphisms-lemma-check-separated-covering},
\ref{stacks-morphisms-lemma-check-quasi-compact-covering}, et
\ref{stacks-morphisms-lemma-check-universal-homeomorphism-covering}.

\medskip\noindent
Pour achever la démonstration, il reste à montrer que $f : \mathcal{X} \to M$
est un espace de modules catégorique uniforme.
Pour cela, il suffit de montrer que, pour tout morphisme plat
$M' \to M$ d'espaces algébriques, le changement de base
$$
M' \times_M \mathcal{X} \longrightarrow M'
$$
est un espace de modules catégorique. Considérons donc un morphisme
$$
\theta : M' \times_M \mathcal{X} \longrightarrow E
$$
où $E$ est un espace algébrique. Pour chaque $i$, nous savons que
$f_i$ est un espace de modules catégorique uniforme. Nous obtenons donc
$$
\xymatrix{
M' \times_M \mathcal{X}_i \ar[d] \ar[r] &
M' \times_M \mathcal{X} \ar[d]^\theta \\
M' \times_M M_i \ar[r]^{\psi_i} &
E
}
$$
Puisque $\{M' \times_M M_i \to M'\}$ est un recouvrement \'etale,
pour obtenir le morphisme souhaité $\psi : M' \to E$, il suffit
de montrer que $\psi_i$ et $\psi_j$ coïncident sur
$M' \times_M M_i \times_M M_j = M' \times_M M_{ij}$.
Cela résulte aussitôt du fait que
$f_{ij} : \mathcal{X}_{ij} =
\mathcal{X}_i \times_\mathcal{X} \mathcal{X}_j \to M_{ij}$ est un quotient catégorique
uniforme ; nous omettons les détails.
Enfin, on montre que $\psi$ s'insère dans le diagramme commutatif
$$
\xymatrix{
M' \times_M \mathcal{X} \ar[d] \ar[rd]^\theta \\
M' \ar[r]^\psi &
E
}
$$
car « $\{M' \times_M \mathcal{X}_i \to M' \times_M \mathcal{X}\}$
est un recouvrement \'etale » et que les morphismes $\psi_i$ s'insèrent dans les
diagrammes commutatifs correspondants par construction.
Ceci achève la démonstration du théorème de Keel--Mori.
\end{proof}

\noindent
Le lemme suivant souligne le caractère local pour la topologie \'etale de la construction.

\begin{lemma}
\label{lemma-etale-separated-over-keel-mori}
Soit $h : \mathcal{X}' \to \mathcal{X}$ un morphisme de champs algébriques.
Supposons que
\begin{enumerate}
\item $\mathcal{I}_\mathcal{X} \to \mathcal{X}$ est fini,
\item $h$ est \'etale et séparé, et induit des isomorphismes sur les
groupes d'automorphismes (Morphismes de champs, remarque
\ref{stacks-morphisms-remark-identify-automorphism-groups}).
\end{enumerate}
Il existe alors un diagramme cartésien
$$
\xymatrix{
\mathcal{X}' \ar[d] \ar[r] &
\mathcal{X} \ar[d] \\
M' \ar[r] &
M
}
$$
où $M' \to M$ est un morphisme d'espaces algébriques séparé et \'etale, et
les flèches verticales sont les espaces de modules construits dans le
théorème \ref{theorem-keel-mori}.
\end{lemma}

\begin{proof}
D'après Morphismes de champs, lemme \ref{stacks-morphisms-lemma-aut-iso-unramified},
nous voyons que
$\mathcal{I}_{\mathcal{X}'} \to
\mathcal{X}' \times_\mathcal{X} \mathcal{I}_\mathcal{X}$
est un isomorphisme. Ainsi, $\mathcal{I}_{\mathcal{X}'} \to \mathcal{X}'$
est fini comme changement de base de $\mathcal{I}_\mathcal{X} \to \mathcal{X}$.
Soient $f' : \mathcal{X}' \to M'$ et $f : \mathcal{X} \to M$ ceux du
théorème \ref{theorem-keel-mori}.
Nous obtenons un diagramme commutatif comme dans le lemme, car
$f'$ est un espace de modules catégorique.
Choisissons $I$ et $g'_i : \mathcal{X}'_i \to \mathcal{X}'$ comme dans le
lemme \ref{lemma-etale-local-finite-inertia}.
Remarquons que $g_i = h \circ g'_i$
est \'etale et séparé, et induit des isomorphismes sur les
groupes d'automorphismes (Morphismes de champs, remarque
\ref{stacks-morphisms-remark-identify-automorphism-groups}).
Soit $f'_i : \mathcal{X}'_i \to M'_i$ le morphisme du
lemme \ref{lemma-well-nigh-affine-moduli-space}.
Dans la démonstration du théorème \ref{theorem-keel-mori},
nous avons vu que les diagrammes
$$
\xymatrix{
\mathcal{X}'_i \ar[d]_{f'_i} \ar[r]_{g'_i} &
\mathcal{X}' \ar[d]^{f'} \\
M'_i \ar[r] &
M'
}
\quad\text{et}\quad
\xymatrix{
\mathcal{X}'_i \ar[d]_{f'_i} \ar[r]_{g_i} &
\mathcal{X} \ar[d]^f \\
M'_i \ar[r] &
M
}
$$
sont cartésiens et que $M'_i \to M'$ et $M'_i \to M$ sont \'etales
(cela résulte aussi directement des propriétés des morphismes
$g'_i, g_i, f', f'_i, f$ énumérées jusqu'ici, en raisonnant exactement de la même façon).
Nous en déduisons d'une part que $M' \to M$ est \'etale et d'autre part que
le diagramme du lemme est cartésien. Il reste à montrer
que $M' \to M$ est séparé. Pour cela, considérons le
diagramme
$$
\xymatrix{
\mathcal{X}' \ar[r] \ar[d] &
\mathcal{X}' \times_\mathcal{X} \mathcal{X}' \ar[d] \\
M' \ar[r] &
M' \times_M M'
}
$$
La flèche horizontale supérieure est universellement fermée, car
$\mathcal{X}' \to \mathcal{X}$ est séparé.
Les flèches verticales sont celles du théorème \ref{theorem-keel-mori}
(en tant que changements de base plats de $\mathcal{X} \to M$)
et sont donc des homéomorphismes universels. La flèche horizontale inférieure
est par conséquent universellement fermée. Comme elle est en outre un monomorphisme
\'etale d'espaces algébriques, c'est l'immersion fermée voulue.
\end{proof}





\section{Propriétés des espaces de modules}
\label{section-properties-moduli-spaces}

\noindent
Une fois établie l'existence d'un espace de modules,
il devient possible (et généralement aisé) d'établir
des propriétés de ces espaces de modules.

\begin{lemma}
\label{lemma-keel-mori-finite-type}
Soit $p : \mathcal{X} \to Y$ un morphisme d'un champ algébrique vers un
espace algébrique. Supposons que
\begin{enumerate}
\item $\mathcal{I}_\mathcal{X} \to \mathcal{X}$ est fini,
\item $Y$ est localement noethérien, et
\item $p$ est localement de type fini.
\end{enumerate}
Soit $f : \mathcal{X} \to M$ l'espace de modules construit dans le
théorème \ref{theorem-keel-mori}.
Alors $M \to Y$ est localement de type fini.
\end{lemma}

\begin{proof}
Puisque $f$ est un espace de modules catégorique uniforme, nous obtenons le
morphisme $M \to Y$. Il suffit de vérifier que $M \to Y$
est localement de type fini, localement pour la topologie \'etale sur $M$ et $Y$.
Puisque $f$ est un espace de modules catégorique uniforme, nous
pouvons d'abord remplacer $Y$ par un schéma affine \'etale sur $Y$.
Ensuite, nous pouvons choisir $I$ et $g_i : \mathcal{X}_i \to \mathcal{X}$
comme dans le lemme \ref{lemma-etale-local-finite-inertia}.
Alors, par le lemme \ref{lemma-etale-separated-over-keel-mori},
nous nous ramenons au cas $\mathcal{X} = \mathcal{X}_i$.
Autrement dit, nous pouvons supposer que $\mathcal{X}$ est presque affine.
Dans ce cas, nous avons $Y = \Spec(A_0)$,
$\mathcal{X} = [U/R]$ avec $U = \Spec(A)$ et
$M = \Spec(C)$, où $C \subset A$ est l'ensemble des fonctions invariantes par $R$
sur $U$. Voir les
lemmes \ref{lemma-well-nigh-affine} et
\ref{lemma-well-nigh-affine-moduli-space}.
Alors $A_0$ est noethérien et $A_0 \to A$ est de type fini.
De plus, $A$ est entier sur $C$ d'après
Groupoïdes, lemme \ref{groupoids-lemma-integral-over-invariants},
donc fini sur $C$
(cet anneau étant de type fini sur $A_0$).
Nous pouvons donc enfin appliquer
Algèbre, lemme \ref{algebra-lemma-Artin-Tate}
pour conclure.
\end{proof}

\begin{lemma}
\label{lemma-keel-mori-diagonal}
Soit $\mathcal{X}$ un champ algébrique. Supposons que
$\mathcal{I}_\mathcal{X} \to \mathcal{X}$ est fini.
Soit $f : \mathcal{X} \to M$ l'espace de modules construit dans le
théorème \ref{theorem-keel-mori}.
\begin{enumerate}
\item Si $\mathcal{X}$ est quasi-séparé, alors $M$ est quasi-séparé.
\item Si $\mathcal{X}$ est séparé, alors $M$ est séparé.
\item Ajouter ici d'autres énoncés, par exemple des versions relatives de ceux qui précèdent.
\end{enumerate}
\end{lemma}

\begin{proof}
Pour le prouver, considérons le diagramme suivant
$$
\xymatrix{
\mathcal{X} \ar[d]_f \ar[r]_{\Delta_\mathcal{X}} &
\mathcal{X} \times \mathcal{X} \ar[d]^{f \times f} \\
M \ar[r]^{\Delta_M} &
M \times M
}
$$
Puisque $f$ est un homéomorphisme universel, nous voyons que
$f \times f$ est un homéomorphisme universel.

\medskip\noindent
Si $\mathcal{X}$ est séparé, alors $\Delta_\mathcal{X}$ est propre,
donc $\Delta_\mathcal{X}$ est universellement fermé ; par conséquent,
$\Delta_M$ est universellement fermé, et $M$ est donc séparé
d'après Morphismes d'espaces, lemme
\ref{spaces-morphisms-lemma-separated-diagonal-proper}.

\medskip\noindent
Si $\mathcal{X}$ est quasi-séparé, alors $\Delta_\mathcal{X}$ est
quasi-compact ; par conséquent, $\Delta_M$ est quasi-compact, et $M$ est donc
quasi-séparé.
\end{proof}

\begin{lemma}
\label{lemma-keel-mori-proper}
Soit $p : \mathcal{X} \to Y$ un morphisme d'un champ algébrique
vers un espace algébrique. Supposons que
\begin{enumerate}
\item $\mathcal{I}_\mathcal{X} \to \mathcal{X}$ est fini,
\item $p$ est propre, et
\item $Y$ est localement noethérien.
\end{enumerate}
Soit $f : \mathcal{X} \to M$ l'espace de modules construit dans le
théorème \ref{theorem-keel-mori}. Alors $M \to Y$ est propre.
\end{lemma}

\begin{proof}
D'après le lemme \ref{lemma-keel-mori-finite-type},
nous voyons que $M \to Y$ est localement de type fini.
D'après le lemme \ref{lemma-keel-mori-diagonal}, nous voyons que
$M \to Y$ est séparé.
Bien entendu, $M \to Y$ est quasi-compact et universellement fermé,
car ces propriétés sont topologiques, le morphisme $\mathcal{X} \to Y$
les possède et $\mathcal{X} \to M$ est un
homéomorphisme universel.
\end{proof}





\section{Champs et recouvrements fpqc}
\label{section-fpqc}

\noindent
Certains champs algébriques satisfont à la descente fpqc.
L'analogue de cette section pour les espaces algébriques
est Propriétés des espaces, section \ref{spaces-properties-section-fpqc}.


\begin{proposition}
\label{proposition-stack-fpqc}
\begin{reference}
Proposition 3.3.6 de « Intro to Algebraic Stacks », par
Anatoly Preygel.
\end{reference}
Soit $\mathcal{X}$ un champ algébrique dont la diagonale est quasi-affine\footnote{Il suffit
de la supposer ind-quasi-affine.}. Alors
$\mathcal{X}$ satisfait à la descente pour les recouvrements fpqc.
\end{proposition}

\begin{proof}
Selon nos conventions, $\mathcal{X}$ est un champ
en groupoïdes $p : \mathcal{X} \to (\Sch/S)_{fppf}$
sur la catégorie des schémas sur un schéma de base $S$, munie
de la topologie fppf. L'énoncé signifie ce qui suit :
étant donné un recouvrement fpqc $\mathcal{U} = \{U_i \to U\}_{i \in I}$
de schémas sur $S$, le foncteur
$$
\mathcal{X}_U \longrightarrow DD(\mathcal{U})
$$
est une équivalence. À gauche se trouve la catégorie
des objets de $\mathcal{X}$ au-dessus de $U$, et à droite
la catégorie des données de descente dans $\mathcal{X}$ relativement à $\mathcal{U}$.
Voir l'explication dans Champs, section \ref{stacks-section-descent-data}.

\medskip\noindent
Pleine fidélité. Supposons que nous ayons deux objets $x, y$ de $\mathcal{X}$
au-dessus de $U$. Alors $J = \mathit{Isom}(x, y)$ est un espace algébrique sur $U$.
Par conséquent, une famille de sections de $J$ sur les $U_i$ dont les restrictions
à $U_i \times_U U_j$ coïncident provient d'une unique section sur $U$, d'après
l'analogue de la proposition pour les espaces algébriques ; voir
Propriétés des espaces, proposition
\ref{spaces-properties-proposition-sheaf-fpqc}.
Notre foncteur est donc pleinement fidèle.

\medskip\noindent
Surjectivité essentielle. On se donne ici des objets
$x_i$ au-dessus de $U_i$ et des isomorphismes
$\varphi_{ij} : \text{pr}_0^*x_i \to \text{pr}_1^*x_j$
sur $U_i \times_U U_j$ satisfaisant à la condition de cocycle
sur $U_i \times_U U_j \times_U U_k$.

\medskip\noindent
Soit $W$ un schéma affine et soit $W \to \mathcal{X}$
un morphisme. Pour chaque $i$, nous pouvons former
$$
W_i = U_i \times_{x_i, \mathcal{X}} W
$$
La projection $W_i \to U_i$ est quasi-affine, car la
diagonale de $\mathcal{X}$ est quasi-affine. Pour toute paire
$i, j \in I$, l'isomorphisme $\varphi_{ij}$ induit
un isomorphisme
$$
W_i \times_U U_j =
(U_i \times_U U_j) \times_{x_i \circ \text{pr}_0, \mathcal{X}} W
\to
(U_i \times_U U_j) \times_{x_j \circ \text{pr}_1, \mathcal{X}} W =
U_i \times_U W_j
$$
De plus, ces isomorphismes satisfont à la condition de cocycle sur
$U_i \times_U U_j \times_U U_k$. Autrement dit, ils
définissent une donnée de descente sur les schémas $W_i/U_i$ relativement à $\mathcal{U}$.
D'après Descente, lemme \ref{descent-lemma-quasi-affine},
nous voyons que cette donnée de descente est effective\footnote{Ou bien utiliser
Compléments sur les groupoïdes, lemme \ref{more-groupoids-lemma-ind-quasi-affine},
dans le cas d'une diagonale ind-quasi-affine.}.
Nous en concluons qu'il existe un morphisme quasi-affine
$W' \to U$ et un diagramme commutatif
$$
\xymatrix{
W' \ar[d] &
W_i \ar[l] \ar[d] \ar[r] &
W \ar[d] \\
U &
U_i \ar[l] \ar[r]^{x_i} &
\mathcal{X}
}
$$
dont les carrés sont cartésiens. Puisque $\{W_i \to W'\}_{i \in I}$
est le changement de base de $\mathcal{U}$ par $W' \to U$, nous concluons
qu'il s'agit d'un recouvrement fpqc. Puisque $W$ satisfait à la condition de faisceau
pour les recouvrements fpqc, nous obtenons un unique morphisme
$W' \to W$ tel que $W_i \to W' \to W$ soit le morphisme donné
$W_i \to W$. Autrement dit, nous avons les diagrammes commutatifs
$$
\xymatrix{
W_i \ar[d] \ar[r] &
W' \ar[d] \ar[r] &
W \ar[d] \\
U_i \ar[r] \ar@/_1pc/[rr]_{x_i} &
U &
\mathcal{X}
}
$$
compatibles avec les isomorphismes $\varphi_{ij}$ et dont le
carré et le rectangle sont cartésiens.

\medskip\noindent
Choisissons une famille de schémas affines $W_\alpha$, $\alpha \in A$,
et des morphismes lisses $W_\alpha \to \mathcal{X}$ tels que
$\coprod W_\alpha \to \mathcal{X}$ soit surjectif.
Par le procédé du paragraphe précédent, nous construisons
un diagramme
$$
\xymatrix{
W_{\alpha, i} \ar[d] \ar[r] &
W_\alpha' \ar[d] \ar[r] &
W_\alpha \ar[d] \\
U_i \ar[r] \ar@/_1pc/[rr]_{x_i} &
U &
\mathcal{X}
}
$$
pour chaque $\alpha$. Les morphismes $W_\alpha' \to U$ sont alors
lisses et conjointement surjectifs.

\medskip\noindent
Notons $x_\alpha$ l'objet de $\mathcal{X}$ au-dessus de $W_\alpha'$
correspondant à $W_\alpha' \to W_\alpha \to \mathcal{X}$.
Puisque $\mathcal{X}$ est un champ pour la topologie fppf et que
$\{W_\alpha' \to U\}$ est un recouvrement fppf, il suffit
de montrer qu'il existe des isomorphismes
$\text{pr}_0^*x_\alpha \to \text{pr}_1^*x_\beta$
sur $W_\alpha' \times_U W'_\beta$ satisfaisant à la
condition de cocycle. Or, après restriction à
$W_{\alpha, i}$, nous disposons bien de tels isomorphismes sur
$W_{\alpha, i} \times_{U_i} W_{\beta, i} =
U_i \times_U (W_\alpha' \times_U W'_\beta)$
car l'image inverse de $x_\alpha$ sur $W_{\alpha, i}$
est isomorphe à l'image inverse de $x_i$ sur $W_{\alpha, i}$.
Puisque $\{U_i \times_U (W_\alpha' \times_U W'_\beta) \to
W_\alpha' \times_U W'_\beta\}_{i \in I}$
est un recouvrement fpqc et que, par la compatibilité susmentionnée
des diagrammes ci-dessus avec $\varphi_{ij}$, ces
isomorphismes descendent à $W_\alpha' \times_U W'_\beta$,
la démonstration est achevée.
\end{proof}










\section{Foncteurs tensoriels}
\label{section-tensor-functors}

\noindent
Soit $f : \mathcal{Y} \to \mathcal{X}$ un morphisme de champs algébriques
noethériens. Le foncteur image inverse
$$
f^* :
\textit{Coh}(\mathcal{O}_\mathcal{X})
\longrightarrow
\textit{Coh}(\mathcal{O}_\mathcal{Y})
$$
est un foncteur tensoriel exact à droite : il est additif, exact à droite
et commute aux produits tensoriels de modules cohérents. On peut
se demander dans quelle mesure tout foncteur tensoriel exact à droite
$F : \textit{Coh}(\mathcal{O}_\mathcal{X}) \to
\textit{Coh}(\mathcal{O}_\mathcal{Y})$ provient d'un
morphisme $f : \mathcal{Y} \to \mathcal{X}$.
Le lecteur pourra consulter
\cite{Hall-Rydh-coherent} pour un résultat très général
de ce type.
Le but de cette section est de donner une courte démonstration du
théorème \ref{theorem-main} en guise d'introduction à ces idées.

\medskip\noindent
Commençons par quelques lemmes.

\begin{lemma}
\label{lemma-extend}
Soient $\mathcal{X}$ et $\mathcal{Y}$ des champs algébriques noethériens.
Tout foncteur tensoriel exact à droite $F : \textit{Coh}(\mathcal{O}_\mathcal{X}) \to
\textit{Coh}(\mathcal{O}_\mathcal{Y})$ se prolonge de manière unique en un
foncteur tensoriel exact à droite
$F : \QCoh(\mathcal{O}_\mathcal{X}) \to \QCoh(\mathcal{O}_\mathcal{Y})$
commutant à toutes les limites inductives.
\end{lemma}

\begin{proof}
L'existence et l'unicité du prolongement sont un fait général ; voir
Catégories, lemme \ref{categories-lemma-extend-functor-by-colim}.
Pour voir que le lemme s'applique, remarquons que les modules cohérents sur des
champs algébriques localement noethériens sont, par définition, des modules
de présentation finie ; voir
Cohomologie des champs, définition \ref{stacks-cohomology-definition-coherent}.
Ainsi, un module cohérent sur $\mathcal{X}$ est un
objet compact au sens catégorique de $\QCoh(\mathcal{O}_\mathcal{X})$ d'après
Cohomologie des champs, lemme
\ref{stacks-cohomology-lemma-finite-presentation-quasi-compact-colimit}.
Enfin, tout module quasi-cohérent est une limite inductive filtrante de
ses sous-modules cohérents d'après Cohomologie des champs, lemme
\ref{stacks-cohomology-lemma-directed-colimit-coherent}.

\medskip\noindent
Puisque $F$ est additif, le prolongement de $F$ l'est aussi (détails omis).
Puisque $F$ est un foncteur tensoriel et que les limites inductives de modules
commutent à la formation des produits tensoriels, le prolongement de $F$
est lui aussi un foncteur tensoriel (détails omis).

\medskip\noindent
Dans ce paragraphe, montrons que le prolongement commute aux sommes directes
arbitraires. Si $\mathcal{F} = \bigoplus_{j \in J} \mathcal{H}_j$
avec $\mathcal{H}_j$ quasi-cohérent, alors
$\mathcal{F} = \colim_{J' \subset J\text{ fini}}
\bigoplus_{j \in J'} \mathcal{H}_j$.
En notant encore $F$ le prolongement de $F$, nous obtenons
\begin{align*}
F(\mathcal{F})
& =
\colim_{J' \subset J\text{ fini}}
F(\bigoplus\nolimits_{j \in J'} \mathcal{H}_j) \\
& =
\colim_{J' \subset J\text{ fini}}
\bigoplus\nolimits_{j \in J'} F(\mathcal{H}_j) \\
& =
\bigoplus\nolimits_{j \in J} F(\mathcal{H}_j)
\end{align*}
Ainsi, $F$ commute aux sommes directes arbitraires.

\medskip\noindent
Dans ce paragraphe, montrons que le prolongement est exact à droite.
Supposons que $0 \to \mathcal{F} \to \mathcal{F}' \to \mathcal{F}'' \to 0$
soit une suite exacte courte de $\mathcal{O}_\mathcal{X}$-modules quasi-cohérents.
Écrivons alors $\mathcal{F}' = \bigcup_{i \in I} \mathcal{F}'_i$ comme
réunion de ses sous-modules cohérents (voir la référence donnée plus haut).
Notons $\mathcal{F}''_i \subset \mathcal{F}''$ l'image de $\mathcal{F}'_i$
et posons $\mathcal{F}_i = \mathcal{F} \cap \mathcal{F}'_i =
\Ker(\mathcal{F}'_i \to \mathcal{F}''_i)$. Il est alors clair que
$\mathcal{F} = \bigcup \mathcal{F}_i$ et
$\mathcal{F}'' = \bigcup \mathcal{F}''_i$
et que nous avons des suites exactes courtes
$$
0 \to \mathcal{F}_i \to \mathcal{F}_i' \to \mathcal{F}_i'' \to 0
$$
Puisque le prolongement commute aux limites inductives filtrantes, nous avons
$F(\mathcal{F}) = \colim_{i \in I} F(\mathcal{F}_i)$,
$F(\mathcal{F}') = \colim_{i \in I} F(\mathcal{F}'_i)$, et
$F(\mathcal{F}'') = \colim_{i \in I} F(\mathcal{F}''_i)$.
Comme les limites inductives filtrantes de faisceaux de modules sont exactes, nous
concluons que le prolongement de $F$ est exact à droite.

\medskip\noindent
La démonstration est achevée, car tout foncteur exact à droite qui commute aux
coproduits commute à toutes les limites inductives ; voir
Catégories, lemme \ref{categories-lemma-colimits-coproducts-coequalizers}.
\end{proof}

\begin{lemma}
\label{lemma-affine}
Soit $\mathcal{X}$ un champ algébrique à diagonale affine.
Soit $B$ un anneau. Soit $F : \QCoh(\mathcal{O}_\mathcal{X}) \to \text{Mod}_B$
un foncteur tensoriel exact à droite qui commute aux sommes directes.
Soit $g : U \to \mathcal{X}$ un morphisme, où $U = \Spec(A)$ est affine. Alors
\begin{enumerate}
\item $C = F(g_{\QCoh, *}\mathcal{O}_U)$ est une $B$-algèbre commutative, et
\item il existe un homomorphisme d'anneaux $A \to C$
\end{enumerate}
tel que $F \circ g_{\QCoh, *} : \text{Mod}_A \to \text{Mod}_B$
envoie $M$ sur $M \otimes_A C$ considéré comme $B$-module.
\end{lemma}

\begin{proof}
Remarquons que $g$ est quasi-compact et quasi-séparé ; voir
Morphismes de champs, lemme
\ref{stacks-morphisms-lemma-quasi-compact-quasi-separated-permanence}.
Dans Cohomologie des champs, proposition
\ref{stacks-cohomology-proposition-direct-image-quasi-coherent}
nous avons construit le foncteur
$g_{\QCoh, *} : \QCoh(\mathcal{O}_U) \to \QCoh(\mathcal{O}_\mathcal{X})$.
D'après Cohomologie des champs, remarques
\ref{stacks-cohomology-remark-direct-image-quasi-coherent-tensor} et
\ref{stacks-cohomology-remark-QCoh-tensor}
nous obtenons une multiplication
$$
\mu :
g_{\QCoh, *}\mathcal{O}_U
\otimes_{\mathcal{O}_\mathcal{X}}
g_{\QCoh, *}\mathcal{O}_U
\longrightarrow
g_{\QCoh, *}\mathcal{O}_U
$$
qui munit $g_{\QCoh, *}\mathcal{O}_U$ d'une structure de
$\mathcal{O}_\mathcal{X}$-algèbre commutative. Ainsi, $C = F(g_{\QCoh, *}\mathcal{O}_U)$
est une algèbre commutative dans $\text{Mod}_B$ ; autrement dit, $C$ est une
$B$-algèbre commutative. Remarquons que nous avons un homomorphisme d'anneaux
$\kappa : A \to \text{End}_{\mathcal{O}_\mathcal{X}}(g_{\QCoh, *}\mathcal{O}_U)$
tel que, pour tout $a \in A$, le diagramme
$$
\xymatrix{
g_{\QCoh, *}\mathcal{O}_U \otimes_{\mathcal{O}_\mathcal{X}}
g_{\QCoh, *}\mathcal{O}_U
\ar[d]_{\kappa(a) \otimes 1} \ar[rr]_-\mu & &
g_{\QCoh, *}\mathcal{O}_U \ar[d]^{\kappa(a)} \\
g_{\QCoh, *}\mathcal{O}_U \otimes_{\mathcal{O}_\mathcal{X}}
g_{\QCoh, *}\mathcal{O}_U
\ar[rr]^-\mu & &
g_{\QCoh, *}\mathcal{O}_U
}
$$
est commutatif. Il s'ensuit que l'on obtient un homomorphisme d'anneaux
$\kappa' = F(\kappa) : A \to \text{End}_B(C)$
tel que $\kappa'(a)(c) c' = \kappa'(a)(cc')$ ; bien entendu, cela signifie que
$a \mapsto \kappa'(a)(1)$ est un homomorphisme d'anneaux $A \to C$.

\medskip\noindent
Le morphisme $g : U \to \mathcal{X}$ est affine ; voir
Morphismes de champs, lemme \ref{stacks-morphisms-lemma-affine-permanence}.
Ainsi, $g_{\QCoh, *}$ est exact et commute aux sommes directes d'après
Cohomologie des champs, lemme
\ref{stacks-cohomology-lemma-quasi-coherent-pushforward-affine}.
Ainsi, $F \circ g_{\QCoh, *} : \text{Mod}_A \to \text{Mod}_B$
est un foncteur exact à droite qui commute aux sommes directes
et envoie $A$ sur $C$.
D'après Foncteurs et morphismes, lemme \ref{functors-lemma-functor},
nous voyons que le foncteur $F \circ g_{\QCoh, *}$ envoie un $A$-module
$M$ sur $M \otimes_A C$ considéré comme un $B$-module.
\end{proof}

\begin{lemma}
\label{lemma-universally-injective}
Avec les notations du lemme \ref{lemma-affine}, supposons que $\mathcal{X}$ soit
noethérien et que $g$ soit surjectif et plat.
Alors $B \to C$ est universellement injectif.
\end{lemma}

\begin{proof}
Considérons le morphisme naturel
$1 : \mathcal{O}_\mathcal{X} \to g_{\QCoh, *}\mathcal{O}_U$
dans $\QCoh(\mathcal{O}_\mathcal{X})$. En prenant son image inverse sur $U$ et
en utilisant l'adjonction, nous constatons que le composé
$$
\mathcal{O}_U =
g^*\mathcal{O}_\mathcal{X} \xrightarrow{g^*1}
g^*g_{\QCoh, *}\mathcal{O}_U \to \mathcal{O}_U
$$
est l'identité dans $\QCoh(\mathcal{O}_U)$.
Écrivons $g_{\QCoh, *}\mathcal{O}_U = \colim \mathcal{F}_i$
sous la forme d'une limite inductive filtrante de $\mathcal{O}_\mathcal{X}$-modules cohérents ; voir
Cohomologie des champs, lemme
\ref{stacks-cohomology-lemma-directed-colimit-coherent}.
Pour $i$ assez grand, le morphisme
$1 : \mathcal{O}_\mathcal{X} \to g_{\QCoh, *}\mathcal{O}_U$
se factorise par $\mathcal{F}_i$ ; voir
Cohomologie des champs, lemme
\ref{stacks-cohomology-lemma-finite-presentation-quasi-compact-colimit}.
Soit $s : \mathcal{O}_\mathcal{X} \to \mathcal{F}_i$ cette
factorisation. Alors
$$
\mathcal{O}_U \xrightarrow{g^*s} g^*\mathcal{F}_i \to
g^*g_{\QCoh, *}\mathcal{O}_U \to \mathcal{O}_U
$$
est l'identité. Autrement dit, nous voyons que $s$ devient
l'inclusion d'un facteur direct après image inverse sur $U$. Posons
$\mathcal{F}_i^\vee = hom(\mathcal{F}_i, \mathcal{O}_\mathcal{X})$
avec les notations de
Cohomologie des champs, lemme
\ref{stacks-cohomology-lemma-internal-hom-fp-into-qcoh}.
En particulier, il existe un morphisme d'évaluation
$ev : \mathcal{F}_i \otimes_{\mathcal{O}_\mathcal{X}} \mathcal{F}_i^\vee
\to \mathcal{O}_\mathcal{X}$.
L'évaluation en $s$ définit un morphisme
$s^\vee : \mathcal{F}_i^\vee \to \mathcal{O}_\mathcal{X}$.
Par dualité avec l'assertion concernant $s$, nous voyons que $g^*(s^\vee)$ est surjectif ;
voir Cohomologie des champs, section
\ref{stacks-cohomology-section-further-remarks}
pour la compatibilité de $hom$ et de $\otimes$ avec la restriction à $U$.
Puisque $g$ est surjectif et plat, nous concluons que $s^\vee$ est surjectif
(voir la référence citée). Puisque $F$ est exact à droite, nous concluons que
$F(\mathcal{F}_i^\vee) \to F(\mathcal{O}_\mathcal{X}) = B$ est surjectif.
Choisissons $\lambda \in F(\mathcal{F}_i^\vee)$
qui s'envoie sur $1 \in B$. Notons $e = F(s)(1) \in F(\mathcal{F}_i)$
l'image de $1$ par le morphisme
$F(s) : B = F(\mathcal{O}_\mathcal{X}) \to F(\mathcal{F}_i)$.
Alors le morphisme
$$
F(ev) :
F(\mathcal{F}_i) \otimes_B F(\mathcal{F}_i^\vee) =
F(\mathcal{F}_i \otimes_{\mathcal{O}_\mathcal{X}} \mathcal{F}_i^\vee)
\longrightarrow
F(\mathcal{O}_\mathcal{X}) = B
$$
envoie $e \otimes \lambda$ sur $1$ par construction. Par conséquent, le
morphisme $B \to F(\mathcal{F}_i)$, $b \mapsto be$,
est universellement injectif, car il admet l'inverse à gauche
$F(\mathcal{F}_i) \to B$, $\xi \mapsto F(ev)(\xi \otimes \lambda)$.
Comme cela vaut pour tout $i$ assez grand, nous concluons.
\end{proof}

\begin{lemma}
\label{lemma-flat}
Soit $B \to C$ un homomorphisme d'anneaux. Si
\begin{enumerate}
\item les coprojections $C \to C \otimes_B C$ sont plates, et
\item $B \to C$ est universellement injectif,
\end{enumerate}
alors $B \to C$ est fidèlement plat.
\end{lemma}

\begin{proof}
Le morphisme $\Spec(C) \to \Spec(B)$ est surjectif, puisque $B \to C$ est universellement
injectif. Il suffit donc de montrer que $B \to C$ est plat,
ce qui résulte de
Descente, théorème \ref{descent-theorem-descend-module-properties}.
\end{proof}

\noindent
La version très simple de la formule de K\"unneth qui suit devrait devenir
obsolète lorsque nous aurons rédigé une section sur les théorèmes de K\"unneth pour la
cohomologie des modules quasi-cohérents sur les champs algébriques.

\begin{lemma}
\label{lemma-easy-kunneth}
Soient $a : \mathcal{Y} \to \mathcal{X}$ et $b : \mathcal{Z} \to \mathcal{X}$
des morphismes représentables par des schémas, quasi-compacts, quasi-séparés et plats.
Alors $a_{\QCoh, *}\mathcal{O}_\mathcal{Y}
\otimes_{\mathcal{O}_\mathcal{X}}
b_{\QCoh, *}\mathcal{O}_\mathcal{Z} =
f_{\QCoh, *}\mathcal{O}_{\mathcal{Y} \times_\mathcal{X} \mathcal{Z}}$
où $f : \mathcal{Y} \times_\mathcal{X} \mathcal{Z} \to \mathcal{X}$
est le morphisme évident.
\end{lemma}

\begin{proof}
Pour abréger, posons $\mathcal{P} = \mathcal{Y} \times_\mathcal{X} \mathcal{Z}$.
Puisque $a \circ \text{pr}_1 = f$ et $b \circ \text{pr}_2 = f$,
nous obtenons des morphismes $a_*\mathcal{O}_\mathcal{Y} \to f_*\mathcal{O}_\mathcal{P}$
et $b_*\mathcal{O}_\mathcal{Z} \to f_*\mathcal{O}_\mathcal{P}$ (en utilisant
les morphismes d'image inverse relatifs ; voir Sites, section \ref{sites-section-pullback}).
Nous obtenons donc un cup-produit relatif
$$
\mu :
a_*\mathcal{O}_\mathcal{Y}
\otimes_{\mathcal{O}_\mathcal{X}}
b_*\mathcal{O}_\mathcal{Z} \longrightarrow
f_*\mathcal{O}_{\mathcal{Y} \times_\mathcal{X} \mathcal{Z}}
$$
En appliquant $Q$ et sa compatibilité avec les produits tensoriels
(Cohomologie des champs, remarque \ref{stacks-cohomology-remark-QCoh-tensor}),
nous obtenons une flèche
$Q(\mu) : a_{\QCoh, *}\mathcal{O}_\mathcal{Y}
\otimes_{\mathcal{O}_\mathcal{X}}
b_{\QCoh, *}\mathcal{O}_\mathcal{Z} \to
f_{\QCoh, *}\mathcal{O}_{\mathcal{Y} \times_\mathcal{X} \mathcal{Z}}$
dans $\QCoh(\mathcal{O}_\mathcal{X})$.
Choisissons ensuite un schéma $U$ et un morphisme lisse surjectif
$U \to \mathcal{X}$. Il suffit de montrer que la restriction de
$Q(\mu)$ à $U_\etale$ est un isomorphisme ; voir
Cohomologie des champs, section \ref{stacks-cohomology-section-further-remarks}.
À son tour, d'après les résultats de la même section,
il suffit de montrer que la restriction de $\mu$ à $U_\etale$
est un isomorphisme (on utilise ici que la source et le but de $\mu$
sont des modules localement quasi-cohérents vérifiant la propriété de changement de base).
En outre, nous pouvons calculer les images directes pour la topologie \'etale ; voir
Cohomologie des champs, proposition
\ref{stacks-cohomology-proposition-loc-qcoh-flat-base-change}.
Enfin, nous voyons que
$a_*\mathcal{O}_\mathcal{Y}|_{U_\etale} = (V \to U)_{small, *}\mathcal{O}_V$
où $V = U \times_\mathcal{X} \mathcal{Y}$.
De même pour $b_*$ et $f_*$. Ainsi, le résultat découle
de la formule de K\"unneth pour les morphismes plats, quasi-compacts et quasi-séparés
de schémas ; voir Catégories dérivées de schémas,
lemme \ref{perfect-lemma-kunneth}.
\end{proof}

\begin{lemma}
\label{lemma-fully-faithful}
Soit $\mathcal{X}$ un champ algébrique à diagonale affine.
Soit $B$ un anneau. Soient $f_i : \Spec(B) \to \mathcal{X}$, $i = 1, 2$,
deux morphismes. Soit $t : f_1^* \to f_2^*$ un isomorphisme
des foncteurs tensoriels
$f_i^* : \QCoh(\mathcal{O}_\mathcal{X}) \to \text{Mod}_B$.
Alors il existe une $2$-flèche $f_1 \to f_2$ qui induit $t$.
\end{lemma}

\begin{proof}
Choisissons un schéma affine $U = \Spec(A)$ et un morphisme lisse surjectif
$g : U \to \mathcal{X}$ ; voir
Propriétés des champs, lemme \ref{stacks-properties-lemma-quasi-compact-stack}.
Puisque la diagonale de $\mathcal{X}$ est affine, nous voyons que
$U_i = \Spec(B) \times_{f_i, \mathcal{X}, g} U$ est affine.
Écrivons $U_i = \Spec(C_i)$. Alors $C_i$ est la $B$-algèbre munie
de l'homomorphisme d'anneaux $A \to C_i$ construit au lemme \ref{lemma-affine}
au moyen du foncteur $F = f_i^*$. Par conséquent, $t$ induit un isomorphisme
$C_1 \to C_2$ de $B$-algèbres, compatible avec les homomorphismes d'anneaux
$A \to C_1$ et $A \to C_2$. Autrement dit, nous avons les diagrammes
commutatifs
$$
\xymatrix{
U_i \ar[r] \ar[d] & U \ar[d]^g \\
\Spec(B) \ar[r]^{f_i} & \mathcal{X}
}
\quad\quad
\xymatrix{
&
U_2 \ar[ld] \ar[d]^{\cong} \ar[rd] \\
\Spec(B) &
U_1 \ar[l] \ar[r] &
U
}
$$
Cela montre déjà que les objets $f_1$ et $f_2$ de $\mathcal{X}$
au-dessus de $\Spec(B)$ deviennent isomorphes après le recouvrement lisse
$\{U_1 \to \Spec(B)\}$. Pour montrer que cet isomorphisme descend en un isomorphisme
de $f_1$ et $f_2$ sur $\Spec(B)$, nous devons montrer que notre isomorphisme
(qui provient des diagrammes commutatifs ci-dessus) est compatible
avec les données de descente sur $U_1 \times_{\Spec(B)} U_1$. À cette fin,
remarquons que $U \times_\mathcal{X} U$ est lui aussi affine, que
nous disposons du morphisme $g' : U \times_\mathcal{X} U \to \mathcal{X}$, et
que
$$
U_i \times_{\Spec(B)} U_i =
\Spec(B) \times_{f_i, \mathcal{X}, g'}
(U \times_\mathcal{X} U)
$$
Il s'ensuit que l'isomorphisme $C_1 \otimes_B C_1 \to C_2 \otimes_B C_2$
induit par l'isomorphisme $C_1 \to C_2$ est compatible avec les
morphismes $U_i \times_{\Spec(B)} U_i \to U \times_\mathcal{X} U$.
Certains détails sont omis.
\end{proof}

\begin{lemma}
\label{lemma-main}
Let $\mathcal{X}$ be a Noetherian algebraic stack with affine diagonal.
Let $B$ be a ring.
Let $F : \QCoh(\mathcal{O}_\mathcal{X}) \to \text{Mod}_B$
be a right exact tensor functor which commutes with direct sums.
Then $F$ comes from a unique morphism $\Spec(B) \to \mathcal{X}$.
\end{lemma}

\begin{proof}
Choose a surjective smooth morphism $g : U \to \mathcal{X}$ with $U = \Spec(A)$
affine, see Properties of Stacks, Lemma
\ref{stacks-properties-lemma-quasi-compact-stack}. Apply
Lemma \ref{lemma-affine} to get the finite type
commutative $B$-algebra $C = F(g_{\QCoh, *}\mathcal{O}_U)$
and the ring map $A \to C$.
By Lemma \ref{lemma-universally-injective}
the ring map $B \to C$ is universally injective.
Consider the algebra
$$
C \otimes_B C =
F(g_{\QCoh, *}\mathcal{O}_U
\otimes_{\mathcal{O}_\mathcal{X}}
g_{\QCoh, *}\mathcal{O}_U)
$$
Since $g$ is flat, quasi-compact, and quasi-separated
by Lemma \ref{lemma-easy-kunneth} we have the first equality in
$$
g_{\QCoh, *}\mathcal{O}_U
\otimes_{\mathcal{O}_\mathcal{X}}
g_{\QCoh, *}\mathcal{O}_U
=
f_{\QCoh, *}\mathcal{O}_{U \times_\mathcal{X} U} =
g_{\QCoh, *}(\text{pr}_{2, *}\mathcal{O}_{U \times_\mathcal{X} U})
$$
where
$f : U \times_\mathcal{X} U \to \mathcal{X}$ is the obvious morphism
and $\text{pr}_2 : U \times_\mathcal{X} U \to U$ is the second
projection. The second equality follows from
Cohomology of Stacks, Lemma \ref{stacks-cohomology-lemma-relative-leray}
and $f = g \circ \text{pr}_2$. Since the diagonal of $\mathcal{X}$
is affine, we see that $U \times_\mathcal{X} U = \Spec(R)$ is affine.
Let us use $\text{pr}_2 : A \to R$ to view $R$ as an $A$-algebra.
All in all we obtain
$$
C \otimes_B C =
F(g_{\QCoh, *}\mathcal{O}_U
\otimes_{\mathcal{O}_\mathcal{X}}
g_{\QCoh, *}\mathcal{O}_U) =
F(g_{\QCoh, *}(\text{pr}_{2, *}\mathcal{O}_{U \times_\mathcal{X} U})) =
R \otimes_A C
$$
where the final equality follows from the final statement of
Lemma \ref{lemma-affine}. Since $A \to R$ is flat (because
$\text{pr}_2$ is flat as a base change of $U \to \mathcal{X}$),
we conclude that $C \otimes_B C$ is flat over $C$.
By Lemma \ref{lemma-flat} we conclude that $B \to C$ is faithfully flat.

\medskip\noindent
We claim there is a solid commutative diagram
$$
\xymatrix{
\Spec(C \otimes_B C) \ar@<1ex>[d] \ar@<-1ex>[d] \ar[r] &
U \times_\mathcal{X} U \ar@<1ex>[d] \ar@<-1ex>[d] \\
\Spec(C) \ar[d] \ar[r] &
U \ar[d] \\
\Spec(B) \ar@{..>}[r] &
\mathcal{X}
}
$$
The arrow $\Spec(C) \to U = \Spec(A)$ comes from the ring map
$A \to C$ in the statement of Lemma \ref{lemma-affine}.
The arrow $\Spec(C \otimes_B C) \to U \times_\mathcal{X} U$
similarly comes from the ring map $R \to C \otimes_B C$.
To verify the top square commutes use Lemma \ref{lemma-fully-faithful};
details omitted.
We conclude we get the dotted arrow $\Spec(B) \to \mathcal{X}$
by Proposition \ref{proposition-stack-fpqc}.

\medskip\noindent
The statement that $F$ is the functor corresponding to pullback
by the dotted arrow is also clear from this and the corresponding
statement in Lemma \ref{lemma-affine}. Details omitted.
\end{proof}

\noindent
For a ring $B$ let us denote $\text{Mod}^{fg}_B$ the category of
finitely generated $B$-modules (AKA finite $B$-modules).

\begin{theorem}
\label{theorem-main}
Let $\mathcal{X}$ be a Noetherian algebraic stack with affine diagonal.
Let $B$ be a Noetherian ring.
Let $F : \text{Coh}(\mathcal{O}_\mathcal{X}) \to \text{Mod}^{fg}_B$
be a right exact tensor functor.
Then $F$ comes from a unique morphism $\Spec(B) \to \mathcal{X}$.
\end{theorem}

\begin{proof}
By Lemma \ref{lemma-extend}
we can extend $F$ uniquely to a right exact tensor functor
$F : \QCoh(\mathcal{O}_\mathcal{X}) \to \text{Mod}_B$
commuting with all direct susms. Then we can apply
Lemma \ref{lemma-main}.
\end{proof}










\begin{multicols}{2}[\section{Autres chapitres}]
\noindent
Préliminaires
\begin{enumerate}
\item \hyperref[introduction-section-phantom]{Introduction}
\item \hyperref[conventions-section-phantom]{Conventions}
\item \hyperref[sets-section-phantom]{Théorie des ensembles}
\item \hyperref[categories-section-phantom]{Catégories}
\item \hyperref[topology-section-phantom]{Topologie}
\item \hyperref[sheaves-section-phantom]{Faisceaux sur les espaces}
\item \hyperref[sites-section-phantom]{Sites et faisceaux}
\item \hyperref[stacks-section-phantom]{Champs}
\item \hyperref[fields-section-phantom]{Corps}
\item \hyperref[algebra-section-phantom]{Algèbre commutative}
\item \hyperref[brauer-section-phantom]{Groupes de Brauer}
\item \hyperref[homology-section-phantom]{Algèbre homologique}
\item \hyperref[derived-section-phantom]{Catégories dérivées}
\item \hyperref[simplicial-section-phantom]{Méthodes simpliciales}
\item \hyperref[more-algebra-section-phantom]{Compléments d'algèbre}
\item \hyperref[smoothing-section-phantom]{Lissage des morphismes d'anneaux}
\item \hyperref[modules-section-phantom]{Faisceaux de modules}
\item \hyperref[sites-modules-section-phantom]{Modules sur les sites}
\item \hyperref[injectives-section-phantom]{Objets injectifs}
\item \hyperref[cohomology-section-phantom]{Cohomologie des faisceaux}
\item \hyperref[sites-cohomology-section-phantom]{Cohomologie sur les sites}
\item \hyperref[dga-section-phantom]{Algèbre différentielle graduée}
\item \hyperref[dpa-section-phantom]{Algèbre à puissances divisées}
\item \hyperref[sdga-section-phantom]{Faisceaux différentiels gradués}
\item \hyperref[hypercovering-section-phantom]{Hyperrecouvrements}
\end{enumerate}
Schémas
\begin{enumerate}
\setcounter{enumi}{25}
\item \hyperref[schemes-section-phantom]{Schémas}
\item \hyperref[constructions-section-phantom]{Constructions de schémas}
\item \hyperref[properties-section-phantom]{Propriétés des schémas}
\item \hyperref[morphisms-section-phantom]{Morphismes de schémas}
\item \hyperref[coherent-section-phantom]{Cohomologie des schémas}
\item \hyperref[divisors-section-phantom]{Diviseurs}
\item \hyperref[limits-section-phantom]{Limites de schémas}
\item \hyperref[varieties-section-phantom]{Variétés}
\item \hyperref[topologies-section-phantom]{Topologies sur les schémas}
\item \hyperref[descent-section-phantom]{Descente}
\item \hyperref[perfect-section-phantom]{Catégories dérivées de schémas}
\item \hyperref[more-morphisms-section-phantom]{Compléments sur les morphismes}
\item \hyperref[flat-section-phantom]{Compléments sur la platitude}
\item \hyperref[groupoids-section-phantom]{Schémas en groupoïdes}
\item \hyperref[more-groupoids-section-phantom]{Compléments sur les schémas en groupoïdes}
\item \hyperref[etale-section-phantom]{Morphismes étales de schémas}
\end{enumerate}
Thèmes de théorie des schémas
\begin{enumerate}
\setcounter{enumi}{41}
\item \hyperref[chow-section-phantom]{Homologie de Chow}
\item \hyperref[intersection-section-phantom]{Théorie de l'intersection}
\item \hyperref[pic-section-phantom]{Schémas de Picard des courbes}
\item \hyperref[weil-section-phantom]{Théories cohomologiques de Weil}
\item \hyperref[adequate-section-phantom]{Modules adéquats}
\item \hyperref[dualizing-section-phantom]{Complexes dualisants}
\item \hyperref[duality-section-phantom]{Dualité pour les schémas}
\item \hyperref[discriminant-section-phantom]{Discriminants et différentes}
\item \hyperref[derham-section-phantom]{Cohomologie de de Rham}
\item \hyperref[local-cohomology-section-phantom]{Cohomologie locale}
\item \hyperref[algebraization-section-phantom]{Géométrie algébrique et formelle}
\item \hyperref[curves-section-phantom]{Courbes algébriques}
\item \hyperref[resolve-section-phantom]{Résolution des surfaces}
\item \hyperref[models-section-phantom]{Réduction semi-stable}
\item \hyperref[functors-section-phantom]{Foncteurs et morphismes}
\item \hyperref[equiv-section-phantom]{Catégories dérivées de variétés}
\item \hyperref[pione-section-phantom]{Groupes fondamentaux de schémas}
\item \hyperref[etale-cohomology-section-phantom]{Cohomologie étale}
\item \hyperref[crystalline-section-phantom]{Cohomologie cristalline}
\item \hyperref[proetale-section-phantom]{Cohomologie pro-étale}
\item \hyperref[relative-cycles-section-phantom]{Cycles relatifs}
\item \hyperref[more-etale-section-phantom]{Compléments de cohomologie étale}
\item \hyperref[trace-section-phantom]{La formule des traces}
\end{enumerate}
Espaces algébriques
\begin{enumerate}
\setcounter{enumi}{64}
\item \hyperref[spaces-section-phantom]{Espaces algébriques}
\item \hyperref[spaces-properties-section-phantom]{Propriétés des espaces algébriques}
\item \hyperref[spaces-morphisms-section-phantom]{Morphismes d'espaces algébriques}
\item \hyperref[decent-spaces-section-phantom]{Espaces algébriques décents}
\item \hyperref[spaces-cohomology-section-phantom]{Cohomologie des espaces algébriques}
\item \hyperref[spaces-limits-section-phantom]{Limites d'espaces algébriques}
\item \hyperref[spaces-divisors-section-phantom]{Diviseurs sur les espaces algébriques}
\item \hyperref[spaces-over-fields-section-phantom]{Espaces algébriques sur des corps}
\item \hyperref[spaces-topologies-section-phantom]{Topologies sur les espaces algébriques}
\item \hyperref[spaces-descent-section-phantom]{Descente et espaces algébriques}
\item \hyperref[spaces-perfect-section-phantom]{Catégories dérivées d'espaces}
\item \hyperref[spaces-more-morphisms-section-phantom]{Compléments sur les morphismes d'espaces}
\item \hyperref[spaces-flat-section-phantom]{Platitude sur les espaces algébriques}
\item \hyperref[spaces-groupoids-section-phantom]{Groupoïdes dans les espaces algébriques}
\item \hyperref[spaces-more-groupoids-section-phantom]{Compléments sur les groupoïdes dans les espaces}
\item \hyperref[bootstrap-section-phantom]{Amorçage}
\item \hyperref[spaces-pushouts-section-phantom]{Sommes amalgamées d'espaces algébriques}
\end{enumerate}
Thèmes de géométrie
\begin{enumerate}
\setcounter{enumi}{81}
\item \hyperref[spaces-chow-section-phantom]{Groupes de Chow des espaces}
\item \hyperref[groupoids-quotients-section-phantom]{Quotients de groupoïdes}
\item \hyperref[spaces-more-cohomology-section-phantom]{Compléments sur la cohomologie des espaces}
\item \hyperref[spaces-simplicial-section-phantom]{Espaces simpliciaux}
\item \hyperref[spaces-duality-section-phantom]{Dualité pour les espaces}
\item \hyperref[formal-spaces-section-phantom]{Espaces algébriques formels}
\item \hyperref[restricted-section-phantom]{Algébrisation des espaces formels}
\item \hyperref[spaces-resolve-section-phantom]{Résolution des surfaces revisitée}
\end{enumerate}
Théorie des déformations
\begin{enumerate}
\setcounter{enumi}{89}
\item \hyperref[formal-defos-section-phantom]{Théorie formelle des déformations}
\item \hyperref[defos-section-phantom]{Théorie des déformations}
\item \hyperref[cotangent-section-phantom]{Le complexe cotangent}
\item \hyperref[examples-defos-section-phantom]{Problèmes de déformation}
\end{enumerate}
Champs algébriques
\begin{enumerate}
\setcounter{enumi}{93}
\item \hyperref[algebraic-section-phantom]{Champs algébriques}
\item \hyperref[examples-stacks-section-phantom]{Exemples de champs}
\item \hyperref[stacks-sheaves-section-phantom]{Faisceaux sur les champs algébriques}
\item \hyperref[criteria-section-phantom]{Critères de représentabilité}
\item \hyperref[artin-section-phantom]{Axiomes d'Artin}
\item \hyperref[quot-section-phantom]{Espaces de Quot et de Hilbert}
\item \hyperref[stacks-properties-section-phantom]{Propriétés des champs algébriques}
\item \hyperref[stacks-morphisms-section-phantom]{Morphismes de champs algébriques}
\item \hyperref[stacks-limits-section-phantom]{Limites de champs algébriques}
\item \hyperref[stacks-cohomology-section-phantom]{Cohomologie des champs algébriques}
\item \hyperref[stacks-perfect-section-phantom]{Catégories dérivées de champs}
\item \hyperref[stacks-introduction-section-phantom]{Introduction aux champs algébriques}
\item \hyperref[stacks-more-morphisms-section-phantom]{Compléments sur les morphismes de champs}
\item \hyperref[stacks-geometry-section-phantom]{La géométrie des champs}
\end{enumerate}
Thèmes de théorie des espaces de modules
\begin{enumerate}
\setcounter{enumi}{107}
\item \hyperref[moduli-section-phantom]{Champs de modules}
\item \hyperref[moduli-curves-section-phantom]{Espaces de modules de courbes}
\end{enumerate}
Divers
\begin{enumerate}
\setcounter{enumi}{109}
\item \hyperref[examples-section-phantom]{Exemples}
\item \hyperref[exercises-section-phantom]{Exercices}
\item \hyperref[guide-section-phantom]{Guide bibliographique}
\item \hyperref[desirables-section-phantom]{Desiderata}
\item \hyperref[coding-section-phantom]{Style de programmation}
\item \hyperref[obsolete-section-phantom]{Obsolète}
\item \hyperref[fdl-section-phantom]{Licence de documentation libre GNU}
\item \hyperref[index-section-phantom]{Index généré automatiquement}
\end{enumerate}
\end{multicols}


\bibliography{my}
\bibliographystyle{amsalpha}

\end{document}
